%%%%%%%%%%%%%%%%%%%%%%%%%%%%%%%%%%%%%%%%%%%%%%%%%%%%%%%%%%%%%%%%%%%%%%
% Lớp 11 | Hàm số mũ và hàm số lôgarít | Phép tính lôgarít
% Dạng 1D6?2-2: Biến đổi, biểu diễn biểu thức chứa lôgarít
% Số câu: 242  (sinh tự động bởi tach_cau.py; từ nay thầy sửa trực tiếp file này)
%%%%%%%%%%%%%%%%%%%%%%%%%%%%%%%%%%%%%%%%%%%%%%%%%%%%%%%%%%%%%%%%%%%%%%

%%%%%%%%%%%%%%%% TRẮC NGHIỆM 4 PHƯƠNG ÁN %%%%%%%%%%%%%%%%

% nguồn: Ok_Toan11_TracNghiem_Lan2
\begin{ex}%[1D6N2-2]%[ma:310549b3]
	Cho $a$, $b$, $c$ là các số thực dương và $a$, $b\ne1$. Khẳng định nào  sau đây là \textbf{sai}?
	\choice 
		 {$\log_{a}b \cdot \log_{b}a=1$}
		{\True  $\log_{a}c=-\log_{c}a$}
		{$\log_{a}c=\dfrac{\log_{b}c}{\log_{b}a}$} 
		 {$\log_{a}c=\log_{a}b \cdot \log_{b}c$}
	\loigiai{Ta  có $\log_{a}c=\dfrac{1}{\log_{c}a}$. Do đó khẳng định $\log_{a}c=-\log_{c}a$ là \textbf{sai}.}
\end{ex}

% nguồn: Ok_Toan11_TracNghiem_Lan2
\begin{ex}%[1D6N2-2]%[ma:5a5457ab]
	Với các  số thực dương $a$, $b$ bất kì. Mệnh đề nào dưới đây đúng
	\choice 
		 {\True $\log_{2}\left(\dfrac{2a^{3}}{b}\right)=1+3\log_{2}a-\log_{2}b$}
		 {$\log_{2}\left(\dfrac{2a^{3}}{b}\right)=1+\dfrac{1}{3}\log_{2}a-\log_{2}b$}
		 {$\log_{2}\left(\dfrac{2a^{3}}{b}\right)=1+3\log_{2}a+\log_{2}b$}
		{$\log_{2}\left(\dfrac{2a^{3}}{b}\right)=1+\dfrac{1}{3}\log_{2}a+\log_{2}b$} 
	 \loigiai{
		Ta có $\log_{2}\left(\dfrac{2a^{3}}{b}\right)=\log_{2}\left(2a^{3}\right)-\log_{2}b  =\log_{2}2+\log_{2}a^{3}-\log_{2}b=1+3\log_{2}a-\log_{2}b$.
	 }
\end{ex}

% nguồn: Ok_Toan11_TracNghiem_Lan2
\begin{ex}%[1D6N2-2]%[ma:8179bc53]
	Cho các  số thực dương $a, b$ thỏa mãn $3\log a+2\log b=1$. Mệnh đề nào sau đây đúng?
	\choice 
		{$a^{3}+b^{2}=1$} 
		{$3a+2b=10$} 
		{\True $a^{3}b^{2}=10$} 
		{$a^{3}+b^{2}=10$} 
	\loigiai{Ta có $3\log a+2\log b=1\Leftrightarrow \log a^{3}+\log b^{2}=1\Leftrightarrow  \log\left(a^{3}b^{2}\right)=1\Leftrightarrow a^{3}b^{2}=10$.}
\end{ex}

% nguồn: Ok_Toan11_TracNghiem_Lan2
\begin{ex}%[1D6N2-2]%[ma:78918c90]
	Với $a$ là số thực dương bất  kỳ, mệnh đề nào dưới đây đúng?
	\choice 
		{\True $\ln  a^{4}=4\ln a$}
		{$\ln\left(4a\right)=4\ln a$} 
		 {$\ln\left(4a\right)=\dfrac{1}{4}\ln a$}
		{$\ln a^{3}=\dfrac{1}{3}\ln  a$}
	\loigiai{Mệnh đề đúng  là $\ln a^{4}=4\ln a$.}
\end{ex}

% nguồn: Ok_Toan11_TracNghiem_Lan2
\begin{ex}%[1D6N2-2]%[ma:cb9ee9bc]
	Cho  $a$, $b$ là các số thực dương tùy ý. Khẳng định nào sau đây là đúng?
	 \choice
		{\True $\ln\left(ab\right)=\ln  a+\ln b$}
		 {$\ln\left(a+b\right)=\ln a+\ln b$}
		{$\ln\left(ab\right)=\ln a \cdot \ln b$} 
		{$\ln\left(a+b\right)=\ln a \cdot \ln b$} 
	\loigiai{Theo  quy tắc logarit ta có $\ln\left(ab\right)=\ln a+\ln b$ là khẳng định đúng.}
\end{ex}

% nguồn: Ok_Toan11_TracNghiem_Lan2
\begin{ex}%[1D6N2-2]%[ma:2b4c2ef6]
	Cho $0<a\ne1$, $x>0$. Mệnh đề  nào sau đây là \textbf{sai}?
	\choice 
		 {$\log_{a}a=1$}
		 {$\log_{a}a^{x}=x$}
		 {$\log_{a}1=0$}
		 {\True $x^{\log_{a}x}=x$}
	\loigiai{ 
		\begin{itemize}
			\item Với $0<a\ne1$, $x>0$  ta có các khẳng định $\log_{a}a=1$, $\log_{a}a^{x}=x$ và $\log_{a}1=0$ là đúng.
			\item Ta có $a^{\log_a x}=x$, do đó khẳng định $x^{\log_{a}x}=x$ là \textbf{sai}. 
		\end{itemize}
	 }
\end{ex}

% nguồn: Ok_Toan11_TracNghiem_Lan2
\begin{ex}%[1D6N2-2]%[ma:cf154284]
	Cho  $a$ là số thực dương khác $1$. Mệnh đề nào sau đây đúng với mọi số thực dương $x$, $y$?
	 \choice
		{$\log_{a}\dfrac{x}{y}=\log_{a}x+\log_{a}y$} 
		{$\log_{a}\dfrac{x}{y}=\log_{a}\left(x-y\right)$} 
		 {\True $\log_{a}\dfrac{x}{y}=\log_{a}x-\log_{a}y$}
		{$\log_{a}\dfrac{x}{y}=\dfrac{\log_{a}x}{\log_{a}y}$} 
	\loigiai{Theo quy tắc tính logarit của một thương, ta có $\log_{a}\dfrac{x}{y}=\log_{a}x-\log_{a}y$,  $\forall x>0$, $y>0$ là khẳng định đúng.}
\end{ex}

% nguồn: Ok_Toan11_TracNghiem_Lan2
\begin{ex}%[1D6N2-2]%[ma:772ba070]
	Cho ba số thực dương $a$, $b$, $c$ và $a\ne1$. Khẳng định nào sau đây  là \textbf{sai}?
	\choice 
		{$\log_{a}\left(bc\right)=\log_{a}b+\log_{a}c$} 
		{$a^{\log_{a}b}=b$} 
		{$\log_{a}b^{\alpha}=\alpha\log_{a}b$} 
		{\True $\log_{a}b=\dfrac{\ln a}{\ln  b}$}
	\loigiai{Ta có $\log_{a}b=\dfrac{\ln b}{\ln a}$, nên khẳng định  $\log_{a}b=\dfrac{\ln a}{\ln b}$ là \textbf{sai}.}
\end{ex}

% nguồn: Ok_Toan11_TracNghiem_Lan2
\begin{ex}%[1D6N2-2]%[ma:0eee29a2]
	Cho số $a>0$, $a\ne 1$ thỏa mãn $a^x=b$. Khi  đó đẳng thức nào sau đây đúng?
	\choice 
		{$a=\log_{x}b$} 
		 {$a=\log_bx$}
		{\True  $x=\log_ab$}
		 {$x=\log_ba$}
	 \loigiai{
		Ta có  $a^x=b\Leftrightarrow x=\log_ab$.
	} 
\end{ex}

% nguồn: Ok_TracNghiemLop11_Lan1
\begin{ex}%[1D6N2-2]%[ma:e2d744ac]
	Cho $a$, $b$, $c>0$, $a\ne 1$ và số $\alpha \in\mathbb{R}$. Trong các khẳng định sau, khẳng  định nào \textbf{sai}?
	\choice 
		 {$\log_aa=1$}
		 {\True $\log_a(b-c)=\log_ab-\log_ac$}
		{$\log_ab^\alpha=\alpha\log_ab$} 
		 {$\log_aa^c=c$}
	 \loigiai{
		Ta có $\log_a\dfrac{b}{c}=\log_ab-\log_ac$ nên $\log_a(b-c)=\log_ab-\log_ac$  là sai.
	 }
\end{ex}

% nguồn: Ok_TracNghiemLop11_Lan1
\begin{ex}%[1D6N2-2]%[ma:abb9b1ba]
	Với $ a$ là số thực khác $0$ tùy  ý, $\log_2a^2$ bằng
	\choice 
		 {$a$}
		 {$2\log_2a$}
		{\True $2\log_2|a|$} 
		 {$\dfrac{1}{2}\log_2a$}
	 \loigiai{
	Ta có $\log_2a^2=2\log_2|a|$  .}
\end{ex}

% nguồn: Ok_TracNghiemLop11_Lan1
\begin{ex}%[1D6N2-2]%[ma:29752f0c]
	Giá trị  của biểu thức $4^{\log_2\sqrt{3}}$ bằng
	\choice 
		{\True $3$} 
		{$\sqrt{3}$} 
		{$2^{\sqrt{3}}$} 
		{$2\sqrt{3}$} 
	 \loigiai{
	Ta có  $4^{\log_2\sqrt{3}}=(2^2)^{\log_2\sqrt{3}}=\left(2^{\log_2\sqrt{3}}\right)^2=(\sqrt{3})^2=3$.}
\end{ex}

% nguồn: Ok_TracNghiemLop11_Lan1
\begin{ex}%[1D6N2-2]%[ma:ac90e208]
	Với $a>0$, $\log(100a) + \log\left(\dfrac{10}{a}\right)$ bằng 
	\choice 
		{$1000$} 
		{$\log\left(100a+\dfrac{10}{a}\right)$} 
		 {\True $3$}
		{$1+2\log a$} 
	\loigiai{ 
		Ta có $\log(100a) + \log\left(\dfrac{10}{a}\right)= \log 100 + \log a + \log  10 - \log a = 3$.
	 }
\end{ex}

% nguồn: Ok_TracNghiemLop11_Lan1
\begin{ex}%[1D6N2-2]%[ma:b35050b2]
	Nếu $\log 4=a$ thì $\dfrac{1}{\log_{256}100}$  bằng
	 \choice
		 {$a^4$}
		 {$\dfrac{a}{8}$}
		{\True  $ 2a$}
		{$ 16a$} 
	 \loigiai{
	 Ta có $\dfrac{1}{\log_{256}100}=\log_{100}256=\dfrac{1}{2}\log{4^4}=2a$.}
\end{ex}

% nguồn: Ok_TracNghiemLop11_Lan1
\begin{ex}%[1D6N2-2]%[ma:6925788a]
	Trong các mệnh  đề sau, mệnh đề nào đúng?
	 \choice
		{\True $\log _a b^\alpha=\alpha \log _a b$ với mọi số $a$, $b$  dương và $a \neq 1$}
		{$\log _a b=\dfrac{1}{\log _b a}$ với mọi số $a$, $b$ dương và $a  \neq 1$}
		{$\log _a b+\log _a c=\log _a b  c$ với mọi số $a$, $b$ dương và $a \neq 1$}
		{$\log _a b=\dfrac{\log _c a}{\log _c b}$  với mọi số $a$, $b$, $c$ dương và $a \neq 1$}
	\loigiai{ 
		Ta có với  mọi số $a$, $b$ dương và $a \neq 1$ thì $\log _a b^\alpha=\alpha \log _a b$.
	} 
\end{ex}

% nguồn: Ok_TracNghiemLop11_Lan1
\begin{ex}%[1D6N2-2]%[ma:c4faf3e4]
	Cho $a$ là số thực dương khác $1$. Mệnh đề nào dưới đây  đúng với mọi số dương $x$, $y$?
	\choice 
		{\True ${\log }_a\dfrac{x}{y}={\log  }_ax-{{\log }_a}y$}
		{${\log }_a\dfrac{x}{y}={\log }_a\left( x-y  \right)$}
		{${\log }_a\dfrac{x}{y}={\log  }_ax+{{\log }_a}y$}
		{${\log  }_a\dfrac{x}{y}=\dfrac{{\log }_ax}{{\log }_ay}$}
	\loigiai{ 
		Ta có ${{\log }_a}\dfrac{x}{y}={{\log  }_a}x-{{\log }_a}y$.
	 }
\end{ex}

% nguồn: Ok_TracNghiemLop11_Lan1
\begin{ex}%[1D6N2-2]%[ma:eee9be3e]
	Cho $a,b,c > 0$, $a\ne  1$ và số $\alpha \in \mathbb{R}$, mệnh đề nào dưới đây {\bf sai}?
	\choice 
		{$\log_a a^c=c$} 
		{$\log_a a=1$} 
		{$\log_a b^{\alpha}=\alpha \log_a b$} 
		{\True $\log_a \left|b-c\right|=\log_a  b-\log_a c$}
	 \loigiai{
		Theo tính chất của logarit, mệnh đề sai là $\log_a  \left|b-c\right|=\log_a b-\log_a c$.
	} 
\end{ex}

% nguồn: Ok_TracNghiemLop11_Lan1
\begin{ex}%[1D6N2-2]%[ma:87c466b8]
	Cho $\log 2=a$  Tính $\log\dfrac{125}{4}$ theo $a$ .
	\choice 
		{\True $3-5a$} 
		 {$6+7a$}
		{$4\left(1+a\right)$} 
		{$2\left(a+5\right)$} 
	 \loigiai{
	Ta có  $\log\dfrac{125}{4}=\log\dfrac{1000}{32}=\log{10^3}-\log{2^5}=3-5a$.}
\end{ex}

% nguồn: Ok_TracNghiemLop11_Lan1
\begin{ex}%[1D6N2-2]%[ma:0c37cc2c]
	Mệnh đề nào sau đây đúng? 
	\choice 
		{$\log_a b^{\alpha}=\alpha \log_a b$ với mọi số thực  dương $a, b$}
		{$\log_a b^{\alpha}=\alpha \log_a b$  với mọi số thực $a, b$}
		{$\log_a b^{\alpha}=\alpha \log_a b$ với mọi số  thực $a, b$ và $a \neq 1$}
		{\True $\log_a b^{\alpha}=\alpha  \log_a b$ với mọi số thực dương $a, b$ và $a \neq 1$}
	\loigiai{ 
		Theo tính chất của logarit ta có mệnh đề  đúng là \lq\lq $\log_a b^{\alpha}=\alpha \log_a b$ với mọi số thực dương $a, b$ và $a \neq 1$\rq\rq.
	 }
\end{ex}

% nguồn: Ok_TracNghiemLop11_Lan1
\begin{ex}%[1D6N2-2]%[ma:071e6773]
	Cho $\log 5=a$.  Tính $\log 50\,000$ theo $ a$.
	 \choice
		{$ 5a$} 
		{\True  $ a+4$}
		 {$ 5a^2$}
		{$  2a^2+1$}
	 \loigiai{
	Ta có  $\log 50\,000=\log\left(5\cdot 10^4\right)=\log 5+4\log 10=a+4$.}
\end{ex}

% nguồn: Ok_TracNghiemLop11_Lan1
\begin{ex}%[1D6N2-2]%[ma:9072eabf]
	Với $a$ là số thực dương tùy ý, $\log\left(10a^2\right)$  bằng
	\choice 
		{$20\log a$} 
		 {\True $1+2\log a$}
		{$1+\left(\log  a\right)^2$}
		 {$10\log a$}
	 \loigiai{
		Với $a$ là số thực dương tùy ý, ta có $\log\left(10a^2\right)=\log  10+\log a^2=1+2\log a$.
	 }
\end{ex}

% nguồn: Ok_TracNghiemLop11_Lan1
\begin{ex}%[1D6N2-2]%[ma:8d3b3c2c]
	Với $a$, $b$ là hai số dương tùy  ý, $\log\left(a b^2\right)$ bằng
	\choice 
		{$2(\log a+\log  b)$}
		{$\log a+\dfrac{1}{2}\log b$} 
		{$2\log a+\log  b$}
		{\True  $\log a+2\log b$}
	 \loigiai{
		Ta  có $\log\left(a b^2\right)=\log a+\log b^2=\log a+2 \log b$.
	} 
\end{ex}

% nguồn: Ok_TracNghiemLop11_Lan1
\begin{ex}%[1D6N2-2]%[ma:c721996a]
	Với mọi  số thực $a$ dương, $\log_4a^4$ bằng
	\choice 
		 {$4$}
		 {$\dfrac{1}{4}$}
		 {$\dfrac{1}{4}\log_4a$}
		{\True  $4\log_4a$}
	\loigiai{ 
	Ta  có $\log_4a^4=4\log_4|a|=4\log_4a$.}
\end{ex}

% nguồn: Ok_TracNghiemLop11_Lan1
\begin{ex}%[1D6N2-2]%[ma:54901e2b]
	Với số thực dương  $m$ bất kỳ, khẳng định nào sau đây {\bf đúng}?
	 \choice
		{$\log_3(3m)=3+\log_3 m$} 
		{$\log_3(3m)=3\log_3  m$}
		 {$\log_3(3m)=9+\log_3 m$}
		{\True  $\log_3(3m)=1+\log_3 m$}
	 \loigiai{
		Ta có $\log_3(3m)=\log_3 3+\log_3  m=1+\log_3 m$.
	 }
\end{ex}

% nguồn: Ok_TracNghiemLop11_Lan1
\begin{ex}%[1D6N2-2]%[5-CTST]%[ma:c9ab0755]
	Trong các mệnh đề sau, mệnh đề nào  đúng?
	 \choice
		{\True $\log_{a}{b^\alpha} = \alpha \log_{a}{b}$ với mọi số $a$, $b$ dương và $a  \neq 1$}
		 {$\log_{a}{b} = \dfrac{1}{\log_{b}{a}}$ với mọi số $a$, $b$ dương và $a \neq 1$}
		 {$\log_{a}{b} + \log_{a}{c} = \log_{a}{bc}$ với mọi số $a$, $b$ dương và $a \neq 1$}
		{$\log_{a}{b} = \dfrac{\log_{c}{a}}{\log_{c}{b}}$  với mọi số $a$, $b$, $c$ dương và $a \neq 1$}
	 \loigiai{
	 }
\end{ex}

% nguồn: Ok_TracNghiemLop11_Lan1
\begin{ex}%[1D6N2-2]%[ma:8a7527d4]
	Cho $a>0$ và $a\ne 1$,  khi đó ${\log }_a\sqrt[4]a$ bằng
	\choice 
		 {$4$}
		{\True $\dfrac{1}{4}$} 
		 {$-\dfrac{1}{4}$}
		 {$-4$}
	\loigiai{ 
		 Ta có ${\log }_a\sqrt[4]a={{\log }_a}{a^{\tfrac{1}{4}}}=\dfrac{1}{4}$.
	} 
\end{ex}

% nguồn: Ok_TracNghiemLop11_Lan1
\begin{ex}%[1D6N2-2]%[ma:ce4a35a9]
	Với $ a$ là  số thực dương tùy ý, $\log_2\left(\dfrac{64}{a^4}\right)$ bằng
	\choice 
		{\True $  6-4\log_2a$}
		{$ 6+4\log_2a$} 
		 {$ 5+4\log_2a$}
		 {$ 6+\dfrac{1}{4}\log_2a$}
	 \loigiai{
	Ta có $\log_2\left(\dfrac{64}{a^4}\right)=\log_264-\log_2a^4=6-4\log_2a$.} 
\end{ex}

% nguồn: Ok_TracNghiemLop11_Lan1
\begin{ex}%[1D6N2-2]%[ma:e0b0453c]
	Cho $a$, $b$ là  hai số thực dương tùy ý khác $1$. Mệnh đề nào dưới đây \textbf{sai} ?
	 \choice
		{$\log_3 a-\log_3 b=\log_3 \dfrac{a}{b}$} 
		{\True $\log_3 a \cdot  \log_3 b=\log_3(a b)$}
		 {$\log_3 a+\log_3 \dfrac{1}{b}=\log_3 \dfrac{a}{b}$}
		 {$\log_3 b=\dfrac{1}{\log_b 3}$}
	\loigiai{$\log_3 a +  \log_3 b=\log_3(a b)$.}
\end{ex}

% nguồn: Ok_TracNghiemLop11_Lan1
\begin{ex}%[1D6N2-2]%[ma:87883691]
	Với  mọi số thực $a$ dương, $\log_2 \dfrac{a}{2}$ bằng
	\choice 
		{$\dfrac{1}{2} \log_2 a$} 
		{$\log_2 a+1$} 
		{\True$\log_2  a-1$}
		{$\log_2 a-2$} 
	 \loigiai{
		Ta có  $\log_2 \dfrac{a}{2}=\log_2 a-\log_2 2=\log_2 a-1$.
	} 
\end{ex}

% nguồn: Ok_Toan11_TracNghiem_Lan2
\begin{ex}%[1D6H2-2]%[ma:8eacbd52]
	[MĐ 103 2020-2021 - ĐỢT 2] Với $a>0$, đặt $\log_3(3a)=b$, khi đó $\log_3\left(9a^3\right)$  bằng
	\choice 
		{$3b$} 
		{\True $3b-1$} 
		 {$3b+2$}
		 {$3b+5$}
	\loigiai{ 
		Ta có $\log_3(3a)=b\Rightarrow 1+\log_3a=b\Rightarrow \log_3a=b-1$.  \\
		Suy ra $\log_3(9a^2)=2+3\log_3a=2+3(b-1)=3b-1$. 
	 }
\end{ex}

% nguồn: Ok_Toan11_TracNghiem_Lan2
\begin{ex}%[1D6H2-2]%[ma:fc14e85e]
	Với mọi $a$, $b$, $x$ là các số thực dương thỏa mãn $\log _3 x=2 \log _3 a+3 \log _3 b$, mệnh đề nào dưới  đây đúng?
	 \choice
		{$x=2 a+3  b$}
		{$x=3 a+2  b$}
		{$x=a^2  b^3$}
		{\True $x=a^2+b^3$} 
	\loigiai{ 
		Từ  giả thiết ta có
		\allowdisplaybreaks 
		\begin{eqnarray*}
			\log _3 x=2 \log _3 a+3 \log _3 b &\Leftrightarrow& \log _3 x=\log _3 a^2+\log  _3 b^3 \\&\Leftrightarrow& \log _3 x=\log _3\left(a^2 b^3\right) \Leftrightarrow x=a^2 b^3
	\end{eqnarray*}}
\end{ex}

% nguồn: Ok_Toan11_TracNghiem_Lan2
\begin{ex}%[1D6H2-2]%[ma:e61e990c]
	[Mã 123 2017] Cho $\log_ax=3$, $\log_bx=4$ với $a$, $b$ là các số thực lớn hơn $1$. Tính $P=\log_{ab}x$. 
	 \choice
		{$P=12$} 
		{\True  $P=\dfrac{12}{7}$}
		 {$P=\dfrac{7}{12}$}
		{$P=\dfrac{1}{12}$} 
	\loigiai{ 
		Ta có $P=\log_{ab}x=\dfrac{1}{\log_xab} 
		 =\dfrac{1}{\log_xa+\log_xb}
		 =\dfrac{1}{\dfrac{1}{3}+\dfrac{1}{4}}
		 =\dfrac{12}{7}$.
	} 
\end{ex}

% nguồn: Ok_Toan11_TracNghiem_Lan2
\begin{ex}%[1D6H2-2]%[ma:228f61ce]
	Với mọi $a$, $b$ thỏa mãn $\dfrac{\log _3 a \cdot \log _2 3}{1+\log _2 5}+\log b=1$. Khẳng định nào dưới  đây đúng?
	 \choice
		{ $a+b=1$} 
		{ $a=1-b \log  _2 5$}
		{\True $a  b=10$}
		{$a  \log _2 5+b=1$}
	\loigiai{ 
		Ta có \allowdisplaybreaks 
		\begin{eqnarray*}
			\dfrac{\log _3 a \cdot \log _2 3}{1+\log _2 5}+\log b=1 &\Leftrightarrow& \dfrac{\log _2  a}{\log _2 10}+\log b=1 \\&\Leftrightarrow& \log a+\log b=1 \\&\Leftrightarrow& \log a b=1 \Leftrightarrow a b=10.
		\end{eqnarray*}
	} 
\end{ex}

% nguồn: Ok_Toan11_TracNghiem_Lan2
\begin{ex}%[1D6H2-2]%[ma:d73170fc]
	Xét tất cả các số dương $a$ và $b$ thỏa mãn $\log _3 a+\log  _3 b=\log _{9}(a b)$. Tính giá trị của $a b$.
	 \choice
		{\True $a  b=1$}
		{$a b=2$} 
		 {$a b=\dfrac{1}{2}$}
		 {$a b=0$}
	\loigiai{ 
		Ta có 
		\allowdisplaybreaks 
		\begin{eqnarray*}
			\log _3 a+\log _3 b=\log _{9}(a  b) &\Leftrightarrow& \log _3(a b)=\log _{3^2}(a b) \\&\Leftrightarrow& \log _3(a b)=\dfrac{1}{2} \log _3(a b)
			\\&\Leftrightarrow&  \dfrac{1}{2} \log _3(a b)=0 \Leftrightarrow a b=1 .
		\end{eqnarray*}
	} 
\end{ex}

% nguồn: Ok_Toan11_TracNghiem_Lan2
\begin{ex}%[1D6H2-2]%[ma:7e41d3b1]
	[Mã 105 2017] Cho $\log_3a=2$  và $\log_2b=\dfrac{1}{2}$. \\
	Tính $I=2\log_3\left[\log_3(3a)\right]+\log_{\frac{1}{4}}b^2$. 
	\choice 
		 {$I=\dfrac{5}{4}$}
		{$I=0$} 
		{$I=4$} 
		{\True  $I=\dfrac{3}{2}$}
	 \loigiai{
		Ta  có $I=2\log_3\left[\log_3(3a)\right]+\log_{\frac{1}{4}}b^2
		 =2\log_3(\log_33+\log_3a)+2\log_{2^{-2}}b
		 =2-\dfrac{1}{2}=\dfrac{3}{2}$.
	} 
\end{ex}

% nguồn: Ok_Toan11_TracNghiem_Lan2
\begin{ex}%[1D6H2-2]%[ma:1d9c132d]
	[Đề Tham Khảo 2019] Đặt $\log_32=a$ khi đó  $\log_{16}27$ bằng
	 \choice
		{$\dfrac{3a}{4}$} 
		{\True $\dfrac{3}{4a}$} 
		 {$\dfrac{4}{3a}$}
		 {$\dfrac{4a}{3}$}
	\loigiai{ 
		Ta có  $\log_{16}27=\dfrac{3}{4}\log_23
		 =\dfrac{3}{4\cdot{\log_3}2}=\dfrac{3}{4a}$.
	 }
\end{ex}

% nguồn: Ok_Toan11_TracNghiem_Lan2
\begin{ex}%[1D6H2-2]%[ma:1f02e13a]
	 [MĐ 104 2020-2021 - ĐỢT 1] Với mọi $a$, $b$ thỏa mãn $\log_2a^3+\log_2b=5$, khẳng định nào dưới đây là \textbf{đúng}?
	 \choice
		 {\True $a^3b=32$}
		 {$a^3b=25$}
		{$a^3+b=25$} 
		 {$a^3+b=32$}
	\loigiai{ 
		Ta  có $\log_2a^3+\log_2b=5
		\Leftrightarrow \log_2 (a^3b)=5 
		\Leftrightarrow a^3b=32$. 
	 }
\end{ex}

% nguồn: Ok_Toan11_TracNghiem_Lan2
\begin{ex}%[1D6H2-2]%[ma:e3afb629]
	Cho các số thực dương $a$ và  $b$ thỏa mãn $a^{2}-16b=0$. Tính giá trị của biểu thức $P=\log_{\sqrt{2}}a-\log_{2}b$.
	 \choice
		{$P=2$} 
		 {\True $P=4$}
		 {$P=16$}
		{$P=\sqrt{2}$} 
	 \loigiai{
		 Ta có $P=\log_{\sqrt{2}}a-\log_{2}\dfrac{a^{2}}{16} =2\log_{2}a-\left(2\log_{2}a-\log_{2}16\right)=\log_{2}16=4$.
	 }
\end{ex}

% nguồn: Ok_Toan11_TracNghiem_Lan2
\begin{ex}%[1D6H2-2]%[ma:2e91fc67]
	[MĐ 104 2020-2021 - ĐỢT 2] Với  $a>0$, đặt $\log_3(3a)=b$ khi đó $\log_3(27a^4)$ bằng
	\choice 
		{$4b+3$} 
		{$4b$} 
		{\True  $4b-1$}
		 {$4b+7$}
	 \loigiai{
		Ta  có $\log_3(3a)=b\Leftrightarrow 1+\log_3a=b
		 \Rightarrow \log_3a=b-1$. \\
		Khi đó  $\log_3(27a^4)=\log_327+\log_3a^4=3+4\log_3a=3+4(b-1)=4b-1$.
	 }
\end{ex}

% nguồn: Ok_Toan11_TracNghiem_Lan2
\begin{ex}%[1D6H2-2]%[ma:feca26c7]
	Với $a$, $b$ là các số thực dương, khác $1$. Mệnh đề nào dưới đây  sai?
	 \choice
		 {$\log_3(ab)=\log_3a+\log_3b$}
		 {$\log_ab=\dfrac{\log_{2022}b}{\log_{2022}a}$}
		{$1-\log_ab=\log_a\dfrac{a}{b}$} 
		{\True $\log_ab^3=\dfrac{1}{3}\log_ab$} 
	 \loigiai{
		 Ta có $\log_ab^3=3\log_ab$.
	} 
\end{ex}

% nguồn: Ok_Toan11_TracNghiem_Lan2
\begin{ex}%[1D6H2-2]%[ma:5c432af3]
	Nếu $\log _2 x=5 \log _2  a+4 \log _2 b$ $(a$, $b>0)$ thì $x$ bằng
	\choice 
		{$a^4 b^5$} 
		 {$5 a+4 b$}
		{$4  a+5 b$}
		{\True  $a^5 b^4$}
	 \loigiai{
		Ta có $$\log _2 x=5 \log _2 a+4 \log _2 b  \Leftrightarrow \log _2 x=\log _2 a^5 b^4 \Leftrightarrow x=a^5 b^4.$$
	} 
\end{ex}

% nguồn: Ok_Toan11_TracNghiem_Lan2
\begin{ex}%[1D6H2-2]%[ma:a932725d]
	[Mã 103 2019] Cho $a$ và $b$ là hai số thực dương thỏa mãn $a^2b^3=16$.  Giá trị của $2\log_2a+3\log_2b$ bằng
	 \choice
		{$2$} 
		{$8$} 
		{$16$} 
		 {\True $4$}
	\loigiai{ 
		Ta có $2\log_2a+3\log_2b=\log_2\left(a^2b^3\right)=\log_216=4$. 
	 }
\end{ex}

% nguồn: Ok_Toan11_TracNghiem_Lan2
\begin{ex}%[1D6H2-2]%[ma:c4c77909]
	[Mã 101  2019] Cho $a$ và $b$ là hai số thực dương thỏa mãn $a^4b=16$. Giá trị của $4\log_2a+\log_2b$ bằng
	\choice 
		{\True  $4$}
		{$2$} 
		{$16$} 
		{$8$} 
	\loigiai{ 
		Ta  có $4\log_2a+\log_2b=\log_2a^4+\log_2b
		=\log_2(a^4b)=\log_216=\log_22^4=4$. 
	 }
\end{ex}

% nguồn: Ok_Toan11_TracNghiem_Lan2
\begin{ex}%[1D6H2-2]%[ma:883175ff]
	[MĐ 102 2020-2021 - ĐỢT  1] Với mọi $a$, $b$ thỏa mãn $\log _2a^3+\log_2b=8$. Khẳng định nào dưới đây \textbf{đúng}?
	\choice 
		{$a^3+b=64$} 
		 {\True $a^3b=256$}
		 {$a^3b=64$}
		{$a^3+b=256$} 
	 \loigiai{
		$\log _2a^3+\log_2b=8 
		 \Leftrightarrow \log_2(a^3b)=8
		\Leftrightarrow  a^3b=2^8
		 \Leftrightarrow a^3b=256$.
	 }
\end{ex}

% nguồn: Ok_Toan11_TracNghiem_Lan2
\begin{ex}%[1D6H2-2]%[ma:8d57b152]
	Cho $a>0$ và $a\ne 1$,  khi đó $\log_a\sqrt[4]{a}$ bằng
	\choice 
		 {$4$}
		{\True  $\dfrac{1}{4}$}
		{$-\dfrac{1}{4}$} 
		 {$4$}
	\loigiai{ 
		Với $a>0$ và $a\ne 1$ ta có: $\log_a\sqrt[4]{a}=\log_aa^{\frac{1}{4}}=\dfrac{1}{4}\log_aa=\dfrac{1}{4}$. 
	} 
\end{ex}

% nguồn: Ok_Toan11_TracNghiem_Lan2
\begin{ex}%[1D6H2-2]%[ma:a0ed50c4]
	[MĐ 102 2020-2021 - ĐỢT 1] Cho $a>0$ và $a\ne  1$, khi đó $\log_a\sqrt[3]{a}$ bằng
	 \choice
		 {$-3$}
		{\True $\dfrac{1}{3}$} 
		{$-\dfrac{1}{3}$} 
		 {$3$}
	 \loigiai{
		Với $a>0$ và $a\ne 1$,  ta có $\log_a\sqrt[3]{a}=\log_aa^{\frac{1}{3}}=\dfrac{1}{3}\log_aa=\dfrac{1}{3}$.
	 }
\end{ex}

% nguồn: Ok_Toan11_TracNghiem_Lan2
\begin{ex}%[1D6H2-2]%[ma:dede3bce]
	Với mọi số thực $a$ dương tuỳ  ý, $4\log_a \sqrt{a}$ bằng
	\choice 
		{$-2\log  a$}
		{\True $2\log  a$}
		{$-4\log  a$}
		 {$8\log a$}
	\loigiai{ 
		Ta có $4\log_a \sqrt{a}=4\log_a  a^{\tfrac{1}{2}}=4\cdot\dfrac{1}{2} \log_a a=2\log a$.
	} 
\end{ex}

% nguồn: Ok_Toan11_TracNghiem_Lan2
\begin{ex}%[1D6H2-2]%[ma:466ce25e]
	Với $a>0$, đặt  $\log_2(2a)=b$, khi đó $\log_2(4a^3)$ bằng
	\choice 
		{$3b+5$} 
		 {$3b$}
		{$3b+2$} 
		 {\True $3b-1$}
	\loigiai{ 
		 Ta có $\log_2(2a)=\log_22+\log_2a=1+\log_2a
		\Rightarrow \log_2a=b-1$.\\ 
		 Khi đó $\log_2(4a^3)=\log_2\left[(2a)^2\cdot a\right]=\log_2(2a)^2+\log_2a=2b+b-1=3b-1$.
	} 
\end{ex}

% nguồn: Ok_Toan11_TracNghiem_Lan2
\begin{ex}%[1D6H2-2]%[ma:818e0a05]
	[Mã 123 2017] Với $a$, $b$ là các số thực dương tùy ý và $a$ khác $1$,\break đặt $P=\log_ab^3+\log_{a^2}b^6$. Mệnh đề nào dưới đây \textbf{đúng}? 
	 \choice
		{\True $P=6\log_ab$} 
		 {$P=27\log_ab$}
		 {$P=15\log_ab$}
		{$P=9\log_ab$} 
	\loigiai{ 
		Ta  có $P=\log_ab^3+\log_{a^2}b^6=3\log_ab+\dfrac{6}{2}\log_ab=6\log_ab$.
	} 
\end{ex}

% nguồn: Ok_Toan11_TracNghiem_Lan2
\begin{ex}%[1D6H2-2]%[ma:b0fb7f45]
	[MĐ 104 2020-2021 - ĐỢT 1] Cho $a>0$ và $\mathrm{a}\ne  1$, khi đó $\log_a\sqrt[5]{a}$ bằng
	\choice 
		{\True  $\dfrac{1}{5}$}
		 {$-\dfrac{1}{5}$}
		{$5$} 
		{$-5$} 
	 \loigiai{
		Ta  có $\log_a\sqrt[5]{a}=\log_aa^{\frac{1}{5}}=\dfrac{1}{5}$.
	 }
\end{ex}

% nguồn: Ok_Toan11_TracNghiem_Lan2
\begin{ex}%[1D6H2-2]%[ma:a98dbee6]
	[Đề Tham Khảo 2020 Lần 2] Xét số thực $a$ và $b$ thỏa mãn $\log_3(3^a\cdot  9^b)=\log_93$. Mệnh đề nào dưới đây \textbf{đúng}?
	\choice 
		 {$a+2b=2$}
		{$4a+2b=1$} 
		{$4ab=1$} 
		{\True $2a+4b=1$} 
	\loigiai{ 
		 Ta có
		$\begin{aligned}
		 \log_3(3^a\cdot 9^b)=\log_93
		&\Leftrightarrow \log_3(3^a\cdot  3^{2b})=\log_{3^2}3
		 \Leftrightarrow \log_33^{a+2b}=\log_33^{\frac{1}{2}} \\
		& \Leftrightarrow a+2b=\dfrac{1}{2} 
		\Leftrightarrow  2a+4b=1.
	\end{aligned}$
 }
\end{ex}

% nguồn: Ok_Toan11_TracNghiem_Lan2
\begin{ex}%[1D6H2-2]%[ma:152cb356]
	[Mã 104 - 2020 Lần 2] Với $a$, $b$ là các số thực dương  tùy ý thỏa mãn $\log_2a-2\log_4b=4$, mệnh đề nào dưới đây \textbf{đúng}?
	 \choice
		{$a=16b^2$} 
		{$a=8b$} 
		 {\True $a=16b$}
		 {$a=16b^4$}
	 \loigiai{
		 Ta có
		$\begin{aligned}
		\log_2a-2\log_4b=4 
		 & \Leftrightarrow {\log_2}a-2\log_{2^2}b=4
		\Leftrightarrow {\log_2}a-2\cdot \dfrac{1}{2}\log_2b=4 \\ 
		 & \Leftrightarrow {\log_2}a-\log_2b=4
		\Leftrightarrow  {\log_2}\dfrac{a}{b}=4 \\
		 & \Leftrightarrow \dfrac{a}{b}=2^4
		\Leftrightarrow  a=16b.
	\end{aligned}$
 }
\end{ex}

% nguồn: Ok_Toan11_TracNghiem_Lan2
\begin{ex}%[1D6H2-2]%[ma:c048aa51]
	Với $a>0$  đặt $\log_2(2a)=b$, khi đó $\log_2(8a^4)$ bằng
	 \choice
		 {$4b+7$}
		 {$4b+3$}
		{$4b$} 
		{\True $4b-1$} 
	 \loigiai{
		Ta  có $\log_2(2a)=b\Leftrightarrow 1+\log_2a=b\Leftrightarrow{\log_2}a=b-1$.\\
		Khi đó  $\log_2(8a^4)=3+\log_2a^4=3+4\log_2a=3+4(b-1)=4b-1$.\\
		 Vậy $\log_2(8a^4)=4b-1$.
	 }
\end{ex}

% nguồn: Ok_Toan11_TracNghiem_Lan2
\begin{ex}%[1D6H2-2]%[ma:0a689b09]
	[Mã 101  - 2020 Lần 1] Với $a$, $b$ là các số thực dương tùy ý và $a\ne 1$, $\log_{a^5}b$ bằng
	 \choice
		{$5\log_ab$} 
		{$\dfrac{1}{5}+\log_ab$} 
		 {$5+\log_ab$}
		{\True  $\dfrac{1}{5}\log_ab$}
	\loigiai{ 
		Ta có $\log_{a^5}b=\dfrac{1}{5}\log_ab$. 
	 }
\end{ex}

% nguồn: Ok_Toan11_TracNghiem_Lan2
\begin{ex}%[1D6H2-2]%[ma:a98f72c0]
	Với mọi $a$, $b$ dương thỏa  mãn $\log_{2}\sqrt{a}-\log_{2}b=3$, khẳng định nào dưới đây đúng?
	\choice 
		{\True $a=64b^{2}$} 
		 {$ab^{2}=64$}
		 {$\sqrt{a}-b=8$}
		 {$\dfrac{\sqrt{a}}{b}=3$}
	\loigiai{Ta có $\log_{2}\sqrt{a}-\log_{2}b=3\Leftrightarrow \log_{2}\dfrac{\sqrt{a}}{b}=3\Leftrightarrow\dfrac{\sqrt{a}}{b}=2^{3}\Leftrightarrow a=64b^{2}$.} 
\end{ex}

% nguồn: Ok_Toan11_TracNghiem_Lan2
\begin{ex}%[1D6H2-2]%[ma:8fc3e8a7]
	Với mọi số  thực $a$ dương, $\log_{4}a^{4}$ bằng
	 \choice
		{$4$} 
		{\True  $4\log_{4}a$}
		{$\dfrac{1}{4}\log_{4}a$} 
		{$\dfrac{1}{4}$} 
	\loigiai{Ta có $\log_{4}a^{4}=4\log_{4}|a|=4\log_{4}a$.} 
\end{ex}

% nguồn: Ok_Toan11_TracNghiem_Lan2
\begin{ex}%[1D6H2-2]%[ma:4f5f60c0]
	[Mã 102 - 2020 Lần 1] Cho $a$ và $b$ là hai số thực dương thỏa  mãn $4^{\log_2(ab)}=3a$. Giá trị của $ab^2$ bằng
	 \choice
		{\True $3$} 
		{$6$} 
		{$2$} 
		 {$12$}
	 \loigiai{
		 Từ giả thiết ta có: \\
		$\begin{aligned}
		4^{\log  _2(ab)}=3a
		 & \Leftrightarrow \log_2(ab) \cdot \log_24=\log_2(3a)
		\Leftrightarrow  2\left(\log_2 a+\log_2 b\right)=\log_2 a+\log_2 3 \\
		& \Leftrightarrow \log_2a+2  \log_2b=\log_23
		\Leftrightarrow \log_2\left(ab^2\right)=\log_23 
		\Leftrightarrow  ab^2=3.
	\end{aligned}$
} 
\end{ex}

% nguồn: Ok_Toan11_TracNghiem_Lan2
\begin{ex}%[1D6H2-2]%[ma:5d458c3d]
	[Mã 110 2017] Cho $\log_ab=2$ và  $\log_ac=3$. Tính $P=\log_a(b^2c^3)$.
	 \choice
		{\True $P=13$} 
		{$P=31$} 
		{$P=30$} 
		{$P=108$} 
	\loigiai{ 
		Ta  có $\log_a(b^2c^3)=2\log_ab+3\log_ac=2\cdot 2+3\cdot 3=13$.
	 }
\end{ex}

% nguồn: Ok_Toan11_TracNghiem_Lan2
\begin{ex}%[1D6H2-2]%[ma:7c15aea6]
	[MĐ 103 2020-2021 - ĐỢT 1] Cho $a>0$ và $a\ne 1$,  khi đó $\log_a\sqrt{a}$ bằng
	\choice 
		 {$2$}
		{$-2$} 
		{$-\dfrac{1}{2}$} 
		{\True  $\dfrac{1}{2}$}
	\loigiai{ 
		Với $a>0$ và $a\ne 1$,  ta có: $\log_a\sqrt{a}=\log_aa^{\frac{1}{2}}=\dfrac{1}{2}\log_aa=\dfrac{1}{2}$.
	} 
\end{ex}

% nguồn: Ok_Toan11_TracNghiem_Lan2
\begin{ex}%[1D6H2-2]%[ma:89ba27e8]
	Với mọi $a$, $b$ thỏa mãn $\log_2a-3\log_2b=2$, khẳng định nào dưới  đây đúng?
	 \choice
		{\True $a=4b^3$} 
		{$a=3b+4$} 
		 {$a=3b+2$}
		 {$a=\dfrac{4}{b^3}$}
	\loigiai{ 
		Điều kiện $a,  b>0$.\\
		Khi đó, \begin{eqnarray*}
		\log_2a-3\log_2b=2&\Leftrightarrow&\log_2a-\log_2b^3=2\\ 
		&\Leftrightarrow&\log_2\dfrac{a}{b^3}=2\\ 
		 &\Leftrightarrow&\dfrac{a}{b^3}=4\\
		 &\Leftrightarrow& a=4b^3.
	\end{eqnarray*}
} 
\end{ex}

% nguồn: Ok_Toan11_TracNghiem_Lan2
\begin{ex}%[1D6H2-2]%[ma:51b749ec]
	Cho hai số thực dương $a$, $b$ bất kì thỏa mãn $4 \ln ^2 a+9 \ln ^2 b=12 \ln a \cdot \ln  b$. Khẳng định nào dưới đây đúng?
	 \choice
		{$3 a=2 b$} 
		{\True  $a^2=b^3$}
		{$2  a=3 b$}
		{$a^3=b^2$} 
	 \loigiai{
		Ta  có
		 \allowdisplaybreaks
		\begin{eqnarray*}
			4 \ln ^2 a+9 \ln  ^2 b=12 \ln a \cdot \ln b &\Leftrightarrow&\left(2 \ln a-3 \ln b\right)^2=0 \\&\Leftrightarrow& 2 \ln a=3 \ln b \Leftrightarrow a^2=b^3.
	\end{eqnarray*}}
\end{ex}

% nguồn: Ok_Toan11_TracNghiem_Lan2
\begin{ex}%[1D6H2-2]%[ma:cb833bac]
	Cho $a$ và $b$ là hai số  thực dương thỏa mãn $a^2b^5 = 64$. Giá trị của $P=2\log_{2}a+5\log_{2}b$ là
	\choice 
		 {$P=7$}
		 {$P=64$}
		{\True $P=6$} 
		{$P=2$} 
	\loigiai{ 
		\begin{itemize}
			\item Theo bài  ra $a$, $b>0$, $a^{2}b^{5}=64\Leftrightarrow \log_{2}\left(a^{2}b^{5}\right)=\log_{2}64\Leftrightarrow2\log_{2}a+5\log_{2}b=6$.
			\item Vậy  $P=6$.
		\end{itemize}
	 }
\end{ex}

% nguồn: Ok_Toan11_TracNghiem_Lan2
\begin{ex}%[1D6H2-2]%[ma:b73f154d]
	Với mọi $a$, $b$ thỏa  mãn $\log_2a^3+\log_2b=6$, khẳng định nào sau đây \textbf{đúng}?
	\choice 
		{\True $a^3b=64$} 
		 {$a^3b=36$}
		{$a^3+b=64$} 
		 {$a^3+b=36$}
	\loigiai{ 
		Ta có $\log_2a^3+\log_2b=6 
		\Leftrightarrow \log_2  (a^3b)=6
		 \Leftrightarrow a^3b=2^6=64$.
	} 
\end{ex}

% nguồn: Ok_Toan11_TracNghiem_Lan2
\begin{ex}%[1D6H2-2]%[ma:826dc64d]
	[Mã 104 2019]  Cho $a$ và $b$ là hai số thực dương thỏa mãn $ab^3=8$. Giá trị của $\log_2a+3\log_2b$ bằng
	\choice 
		 {$6$}
		 {$2$}
		{\True  $3$}
		 {$8$}
	 \loigiai{
		Ta có $\log_2a+3\log_2b=\log_2a+\log_2b^3=\log_2(ab^3)=\log_28=3$. 
	} 
\end{ex}

% nguồn: Ok_Toan11_TracNghiem_Lan2
\begin{ex}%[1D6H2-2]%[ma:3515aff8]
	[Mã 102 - 2020 Lần 1] Với $a, b$ là các số thực dương tùy  ý và $a\neq 1$, $\log_{a^2}b$ bằng
	\choice 
		{$\dfrac{1}{2}+\log _a  b$}
		 {\True $\dfrac{1}{2} \log _a b$}
		{$2+\log _a b$} 
		 {$2 \log _a b$}
	\loigiai{ 
		Ta có $\log _{a^2} b=\dfrac{1}{2}\log  _a b$.
	} 
\end{ex}

% nguồn: Ok_Toan11_TracNghiem_Lan2
\begin{ex}%[1D6H2-2]%[ma:38211f97]
	Cho $\log_{a}x=3$, $\log_{b}x=4$ với $a$, $b$ là các số thực lớn hơn 1. Tính $P=\log_{ab}x$. 
	 \choice
		 {$P=12$}
		 {$P=\dfrac{7}{12}$}
		{$P=\dfrac{1}{12}$} 
		 {\True $P=\dfrac{12}{7}$}
	 \loigiai{
		Ta  có $P=\log_{ab}x=\dfrac{1}{\log_{x}\left(ab\right)} =\dfrac{1}{\log_{x}a+\log_{x}b} =\dfrac{1}{\dfrac{1}{\log_{a}x}+\dfrac{1}{\log_{b}x}}=\dfrac{\log_{a}x \cdot \log_{b}x}{\log_{a}x+\log_{b}x}=\dfrac{12}{7} $.
	} 
\end{ex}

% nguồn: Ok_Toan11_TracNghiem_Lan2
\begin{ex}%[1D6H2-2]%[ma:15a64361]
	[Mã 101 2018]  Với $a$ là số thực dương tùy ý, $\ln (5a)-\ln (3a)$ bằng
	 \choice
		{\True $\ln  \dfrac{5}{3}$}
		{$\dfrac{\ln 5}{\ln 3}$} 
		{$\dfrac{\ln (5a)}{\ln (3a)}$} 
		 {$\ln (2a)$}
	\loigiai{ 
		Ta có $\ln (5a)-\ln  (3a)=\ln \dfrac{5}{3}$.
	} 
\end{ex}

% nguồn: Ok_Toan11_TracNghiem_Lan2
\begin{ex}%[1D6H2-2]%[ma:2dd71ab2]
	Cho $a,b,c>0$, $a\ne1$ và $\log_{a}b=\sqrt{2\,022}$.  Tính $\log_{\sqrt[6]{a}}\left(a^{\tfrac{7}{4}}\cdot\sqrt[6]{b}\right)$.
	\choice 
		 {$\dfrac{42+\sqrt{2\,022}}{6}$}
		{$\dfrac{7}{4}+6\sqrt{2\,022}$} 
		{\True $\dfrac{21}{2}+\sqrt{2\,022}$} 
		{$\dfrac{21}{2}+\dfrac{\sqrt{2\,022}}{6}$} 
	\loigiai{Ta có $\log_{\sqrt[6]{a}}\left(a^{\tfrac{7}{4}}\cdot\sqrt[6]{b}\right)=\log_{\sqrt[6]{a}}a^{\tfrac{7}{4}}+\log_{\sqrt[6]{a}}\sqrt[6]{b}=6\cdot\dfrac{7}{4}+\sqrt{2\,022}=\dfrac{21}{2}+\sqrt{2\,022}$.} 
\end{ex}

% nguồn: Ok_Toan11_TracNghiem_Lan2
\begin{ex}%[1D6H2-2]%[ma:2b51c20e]
	[Đề Minh Họa 2017] Đặt $a=\log_23$, $b=\log_53$. Hãy biểu diễn $\log_645$ theo $a$ và  $b$.
	\choice 
		{$\log_645=\dfrac{2a^2-2ab}{ab}$} 
		{\True $\log_645=\dfrac{a+2ab}{ab+b}$} 
		 {$\log_645=\dfrac{2a^2-2ab}{ab+b}$}
		{$\log_645=\dfrac{a+2ab}{ab}$} 
	\loigiai{ 
		Ta có 
		$\begin{aligned}
		\log_645&=\dfrac{\log_2(3^2\cdot  5)}{\log_2(2\cdot 3)}
		 =\dfrac{2\log_23+\log_25}{1+\log_23} \\
		 &=\dfrac{2a+\log_23\cdot{\log_3}5}{1+a}
		=\dfrac{2a+\dfrac{\log_23}{\log_53}}{1+a}  \\
		 &=\dfrac{2a+\dfrac{a}{b}}{1+a}=\dfrac{a+2ab}{ab+b}\cdot
	\end{aligned}$
 }
\end{ex}

% nguồn: Ok_Toan11_TracNghiem_Lan2
\begin{ex}%[1D6H2-2]%[ma:9eb7cba1]
	[MĐ 103 2020-2021 - ĐỢT 1] Với mọi $a$, $b$ thỏa mãn $\log_2a^3+\log_2b=7$. Khẳng định nào dưới đây  \textbf{đúng}?
	 \choice
		{$a^3+b=49$} 
		 {\True $a^3b=128$}
		 {$a^3+b=128$}
		 {$a^3b=49$}
	\loigiai{ 
		Điều kiện $a>0$, $b>0$.  \\
		 Ta có
		$\begin{aligned}
		 \log_2a^3+\log_2b=7
		&\Leftrightarrow  \log_2 (a^3b)=7 \\
		 &\Leftrightarrow a^3b=2^7
		\Leftrightarrow a^3b=128. 
	\end{aligned}$
} 
\end{ex}

% nguồn: Ok_Toan11_TracNghiem_Lan2
\begin{ex}%[1D6H2-2]%[ma:b88a9243]
	Cho $a$, $b$ là các số thực dương và $a$ khác $1$, thỏa mãn $\log_{a^3}\dfrac{a^5}{\sqrt[4]{b}}=2$. Giá trị của  biểu thức $\log_ab$ bằng
	\choice 
		 {$4$}
		{$\dfrac{1}{4}$} 
		 {$-\dfrac{1}{4}$}
		 {\True $-4$}
	\loigiai{ 
		Ta có $2=\log_{a^3}\dfrac{a^5}{\sqrt[4]{b}}=\dfrac{1}{3}\log_a\dfrac{a^5}{b^{\tfrac{1}{4}}}=\dfrac{1}{3}\left(\log_aa^5-\log_ab^{\tfrac{1}{4}}\right)=\dfrac{1}{3}\left(5-\dfrac{1}{4}\log_ab\right)$.\\ 
		Suy  ra $5-\dfrac{1}{4}\log_ab=6\Rightarrow \log_ab=-4$.
	 }
\end{ex}

% nguồn: Ok_Toan11_TracNghiem_Lan2
\begin{ex}%[1D6H2-2]%[ma:5b1e262b]
	[TK 2020-2021] Với $a$ là số thực dương tùy ý, $\log_3(9a)$ bằng 
	\choice 
		{$\dfrac{1}{2}+\log_3a$} 
		{$2\log_3a$} 
		{$(\log_3a)^2$} 
		 {\True $2+\log_3a$}
	 \loigiai{
		Ta  có $\log_3(9a)=\log_39+\log_3a=2+\log_3a$.
	} 
\end{ex}

% nguồn: Ok_Toan11_TracNghiem_Lan2
\begin{ex}%[1D6H2-2]%[ma:cdccbbe7]
	[Mã 104 2017] Với mọi $a$, $b$, $x$ là các số thực  dương thoả mãn \break $\log{_2}x=5\log_2a+3\log_2b$. Mệnh đề nào dưới đây \textbf{đúng}?
	 \choice
		 {$x=5a+3b$}
		 {$x=a^5+b^3$}
		{\True  $x=a^5b^3$}
		{$x=3a+5b$} 
	 \loigiai{
		 Có $\log_2x=5\log_2a+3\log_2b=\log_2a^5+\log_2b^3=\log_2a^5b^3
		 \Leftrightarrow x=a^5b^3$.
	 }
\end{ex}

% nguồn: Ok_Toan11_TracNghiem_Lan2
\begin{ex}%[1D6H2-2]%[ma:5175c01a]
	Cho số thực  dương $a$ khác $1$. Giá trị của biểu thức $\log_2(4a)$ bằng
	 \choice
		 {$4+\log_2a$}
		{\True  $2+\log_2a$}
		 {$2\log_2a$}
		{$4\log_2a$} 
	\loigiai{ 
		Ta  có $\log_2(4a)=\log_24+\log_2a=2+\log_2a$.
	 }
\end{ex}

% nguồn: Ok_Toan11_TracNghiem_Lan2
\begin{ex}%[1D6H2-2]%[ma:8c17018e]
	[Mã 103 2018]  Với $a$ là số thực dương tùy ý, $\ln (7a)-\ln (3a)$ bằng
	\choice 
		{$\dfrac{\ln 7}{\ln  3}$}
		{\True  $\ln \dfrac{7}{3}$}
		{$\ln (4a)$} 
		{$\dfrac{\ln (7a)}{\ln  (3a)}$}
	 \loigiai{
		Ta có $\ln (7a)-\ln  (3a) =\ln \dfrac{7a}{3a}=\ln \dfrac{7}{3}$.
	 }
\end{ex}

% nguồn: Ok_Toan11_TracNghiem_Lan2
\begin{ex}%[1D6H2-2]%[ma:1bf710d0]
	Với $a$  là số thực dương tùy ý, $\log (100a)$ bằng
	 \choice
		{$1-\log  a$}
		{\True  $2+\log a$}
		{$2-\log a$} 
		{$1+\log  a$}
	 \loigiai{
		Ta có $\log (100a)=\log (100)+\log a=2+\log  a$
	 }
\end{ex}

% nguồn: Ok_Toan11_TracNghiem_Lan2
\begin{ex}%[1D6H2-2]%[ma:94459c5d]
	[Mã 103 - 2020 Lần 2] Với $a$, $b$ là các số thực  dương tùy ý thỏa mãn $\log_3a-2\log_9b=3$, mệnh đề nào dưới đây \textbf{đúng}?
	\choice 
		 {\True $a=27b$}
		{$a=9b$} 
		 {$a=27b^4$}
		{$a=27b^2$} 
	\loigiai{ 
		Ta có $\log_3a-2\log_9b=3 
		\Leftrightarrow {\log_3}a-\log_3b=3 
		\Leftrightarrow {\log_3}\dfrac{a}{b}=3 
		\Leftrightarrow \dfrac{a}{b}=27 
		\Leftrightarrow a=27b$. 
	 }
\end{ex}

% nguồn: Ok_Toan11_TracNghiem_Lan2
\begin{ex}%[1D6H2-2]%[ma:05f4470b]
	Cho $a,b,c$ là các số thực dương, trong đó $a,b >  1$ và thỏa mãn $\log_{a}c=3$, $\log_{b}c=4$. Tính giá trị biểu thức $P=\log_{ab}c$.
	\choice 
		{\True  $P=\dfrac{12}{7}$}
		 {$P=\dfrac{7}{12}$}
		{$P=\dfrac{1}{12}$} 
		 {$P=12$}
	 \loigiai{
		\begin{itemize}
			\item Từ giả thiết ta suy ra $c\ne1$ và $\log_{c}a=\dfrac{1}{\log_{a}c}=\dfrac{1}{3}$;  $\log_{c}b=\dfrac{1}{\log_{b}c}=\dfrac{1}{4}$.
			 \item Khi đó $P=\log_{ab}c=\dfrac{1}{\log_{c}\left(ab\right)}=\dfrac{1}{\log_{c}a+\log_{c}b}=\dfrac{1}{\dfrac{1}{3}+\dfrac{1}{4}}=\dfrac{12}{7}$.
		\end{itemize}
	 }
\end{ex}

% nguồn: Ok_Toan11_TracNghiem_Lan2
\begin{ex}%[1D6H2-2]%[ma:d4e6a1f9]
	[Đề Tham Khảo 2017] Cho $a$, $b$ là các số thực dương thỏa mãn  $a\ne 1$, $a\ne \sqrt{b}$ và $\log_ab=\sqrt{3}$. \break Tính $P=\log_{\frac{\sqrt{b}}{a}}\sqrt{\dfrac{b}{a}}$.
	\choice 
		 {$P=-5+3\sqrt{3}$}
		 {$P=-1+\sqrt{3}$}
		 {\True $P=-1-\sqrt{3}$}
		{$P=-5-3\sqrt{3}$} 
	\loigiai{ 
		Cách 1:  Phương pháp tự luận.\\
		 Ta có
		$\begin{aligned}
		 P&=\dfrac{\log_a\sqrt{\dfrac{b}{a}}}{\log_a\dfrac{\sqrt{b}}{a}}
		=\dfrac{\dfrac{1}{2}(\log_ab-1)}{\log_a\sqrt{b}-1}  \\
		&=\dfrac{\dfrac{1}{2}(\sqrt{3}-1)}{\dfrac{1}{2}\log_ab-1} 
		=\dfrac{\sqrt{3}-1}{\sqrt{3}-2} 
		=-1-\sqrt{3}. 
	\end{aligned}$ \\
	Cách 2: Phương pháp  trắc nghiệm.\\
	Chọn $a=2$, $b=2^{\sqrt{3}}$. Bấm máy  tính ta được $P=-1-\sqrt{3}$.
 }
\end{ex}

% nguồn: Ok_Toan11_TracNghiem_Lan2
\begin{ex}%[1D6H2-2]%[ma:13c29cc5]
	Đặt $a=\log_23$  và $ b=\log_53$. Hãy biểu diễn $\log_645$ theo $a$ và $b$.
	 \choice
		 {\True $\log_645=\dfrac{a+2ab}{ab+b}$}
		 {$\log_645=\dfrac{2a^2-2ab}{ab}$}
		{$\log_645=\dfrac{a+2ab}{ab}$} 
		 {$\log_645=\dfrac{2a^2-2ab}{ab+b}$}
	 \loigiai{
		 Ta có $\log_645=\dfrac{\log_3(5\cdot 3^2)}{\log_3(2\cdot 3)}=\dfrac{\log_35+2}{\log_32+1}=\dfrac{\dfrac{1}{b}+2}{\dfrac{1}{a}+1}=\dfrac{a+2ab}{ab+b}$.
	 }
\end{ex}

% nguồn: Ok_TracNghiemLop11_Lan1
\begin{ex}%[1D6H2-2]%[ma:61b7b49a]
	Cho $ a,b$ là các số thực dương thỏa mãn $ a\ne 1$, $ a\ne \sqrt  {b}$ và ${{\log}_a}b=\sqrt{3}$. Tính $P={{\log}_{\tfrac{\sqrt {b}}{a}}}\sqrt{\dfrac{b}{a}}$.
	\choice 
		 {$P=-5+3\sqrt{3}$}
		{$P=-1+\sqrt{3}$} 
		 {\True $P=-1-\sqrt{3}$}
		 {$P=-5-3\sqrt{3}$}
	\loigiai{ 
		 Cách 1: Phương pháp tự luận\\
		$P=\dfrac{{{\log}_a}\sqrt{\dfrac{b}{a}}}{{{\log }_a}\dfrac{\sqrt {b}}{a}}=\dfrac{\dfrac{1}{2}\left( {{\log }_a}b-1 \right)}{{{\log}_a}\sqrt {b}-1}=\dfrac{\dfrac{1}{2}\left(\sqrt {3}-1 \right)}{\dfrac{1}{2}{{\log }_a}b-1}=\dfrac{\sqrt {3}-1}{\sqrt {3}-2}$$=-1-\sqrt  {3}$.\\
		Cách 2: Phương pháp trắc  nghiệm.\\
		Chọn $a=2$,  $b={2^{\sqrt{3}}}$. Bấm máy tính ta được $P=-1-\sqrt{3}$.
	} 
\end{ex}

% nguồn: Ok_TracNghiemLop11_Lan1
\begin{ex}%[1D6H2-2]%[ma:ee63e0d0]
	Với $ a,b$ là các số thực  dương tùy ý và ($ a\ne 1$) thì $\log_{a^2}b$ bằng
	\choice 
		{$ 2\log_ab$} 
		 {$\dfrac{1}{2}+\log_ab$}
		{\True $\dfrac{1}{2}{\log_a}b$} 
		{$ 2+\log_ab$} 
	\loigiai{ 
	Ta có $\log_{a^2}b=\dfrac{1}{2}{\log_a}b$.} 
\end{ex}

% nguồn: Ok_TracNghiemLop11_Lan1
\begin{ex}%[1D6H2-2]%[ma:4d754547]
	Cho $a=\log_35$, $b=\log_310$. Khi đó giá  trị của $\log_{\sqrt{3}}50$ được tính theo $a$, $b$ là
	 \choice
		{\True $2(a+b)$} 
		{$2(a-b)$} 
		 {$2a-b$}
		 {$2a+b$}
	 \loigiai{
		 Ta có $\log_{\sqrt{3}}50=2\log_350=2\left(\log_35+\log_310\right)=2(a+b)$.
	 }
\end{ex}

% nguồn: Ok_TracNghiemLop11_Lan1
\begin{ex}%[1D6H2-2]%[ma:c2bc9291]
	Với hai số thực dương $a$, $b$ tùy ý  và $\dfrac{\log_3 a}{\log_3 6}=\log_6 b +2$. Khi đó
	 \choice
		 {$a+36b=0$}
		 {$36a+b=0$}
		 {$36a-b=0$}
		{\True  $a-36=0$}
	\loigiai{ 
		Ta có $\dfrac{\log_3 a}{\log_3 6}=\log_6 b +2 \Leftrightarrow \log_6 a = \log_6 b +2 \Leftrightarrow \log_6 a - \log_6 b=2 \Leftrightarrow \log \dfrac{a}{b}=2  \Leftrightarrow a=36b \Leftrightarrow a-36b=0$.
	} 
\end{ex}

% nguồn: Ok_TracNghiemLop11_Lan1
\begin{ex}%[1D6H2-2]%[ma:3f454fe4]
	Với các số thực dương  $a, b$ bất kì, mệnh đề nào sau đây là đúng?
	\choice 
		{\True $\log (a b)=\log  a+\log b$}
		{$\log  \dfrac{a}{b}=\dfrac{\log a}{\log b}$}
		{$\log  (a b)=\log a \cdot \log b$}
		{$\log  \dfrac{a}{b}=\log b-\log a$}
	 \loigiai{
		Với các số thực dương $a, b$ bất kì, ta có $\log (ab)=\log a+\log  b$.
	 }
\end{ex}

% nguồn: Ok_TracNghiemLop11_Lan1
\begin{ex}%[1D6H2-2]%[ma:d30e9abc]
	Xét tất cả các số dương $a$ và $b$ thỏa mãn $\log_2 a=\log_{16}\left(ab^5\right)$. Hãy chọn khẳng  định đúng.
	\choice 
		{\True $a^3=b^5$} 
		{$a=b$} 
		{$a^5=b^3$} 
		 {$a^4=b^5$}
	\loigiai{\allowdisplaybreaks\begin{eqnarray*}
		&& \log_2 a=\log_{16}\left(ab^5\right)\\ 
		 & \Rightarrow & \log_2 a=\log_{2^4}\left(ab^5\right)\\
		& \Rightarrow &  \log_2 a=\dfrac{1}{4}\log_{2}\left(ab^5\right)\\
		& \Rightarrow &  4\log_2 a=\log_{2}\left(ab^5\right)\\
		& \Rightarrow & 4\log_2 a=\log_{2}a+\log_{2}\left(b^5\right)\\ 
		& \Rightarrow  & 3\log_2 a=\log_{2}\left(b^5\right)\\
		& \Rightarrow & \log_2  \left( a^3\right) =\log_{2}\left(b^5\right)\\
		& \Rightarrow &  a^3=b^5.
\end{eqnarray*}}
\end{ex}

% nguồn: Ok_TracNghiemLop11_Lan1
\begin{ex}%[1D6H2-2]%[ma:511f6c48]
	Đặt $\log _5 2=a$, $\log _5 3=b$. Biểu thị theo $a$ và $b$ ta  được $\log _8 27$ bằng
	 \choice
		{\True $\dfrac{b}{a}$} 
		 {$\dfrac{a}{b}$}
		{$\dfrac{3 a}{b}$} 
		 {$\dfrac{a}{3b}$}
	 \loigiai{
		Ta có $\log _8 27=\log _{2^{3}}  3^{3}=\log_2 3=\dfrac{\log _5 3}{\log _5 2}=\dfrac{b}{a}$.
	} 
\end{ex}

% nguồn: Ok_TracNghiemLop11_Lan1
\begin{ex}%[1D6H2-2]%[ma:ebebb8bd]
	Cho hàm số $f(x)=\log_2 x$. Biết rằng $m$, $n$ là các số thực dương thỏa mãn $f(n)-f(m)=1$  và $9^m \cdot 27^n=9$. Giá trị của biểu thức $16 m^2+4 n^2$ bằng
	\choice 
		{$4$}{$1$}{$0$}{\True  $2$}
	\loigiai{ 
		Ta  có
		 \allowdisplaybreaks
		\begin{eqnarray*}
			f(n)-f(m)=1\Leftrightarrow \log_2 n-\log_2 m=1 &\Leftrightarrow & \dfrac{n}{m}=2\Leftrightarrow n=2m\Leftrightarrow n^2=4m^2\Leftrightarrow  4n^2=16m^2.
		\end{eqnarray*}
		Lại  có
		 $9^m \cdot 27^n=9\Leftrightarrow 9^m \cdot 27^{2m}=9\Leftrightarrow \left( 9\cdot 27^2\right) ^m=9\Leftrightarrow m=\log_{9\cdot 27^2} 9$.\\
		Ta có  $16 m^2+4 n^2=32m^2=32\cdot \left( \log_{9\cdot 17^2} 9\right)^2=2$.
	} 
\end{ex}

% nguồn: Ok_TracNghiemLop11_Lan1
\begin{ex}%[1D6H2-2]%[ma:693392b1]
	Với $a>0$, $a \neq 1$, cho $f(a)=2\ln a+3 \log_a  e-\dfrac{3}{\ln a}-\dfrac{2}{\log_a e}$. Giá trị của $f\left(5^{302}+\sqrt{7}\right)$ là
	\choice 
		 {$f\left(5^{302}+\sqrt{7}\right)=1$}
		{$f\left(5^{302}+\sqrt{7}\right)=5^{302}+\sqrt{7}$} 
		 {$f\left(5^{302}+\sqrt{7}\right)=\ln \left(5^{302}+\sqrt{7}\right)$}
		{\True  $f\left(5^{302}+\sqrt{7}\right)=0$}
	\loigiai{$f(a)=2\ln  a+3 \log_a e-\dfrac{3}{\ln a}-\dfrac{2}{\log_a e}
		=2\ln a+3  \log_a e-3\log_a e-2\ln a=0.\\
	f\left(5^{302}+\sqrt{7}\right)=0$.} 
\end{ex}

% nguồn: Ok_TracNghiemLop11_Lan1
\begin{ex}%[1D6H2-2]%[ma:d16af600]
	Cho các số thực dương $ a$, $ b$ thỏa mãn $\log_2a=x$, $\log_2b=y$.\\ Tính $ P=\log_2\left(a^2b^3\right)$. 
	 \choice
		 {$ P=x^2y^3$}
		{$  P=x^2+y^3$}
		{$ P=6xy$} 
		{\True  $ P=2x+3y$}
	 \loigiai{
	$ P=\log_2\left(a^2b^3\right)=\log_2a^2+\log_2b^3=2\log_2a+3\log_2b=2x+3y$.} 
\end{ex}

% nguồn: Ok_TracNghiemLop11_Lan1
\begin{ex}%[1D6H2-2]%[ma:1d664c85]
	Cho $ a > 0$, $b  > 0 $ thỏa mãn $ a^2+b^2 = 7ab $. Chọn mệnh đề đúng trong các mệnh đề sau:
	\choice 
		{$\log (a+b) = \dfrac{3}{2} \left(\log  a+\log b\right)$}
		{$2 \left(\log a+\log b\right) = \log  (7ab) $}
		{$3 \log (a+b) = \dfrac{1}{2} \left(\log  a+\log b\right)$}
		{\True $\log  \left(\dfrac{a+b}{3} \right) = \dfrac{1}{2} \left(\log a+\log b\right)$}
	 \loigiai{
		Ta có 
		\begin{eqnarray*}
			& &a^2+b^2 =  7ab \\
			&\Leftrightarrow&  \left(a+b\right)^2 = 9ab\\
			&\Leftrightarrow& 2 \log (a+b) = \log  (9ab) \\
			 &\Leftrightarrow& 2 \log (a+b) = 2 \log 3+\log a+\log b\\
			&\Leftrightarrow& \log (a+b)-\log 3  = \dfrac{1}{2} \left(\log a+\log b\right) \\
			&\Leftrightarrow&  \log \dfrac{a+b}{3} = \dfrac{1}{2} \log ab.
		\end{eqnarray*}
	} 
\end{ex}

% nguồn: Ok_TracNghiemLop11_Lan1
\begin{ex}%[1D6H2-2]%[ma:617bf2fb]
	Cho $\log_ab=3,\log_ac=-2$. Khi đó $\log_a\left(a^3b^2\sqrt c\right)$  bằng
	\choice 
		{$  5$}
		{\True  $ 8$}
		 {$ 10$}
		{$ 13$} 
	 \loigiai{
	Ta có $\log_a\left(a^3b^2\sqrt c\right)=\log_a{a^3}+\log_a{b^2}+\log_a\sqrt c =3+2\log_ab+\dfrac{1}{2}{\log_a}c=3+2\cdot 3-\dfrac{1}{2}\cdot 2=8$.} 
\end{ex}

% nguồn: Ok_TracNghiemLop11_Lan1
\begin{ex}%[1D6H2-2]%[ma:764ae8d1]
	Với $a$ và $b$ là các số thực dương và $a \neq 1$. Biểu thức $\log _{a}\left(a^2 b\right)$  bằng
	\choice 
		{$2 \log  _{a}b$}
		{$1+2 \log _{a}b$} 
		 {$2-\log _{a}b$}
		{\True $2+\log  _{a}b$}
	 \loigiai{
		Ta có $\log _{a}\left(a^2 b\right)= \log _{a} a^2 + \log _{a} b=2\log _{a} a+\log  _{a}b=2+\log _{a}b$.
	} 
\end{ex}

% nguồn: Ok_TracNghiemLop11_Lan1
\begin{ex}%[1D6H2-2]%[ma:0320f48c]
	 Cho $\log_25=a$; $\log_53=b$. Tính $\log_524$ theo $a$ và $b$.
	\choice 
		{\True  $\log_524=\dfrac{3+ab}{a}$}
		{$\log_524=\dfrac{a+3b}{a}$} 
		{$\log_524=\dfrac{3a+b}{b}$} 
		 {$\log_524=\dfrac{a+b}{3ab}$}
	 \loigiai{
		Ta có \allowdisplaybreaks \begin{eqnarray*}
		\log_524&=&\log_5(8\cdot 3)=\log_58+\log_53\\ 
		&=&3\cdot\log_52+\log_53=\dfrac{3}{\log_25}+\log_53\\ 
		 &=&\dfrac{3}{a}+b=\dfrac{3+ab}{a}.
	\end{eqnarray*}
 }
\end{ex}

% nguồn: Ok_TracNghiemLop11_Lan1
\begin{ex}%[1D6H2-2]%[ma:cb4aadc7]
	Cho $a>0$ và $a\ne 1$ khi đó  $\log_a\sqrt[3]{a}$ bằng
	\choice 
		 {$-3$}
		 {\True $\dfrac{1}{3}$}
		 {$-\dfrac{1}{3}$}
		 {$3$}
	\loigiai{ 
		 $\log_a\sqrt[3]{a}=\dfrac{1}{3}\log_aa=\dfrac{1}{3}$.
	} 
\end{ex}

% nguồn: Ok_TracNghiemLop11_Lan1
\begin{ex}%[1D6H2-2]%[ma:1856d197]
	Cho $a$ và $b$ là hai  số thực dương thỏa mãn $a^4 b=16$. Giá trị của $4\log_2 a+\log_2 b$ bằng
	 \choice
		 {\True $4$}
		 {$2$}
		{$16$} 
		{$8$} 
	\loigiai{ 
		 $4\log_2 a+\log_2 b=\log_2 a^4+\log_2 b=\log_2 \left(a^4 b\right)=\log_2 16=\log_2 2^4=4$.
	 }
\end{ex}

% nguồn: Ok_TracNghiemLop11_Lan1
\begin{ex}%[1D6H2-2]%[ma:69253415]
	Cho các số thực dương  $a$, $b$ với $a\ne 1$. Khẳng định nào sau đây là khẳng định đúng?
	\choice 
		{$\log _{a^2}  \left(ab\right)=\dfrac{1}{4} \log_a b$}
		{\True$\log _{a^2} \left(ab\right)=\dfrac{1}{2}+\dfrac{1}{2} \log_a b$  }
		 {$\log _{a^2} \left(ab\right)=\dfrac{1}{2} \log_a b$}
		 {$\log _{a^2} \left(ab\right)=2+2\log_a b$}
	 \loigiai{
		Ta có $\log _{a^2} \left(ab\right)=\log _{a^2} a+\log _{a^2} b=\dfrac{1}{2}\cdot\log_a  a+\dfrac{1}{2}\cdot\log_a b=\dfrac{1}{2}+\dfrac{1}{2}\cdot\log_a b$.
	} 
\end{ex}

% nguồn: Ok_TracNghiemLop11_Lan1
\begin{ex}%[1D6H2-2]%[ma:23cc81f1]
	Cho $\log_25=a$. Giá trị của $\log_425$  theo $ a$ bằng
	 \choice
		 {$\dfrac{3}{2}a$}
		{\True  $a$}
		{$  3a$}
		{$  2a$}
	 \loigiai{
	 $\log_425=\log_{2^2}{5^2}=\dfrac{2}{2}\log_25=a$.}
\end{ex}

% nguồn: Ok_TracNghiemLop11_Lan1
\begin{ex}%[1D6H2-2]%[ma:e669c187]
	 Đặt $a=\log _2 7, b=\log _2 3$. Giá trị $\log _2\left(\dfrac{56}{9}\right)$ theo $a$ và $b$ là
	 \choice
		{\True $3+a-2  b$}
		 {$3+a-b^2$}
		 {$\dfrac{3 a}{2 b}$}
		{$\dfrac{3  a}{b^2}$}
	 \loigiai{
		Ta có $\log _2\left(\dfrac{56}{9}\right)=\log _2( 2^3\cdot 7)-\log _23^2=3+\log _2  7-2\log _2 3 =3+a-2b$.
	} 
\end{ex}

% nguồn: Ok_Toan11_TracNghiem_Lan2
\begin{ex}%[1D6V2-2]%[ma:6b846eff]
	[Mã 110 2017] Cho $x$, $y$ là các số thực lớn hơn $1$  thoả mãn $x^2+9y^2=6xy$. Tính $M=\dfrac{1+\log_{12}x+\log_{12}y}{2\log_{12}(x+3y)}$.
	\choice 
		 {$M=\dfrac{1}{2}$}
		{$M=\dfrac{1}{3}$} 
		 {$M=\dfrac{1}{4}$}
		{\True $M=1$} 
	 \loigiai{
		Ta có $x^2+9y^2=6xy\Leftrightarrow{(x-3y)^2}=0\Leftrightarrow  x=3y$.\\
		Khi  đó $M=\dfrac{1+\log_{12}x+\log_{12}y}{2\log_{12}(x+3y)}
		 =\dfrac{\log_{12}(12xy)}{\log_{12}(x+3y)^2}
		 =\dfrac{\log_{12}(36y^2)}{\log_{12}(36y^2)}=1$.
	 }
\end{ex}

% nguồn: Ok_Toan11_TracNghiem_Lan2
\begin{ex}%[1D6V2-2]%[ma:288869cc]
	Ba số $a+\log _2 3$, $a+\log _4 3; a+\log _8 3$ theo thứ tự lập thành cấp số nhân. Công bội của cấp  số nhân này bằng
	 \choice
		 {$\dfrac{1}{4}$}
		{$1$} 
		 {\True $\dfrac{1}{3}$}
		{$\dfrac{1}{2}$} 
	\loigiai{Theo giả thiết,  ta có
		\allowdisplaybreaks 
		\begin{eqnarray*}
			\left(a+\log _4 3\right)^2=\left(a+\log _2 3\right)\left(a+\log _8 3\right) &\Leftrightarrow& a \log _2 3+\left(\dfrac{1}{2} \log  _2 3\right)^2=\dfrac{4}{3} a \log _2 3+\dfrac{1}{3}\left(\log _2 3\right)^2 \\&\Leftrightarrow& \dfrac{1}{3} a \log _2 3=-\dfrac{1}{12}\left(\log _2 3\right)^2 \\&\Leftrightarrow& a=-\dfrac{1}{4} \log _2 3.
		\end{eqnarray*}
		Vậy $q=\dfrac{a+\log _4 3}{a+\log _2 3}=\dfrac{-\dfrac{1}{4} \log  _2 3+\dfrac{1}{2} \log _2 3}{-\dfrac{1}{4} \log _2 3+\log _2 3}=\dfrac{1}{3}$.
	} 
\end{ex}

% nguồn: Ok_Toan11_TracNghiem_Lan2
\begin{ex}%[1D6V2-2]%[ma:db837458]
	Với hai số thực dương $a$, $b$ tùy ý và $\dfrac{\log _3 5 \log _5 a}{1+\log _3 2}-\log _6 b=2$. Khẳng định  nào dưới đây là khẳng định đúng?
	 \choice
		{$a=b \log _6 2$} 
		{\True $a=36  b$}
		{$2  a+3 b=0$}
		{$a=b \log _6 3$} 
	\loigiai{ 
		Ta  có
		\allowdisplaybreaks 
		\begin{eqnarray*}
			\dfrac{\log _3 5 \log _5 a}{1+\log _3 2}-\log _6 b=2 &\Leftrightarrow& \dfrac{\log _3 a}{\log _3 6}-\log _6 b=2 \\&\Leftrightarrow& \log _6  a-\log _6 b=2
			\Leftrightarrow \log  _6 \dfrac{a}{b}=2 \\&\Leftrightarrow& \dfrac{a}{b}=36 \Leftrightarrow a=36 b.
		\end{eqnarray*}
	} 
\end{ex}

% nguồn: Ok_Toan11_TracNghiem_Lan2
\begin{ex}%[1D6V2-2]%[ma:41997516]
	Đặt $a = \ln 2$, $b  = \ln 5$ và
	$ 
	P = \ln \dfrac{2}{5}  + \ln \dfrac{5}{8} + \ln \dfrac{8}{11} + \ln \dfrac{11}{4} + \cdots + \ln \dfrac{7997}{8000}.
	$ 
	Khẳng định nào sau đây là đúng? 
	\choice 
		{  $P = 6a - 3b$ }
		{ $P =  5a - 3b$ }
		{  $P = 3a - 6b$ }
		{ \True $P = -5a - 3b$  }
	\loigiai{ 
		Ta có 
		 $$
		P = \ln \dfrac{2}{5} + \ln \dfrac{5}{8} + \ln  \dfrac{8}{11} + \cdots + \ln \dfrac{7997}{8000}
		= \ln\!\left(\dfrac{2}{5} \cdot \dfrac{5}{8} \cdot \dfrac{8}{11} \cdots \dfrac{7997}{8000}\right)  = \ln \dfrac{1}{4000}.
		$$ 
		Suy ra 
		 $$
		 P = -\ln(4000)
		 = -\ln(2^5 \cdot 5^3)
		= -5\ln 2 - 3\ln  5
		= -5a  - 3b.
		 $$
	 }
\end{ex}

% nguồn: Ok_Toan11_TracNghiem_Lan2
\begin{ex}%[1D6V2-2]%[ma:3703c080]
	Với mọi số thực dương $a$, $b$ thỏa mãn $\log_2 a + \log_4  b = 1$. Khẳng định nào sau đây là đúng?
	 \choice
		{ $a^2 b = 1$ } 
		{ $ab^2 = 4$ } 
		{ $ab^2 = 1$ } 
		 { \True $a^2 b = 4$ }
	\loigiai{ 
		Ta  có
		\begin{eqnarray*}
			&&\log_2 a + \log_4  b = 1\\
			&\Leftrightarrow& \log_2 a + \dfrac{1}{2} \log_2 b = 1\\ 
			&\Leftrightarrow& 2\log_2 a + \log_2  b = 2\\
			 &\Leftrightarrow& \log_2 (a^2 b) = 2\\
			 &\Leftrightarrow& a^2 b = 4.
		\end{eqnarray*}
	} 
\end{ex}

% nguồn: Ok_Toan11_TracNghiem_Lan2
\begin{ex}%[1D6V2-2]%[ma:8fe2ea8e]
	Cho $\log_{12} 3 = a$. Khẳng định nào  sau đây là đúng?
	\choice 
		{ \True $\log_{24} 18=\dfrac{3a+1}{3-a}$ } 
		{ $\log_{24} 18=\dfrac{3a+1}{3+a}$  }
		 { $\log_{24} 18=\dfrac{3a-1}{3+a}$ }
		{ $\log_{24} 18=\dfrac{3a-1}{3-a}$  }
	 \loigiai{
		Ta có 
		$$ 
		a =  \log_{12} 3
		= \dfrac{1}{\log_3  12}
		= \dfrac{1}{1 + 2\log_3  2}
		 \Leftrightarrow
		\log_2  3 = \dfrac{2a}{1-a}.
		$$ 
		Khi đó 
		\begin{eqnarray*}
			\log_{24} 18 &=& \dfrac{\log_2 \left(3^2 \cdot 2\right)}{\log_2  \left(2^3 \cdot 3\right)}\\
			&=& \dfrac{1 + 2\log_2 3}{3 + \log_2  3}\\
			&=& \dfrac{1 + \dfrac{4a}{1-a}}{3 + \dfrac{2a}{1-a}} \\ 
			&=& \dfrac{3a+1}{3-a}. 
		\end{eqnarray*}
	 }
\end{ex}

% nguồn: Ok_Toan11_TracNghiem_Lan2
\begin{ex}%[1D6V2-2]%[ma:5945ca9e]
	Cho $\log_{a}x=2$,  $\log_{b}x=3$ với $a$, $b$ là các số thực lớn hơn 1. Tính $P=\log_{\tfrac{a}{b^{2}}}x$.
	 \choice
		{$6$} 
		 {\True $-6$}
		 {$\dfrac{1}{6}$}
		{$-\dfrac{1}{6}$} 
	\loigiai{ 
		Vì $a$, $b$ là các số thực lớn hơn $1$ nên  ta có
		\begin{itemize}
			\item $\heva{&\log_{a}x=2\\ &\log_{b}x=3}\Leftrightarrow\heva{&x=a^{2}\\ &x=b^{3}}\Leftrightarrow a^{2}=b^{3}\Leftrightarrow a=\sqrt{b^{3}}\Leftrightarrow  a=b^{\tfrac{3}{2}}$.
			\item $P=\log_{\tfrac{a}{b^{2}}}x=\log_{\tfrac{b^{\frac{3}{2}}}{b^{2}}}x=\log_{b^{-\frac{1}{2}}}x=-2\log_{b}x=-6$. 
		\end{itemize}
	 }
\end{ex}

% nguồn: Ok_Toan11_TracNghiem_Lan2
\begin{ex}%[1D6V2-2]%[ma:9566c45b]
	Cho ba số thực dương $a$, $b$,  $c$ đều khác $1$ thoả mãn $\log_{a}b=2\log_{b}c=4\log_{c}a$ và $a+2b+3c=48$. Khi đó $P=abc$ bằng bao nhiêu?
	\choice 
		{\True $243$} 
		 {$521$}
		 {$512$}
		 {$324$}
	\loigiai{ 
		\begin{itemize}
			\item Do $a$, $b$, $c$ đều khác $1$ nên $\log_{a}b$, $\log_{b}c$, $\log_{c}a$ đều khác  $0$. Từ giả thiết $\log_{a}b=2\log_{b}c=4\log_{c}a$.
			\item Ta có 
			$\log_{a}b=2\log_{b}c\Leftrightarrow \log_{a}c \cdot  \log_{c}b=2\log_{b}c\Leftrightarrow \log_{a}c=2\log_{b}^{2}c$.
			\item $\log_{a}b=4\log_{c}a\Leftrightarrow \log_{a}c \cdot  \log_{c}b=4\log_{c}a\Leftrightarrow \log_{c}b=4\log_{c}^{2}a$.
			\item $1\log_{a}c \cdot \log_{c}b=8\log_{b}^{2}c \cdot \log_{c}^{2}a\Leftrightarrow \log_{a}b=8\log_{b}^{2}a\Leftrightarrow \log_{a}^{3}b=8\Leftrightarrow \log_{a}b=2\Leftrightarrow  b=a^{2}$.
			\item Mặt khác $\log_{a}b=2\log_{b}c\Leftrightarrow  \log_{a}b=2\log_{a^{2}}c\Leftrightarrow b=c$.
			\item Ta lại  có $a+2b+3c=48\Leftrightarrow a+2a^{2}+3a^{2}=48\Leftrightarrow5a^{2}+a-48=0\Leftrightarrow\hoac{&a=-\dfrac{16}{5}\\ &a=3.}$.
			\item Do $a$, $b$, $c$  đều là số thực dương $\Rightarrow a=3\Rightarrow b=9$, $c=9\Rightarrow P=abc=243$.
		\end{itemize}
	} 
\end{ex}

% nguồn: Ok_Toan11_TracNghiem_Lan2
\begin{ex}%[1D6V2-2]%[ma:7e990173]
	Cho $a$, $b>0$. Nếu $\ln x = 5\ln a + 2\ln \sqrt{b}$ thì $x$ bằng 
	\choice 
		{ $10a\sqrt{b}$ } 
		{ $\dfrac{a^5}{b}$ } 
		{  \True $a^5 b$ }
		{ $a^5 + b$ } 
	\loigiai{ 
		Ta  có
		\begin{eqnarray*}
			&&\ln  x = 5\ln a + 2\ln\sqrt{b}\\
			&\Leftrightarrow& \ln x = 5\ln a + \ln (\sqrt{b})^2\\ 
			&\Leftrightarrow& \ln x = \ln a^5  + \ln b\\
			&\Leftrightarrow& \ln x  = \ln(a^5 b)\\
			 &\Leftrightarrow& x = a^5 b.
		\end{eqnarray*}
	} 
\end{ex}

% nguồn: Ok_Toan11_TracNghiem_Lan2
\begin{ex}%[1D6V2-2]%[ma:889433b9]
	Cho $a = \log_{25}7$; $b = \log_2 5$. Tính $\log_{5}\dfrac{49}{8}$ theo  $a, b$.
	 \choice
		 {$\dfrac{4a-3}{b}$}
		{$\dfrac{4ab+3}{b}$} 
		{$\dfrac{5ab-3}{b}$} 
		{\True  $\dfrac{4ab-3}{b}$}
	 \loigiai{
		\begin{itemize}
			\item Ta có $a=\log_{25}7=\log_{5^{2}}7=\dfrac{1}{2}\log_{5}7\Rightarrow \log_{5}7=2a$ và  $b=\log_{2}5\Rightarrow \log_{5}2=\dfrac{1}{b}$.
			\item $\log_{5}\dfrac{49}{8}=\log_{5}49-\log_{5}8=\log_{5}7^{2}-\log_{5}2^{3}=2\log_{5}7-3\log_{5}2=2  \cdot 2a-3 \cdot \dfrac{1}{b}=\dfrac{4ab-3}{b}$.
		\end{itemize}
	 }
\end{ex}

% nguồn: Ok_Toan11_TracNghiem_Lan2
\begin{ex}%[1D6V2-2]%[ma:bcc896e2]
	Đặt $a=\log_2 3$, $b=\log_3 5$. Khẳng định  nào sau đây là đúng?
	\choice 
		{  $\log_{20} 12 =\dfrac{a+b}{b+2}$ }
		{ $\log_{20} 12  =\dfrac{ab+1}{b-2}$ }
		 { $\log_{20} 12 =\dfrac{a+1}{b-2}$ }
		{ \True $\log_{20} 12  =\dfrac{a+2}{ab+2}$ }
	 \loigiai{
		 Ta có
		$$ 
		\log_{20}  12
		=  \dfrac{\log_2 12}{\log_2 20}
		= \dfrac{\log_2 (4\cdot 3)}{\log_2  (4\cdot 5)}
		= \dfrac{2 + \log_2 3}{2 + \log_2  5}
		= \dfrac{2 +  a}{2 + a b}.
		 $$
	 }
\end{ex}

% nguồn: Ok_Toan11_TracNghiem_Lan2
\begin{ex}%[1D6V2-2]%[ma:82e354a7]
	Cho $x$, $y$ là các số thực lớn hơn $1$ thỏa mãn $x^{2}+9y^{2}=6xy$. Tính $M=\dfrac{1+\log_{12}x+\log_{12}y}{2\log_{12}\left(x+3y\right)}$. 
	\choice 
		{$M=\dfrac{1}{4}$} 
		 {$M=\dfrac{1}{2}$}
		{$M=\dfrac{1}{3}$} 
		{\True  $M=1$}
	 \loigiai{
		\begin{itemize}
			\item Ta  có $x^{2}+9y^{2}=6xy\Leftrightarrow\left(x-3y\right)^{2}=0\Leftrightarrow x=3y$.
			 \item Suy ra $M=\dfrac{1+\log_{12}\left(3y\right)+\log_{12}y}{2\log_{12}\left(6y\right)}=\dfrac{\log_{12}\left(36y^{2}\right)}{\log_{12}\left(36y^{2}\right)}=1$.
		\end{itemize}
	 }
\end{ex}

% nguồn: Ok_Toan11_TracNghiem_Lan2
\begin{ex}%[1D6V2-2]%[ma:b36909af]
	Cho  $\log_{2}x=\sqrt{5}$. Giá trị của biểu thức $P=\log_{2x}x$ bằng
	 \choice
		{$P=1+\sqrt{5}$} 
		{\True $P=\dfrac{\sqrt{5}}{\sqrt{5}+1}$} 
		{$P=\dfrac{5}{1+\sqrt{5}}$} 
		{$P=\dfrac{1}{\sqrt{5}}$} 
	 \loigiai{
		\begin{itemize}
			\item Ta có $\log_{2}x=\sqrt{5}\Leftrightarrow x=2^{\sqrt{5}}>1$ nên $\log_{x}2$, $\log_{x}x$, $\log_{x}\left(2x\right)$ đều xác định  và $\log_{x}\left(2x\right)\ne0$.
			\item  Cách 1. $P=\log_{2x}x=\dfrac{1}{\log_{x}\left(2x\right)}=\dfrac{1}{\log_{x}2+\log_{x}x}=\dfrac{1}{\dfrac{1}{\log_{2}x}+1}=\dfrac{1}{\dfrac{1}{\sqrt{5}}+1}=\dfrac{\sqrt{5}}{\sqrt{5}+1}$.
			\item Cách 2. $P=\log_{2x}x=\dfrac{\log_{2}x}{\log_{2}\left(2x\right)}=\dfrac{\log_{2}x}{\log_{2}2+\log_{2}x}=\dfrac{\sqrt{5}}{1+\sqrt{5}}$. 
		\end{itemize}
	 }
\end{ex}

% nguồn: Ok_Toan11_TracNghiem_Lan2
\begin{ex}%[1D6V2-2]%[ma:7385b3e5]
	Có bao nhiêu số thực dương $n\ne 1$ để $\log_n265$  là một số nguyên?
	 \choice
		{$2$} 
		 {$4$}
		{$6$} 
		{\True  $8$}
	\loigiai{ 
		 Ta có $\log_n256=\log_n{2^8}=8\log_n2=\dfrac{8}{\log_2n}$.\\
		 Để $\log_n265$ là một số nguyên thì
		\[\log_2n\in\{\pm 1;\pm 2;\pm 4;\pm  8\}\Leftrightarrow n\in\left\{\dfrac{1}{2};2;\dfrac{1}{4};4;\dfrac{1}{16};16;\dfrac{1}{256};256\right\}.\]
		 Vậy có tất cả 8 số thực dương $n\ne 1$ thỏa mãn điều kiện bài toán.
	} 
\end{ex}

% nguồn: Ok_Toan11_TracNghiem_Lan2
\begin{ex}%[1D6V2-2]%[ma:a89bf253]
	Cho $\log_{2}5=a$; $\log_{5}3=b$. Tính $\log_{5}24$  theo $a$ và $b$.
	\choice 
		{$\log_{5}24=\dfrac{3a+b}{b}$} 
		{$\log_{5}24=\dfrac{a+3b}{a}$} 
		{\True $\log_{5}24=\dfrac{3+ab}{a}$} 
		{$\log_{5}24=\dfrac{a+b}{3ab}$} 
	\loigiai{ 
		Ta có 
		\begin{eqnarray*}
			\log_{5}24&=&\log_{5}\left(8 \cdot 3\right)=\log_{5}8+\log_{5}3  =3\log_{5}2+\log_{5}3 =\dfrac{3}{\log_{2}5}+\log_{5}3 =\dfrac{3}{a}+b=\dfrac{3+ab}{a}.
		\end{eqnarray*}
	} 
\end{ex}

% nguồn: Ok_Toan11_TracNghiem_Lan2
\begin{ex}%[1D6V2-2]%[ma:0409ffa2]
	Giả sử $\log_{27} 5 =  a$, $\log_8 7 = b$, $\log_2 3 = c$. Khẳng định nào sau đây là đúng?
	 \choice
		{ \True $\log_{12} 35=\dfrac{3b + 3ac}{c+2}$  }
		{ $\log_{12} 35=\dfrac{3b + 3ac}{c+1}$  }
		{ $\log_{12} 35=\dfrac{3b  + 2ac}{c+3}$ }
		{ $\log_{12} 35=\dfrac{3b +  2ac}{c+2}$ }
	\loigiai{ 
		Ta có 
		$$\log_{27}  5 = a
		\Leftrightarrow \dfrac{1}{3}\log_3 5 = a 
		 \Leftrightarrow \dfrac{\log_2 5}{\log_2 3} = 3a
		\Leftrightarrow \log_2 5  = 3ac.$$
		Ta cũng có 
		$$\log_8  7
		\Leftrightarrow \dfrac{1}{3}\log_2 7 = b 
		 \Leftrightarrow \log_2 7 = 3b.$$
		 Khi đó
		 $$
		\log_{12}  35
		= \dfrac{\log_2 35}{\log_2  12}
		 = \dfrac{\log_2 5 + \log_2 7}{\log_2 3 + 2}
		= \dfrac{3ac +  3b}{c+2}.
		$$ 
	 }
\end{ex}

% nguồn: Ok_Toan11_TracNghiem_Lan2
\begin{ex}%[1D6V2-2]%[ma:f66c52dc]
	Cho $\log_2 5 = m$, $\log_3 5 = n$. Khẳng định nào sau đây là đúng? 
	\choice 
		{  $\log_6 5 = m^2 + n^2$ }
		{ \True $\log_6 5 =\dfrac{mn}{m+n}$ } 
		{  $\log_6 5 =m+n$ }
		{ $\log_6 5 =\dfrac{1}{m+n}$ } 
	\loigiai{ 
		 Ta có
		$$ 
		 \log_6 5
		= \dfrac{1}{\log_5  6}
		= \dfrac{1}{\log_5 (2\cdot 3)} 
		= \dfrac{1}{\log_5 2 + \log_5 3} 
		= \dfrac{1}{\dfrac{1}{m} +  \dfrac{1}{n}}
		= \dfrac{mn}{m+n}. 
		 $$
	 }
\end{ex}

% nguồn: Ok_Toan11_TracNghiem_Lan2
\begin{ex}%[1D6V2-2]%[ma:68e3a23b]
	Cho $a$, $b$ là các số thực dương và thỏa mãn $\dfrac{\log_3 a \cdot \log_2 3}{1 + \log_2 5} + \log b = 1$. Khẳng định nào sau đây là  đúng?
	\choice 
		{  $a = 1 - b \log_2 5$ }
		{ \True $ab = 10$  }
		{ $a \log_2 5 + b = 1$  }
		{ $a  + b = 1$ }
	 \loigiai{
		 Ta có
		\begin{eqnarray*}
			&&\dfrac{\log_3  a \cdot \log_2 3}{1+\log_2 5} + \log b = 1\\
			&\Leftrightarrow& \dfrac{\log_2 a}{\log_2 10} + \log  b = 1\\
			&\Leftrightarrow& \log a + \log b =  1\\
			&\Leftrightarrow& \log(ab) = 1\\ 
			&\Leftrightarrow&  ab = 10.
		\end{eqnarray*}
	 }
\end{ex}

% nguồn: Ok_Toan11_TracNghiem_Lan2
\begin{ex}%[1D6V2-2]%[ma:d8b7fb5c]
	Cho $a=\log_35$, $b=\log_27$, $c=\log_23$ và $$I=\dfrac{-1}{\log 126}\left(\log  \dfrac{1}{2}+\log \dfrac{2}{3}+\log \dfrac{3}{4}+\cdots+\log \dfrac{149}{150}\right).$$ Tính $I$ theo $a$, $b$, $c$.
	 \choice
		{\True  $\dfrac{1+c+2ac}{1+2c+b}$}
		{$\dfrac{1+c-2ac}{1+2c+b}$} 
		{$\dfrac{1+c-2ac}{1+2c-b}$} 
		 {$\dfrac{1+c+2ac}{1+2c-b}$}
	 \loigiai{
		Từ giả  thiết suy ra $\log_25=\log_23\cdot\log_35=ac$.\\
		Ta có \begin{align*}
		I&=\dfrac{-1}{\log  126}\cdot\log \left(\dfrac{1}{2}\cdot\dfrac{2}{3}\cdot\dfrac{3}{4}\cdots\dfrac{149}{150}\right)\\
		&=\dfrac{\log 150}{\log 126}\\ 
		 &=\log_{126}150=\dfrac{\log_2150}{\log_2126}\\
		 &=\dfrac{1+\log_23+2\log_25}{1+2\log_23+\log_27}\\
		&=\dfrac{1+c+2ac}{1+2c+b}. 
	\end{align*}
 }
\end{ex}

% nguồn: Ok_Toan11_TracNghiem_Lan2
\begin{ex}%[1D6V2-2]%[ma:633f996d]
	Cho số thực dương $a$, $b$ thỏa mãn $\log_{16} a  = \log_{20} b = \log_{25} \dfrac{2a-b}{3}$. Tỉ số $\dfrac{a}{b}$ thuộc khoảng nào sau đây?
	 \choice
		{ $\left(0;\dfrac12\right)$ } 
		{ $\left(\dfrac12;\dfrac23\right)$  }
		{ \True $(1;2)$ } 
		{ $(-2;0)$  }
	\loigiai{ 
		Đặt 
		$$ 
		\log_{16}  a = \log_{20} b = \log_{25} \dfrac{2a-b}{3} = x
		\Leftrightarrow 
		\heva{ 
		& a  = 16^x \\
		& b  = 20^x \\
		& 2a - b  = 3\cdot 25^x.
		} 
		$$ 
		 Khi đó, ta có
		\begin{eqnarray*}
			&&2\cdot 16^x - 20^x = 3\cdot 25^x\\ 
			 &\Leftrightarrow&
			2\cdot\left(\dfrac{16}{20}\right)^x - 3\cdot\left(\dfrac{25}{20}\right)^x  - 1 = 0\\
			&\Leftrightarrow& 
			 2\cdot\left(\dfrac45\right)^x - 3\cdot\left(\dfrac54\right)^x - 1 = 0\\
			&\Leftrightarrow& 
			2\cdot\left(\dfrac45\right)^{2x} -\left(\dfrac45\right)^x - 3 = 0.\\ 
			&\Leftrightarrow& 
			 \hoac{
			&\left(\dfrac{4}{5}\right)^x = -1 \text{ (loại)}\\ 
			&\left(\dfrac{4}{5}\right)^x =  \dfrac{3}{2}.}
		\end{eqnarray*}
		Vậy 
		 $$
		\dfrac{a}{b} =  \dfrac{16^x}{20^x}
		=  \left(\dfrac{4}{5}\right)^x
		=  \dfrac{3}{2}\in (1;2).
		$$ 
	} 
\end{ex}

% nguồn: Ok_Toan11_TracNghiem_Lan2
\begin{ex}%[1D6V2-2]%[ma:5e7dcada]
	Cho $\log_725=m$ và  $\log_25=n$. Tính $A=\log_{\sqrt[3]{5}}\dfrac{49}{8}$ theo $m$ và $n$
	\choice 
		{$\dfrac{9m+12n}{mn}$} 
		 {\True $\dfrac{-9m+12n}{mn}$}
		 {$-9m+12n+mn$}
		 {$\dfrac{-3m+4n}{3mn}$}
	 \loigiai{
		Ta có \begin{itemize}
		\item 
		 $\log_{\sqrt{5}}\dfrac{49}{8}=3\left(\log_549-\log_58\right)=3\left(2\log_57-3\log_52\right)$,
		\item $\log_725=2\log_75=\dfrac{2}{\log_57}=m\Rightarrow \log_57=\dfrac{2}{m}$, 
		\item $\log_25=n  \Leftrightarrow \dfrac{1}{\log_52}=n \Leftrightarrow \log\limits_52=\dfrac{1}{n}$.
	\end{itemize}
	Vậy $A=3\left(\dfrac{4}{m}-\dfrac{3}{n}\right)=\dfrac{12n-9m}{mn}$. 
 }
\end{ex}

% nguồn: Ok_Toan11_TracNghiem_Lan2
\begin{ex}%[1D6V2-2]%[ma:c98744f1]
	Cho $\log_3 5 =  a$, $\log_3 6 = b$, $\log_3 22 = c$ và $P = \log_3 \left(\dfrac{90}{11}\right)$. Khẳng định nào sau đây là đúng?
	\choice 
		{ $P  = 2a + b - c$ }
		{ \True $P = a + 2b - c$  }
		{ $P = 2a +  b + c$ }
		{ $P = 2a - b  + c$ }
	 \loigiai{
		 Ta có
		\begin{eqnarray*}
			P &=&  \log_3\left(\dfrac{90}{11}\right)\\
			&=& \log_3 180  - \log_3 22\\
			&=& \log_3 (36\cdot 5) - \log_3  22\\
			&=&  \log_3 36 + \log_3 5 - \log_3 22 \\
			&=& 2\log_3 6  + a - c \\
			 &=& a + 2b - c.
		\end{eqnarray*}
	} 
\end{ex}

% nguồn: Ok_Toan11_TracNghiem_Lan2
\begin{ex}%[1D6V2-2]%[ma:c9854e45]
	Cho $a=\log_7 5$; $b=\log_3 5$. Biểu thức $M=\log_{21} 5$ bằng 
	\choice 
		{  \True $\dfrac{ab}{a+b}$ }
		{ $ab$ } 
		{ $\dfrac{1}{ab}$  }
		{  $\dfrac{a+b}{ab}$ }
	\loigiai{ 
		Ta có 
		$$ 
		 M=\log_{21} 5
		=\dfrac{1}{\log_5  21}
		=\dfrac{1}{\log_5  3 + \log_5 7}
		=\dfrac{1}{\dfrac{1}{\log_3  5} + \dfrac{1}{\log_7 5}}
		=\dfrac{1}{\dfrac{1}{b} +  \dfrac{1}{a}}
		=\dfrac{1}{\dfrac{a+b}{ab}} 
		 =\dfrac{ab}{a+b}.
		 $$
	} 
\end{ex}


%%%%%%%%%%%%%%%% ĐÚNG/SAI %%%%%%%%%%%%%%%%

% nguồn: Ok_dungsai_lop11_Lan1
\begin{ex}%[1D6H2-2]%[ma:db661ebb]
 Cho  biểu thức $A=\log _2 x^2+\log _{\frac{1}{2}} x^3+\log _4 x$. Vậy
 \choiceTF 
  {\True Khi ${{\log }_{2}}x=1$  thì $A=-\dfrac{1}{2}$}
  {Khi ${{\log }_{2}}x=2$  thì $A=1$}
  {\True Khi ${{\log }_{2}}x=\sqrt{3}$ thì $A=-\dfrac{\sqrt{3}}{2}$} 
  {Khi ${\log _2 x=\sqrt{2}}$ thì  $A=\dfrac{\sqrt{2}}{2}$}
 \loigiai{ 
  Ta có:\\ 
  $\begin{aligned}
  A  & ={{\log }_{2}}{{x}^{2}}+{{\log }_{\frac{1}{2}}}{{x}^{3}}+{{\log }_{4}}x \\
  {} & =2{{\log  }_{2}}x+3{{\log }_{{{2}^{-1}}}}x+{{\log }_{{{2}^{2}}}}x=2{{\log }_{2}}x-3{{\log }_{2}}x+\dfrac{1}{2}{{\log }_{2}}x=-\dfrac{1}{2}{{\log }_{2}}x. \\
 \end{aligned}$
 }
\end{ex}

% nguồn: Ok_dungsai_lop11_Lan1
\begin{ex}%[1D6H2-2]%[ma:955968ee]
 Cho các biểu  thức sau $P=\log _2 8+\log _3 27-\log _5 5^3$; $Q=\ln (2 \mathrm{e})-\log 100$. Khi đó
 \choiceTF 
  {$P+Q=2 \ln 2$} 
  {\True $Q-P=\ln  2-4$}
  {\True $3  Q+P=3 \ln 2$}
   {\True $2 Q+P=2 \ln 2+1$}
  \loigiai{
  Ta có $ P=\log _2 8+\log _3 27-\log _5 5^3=\log _2 2^3+\log _3 3^3-\log _5 5^3=3+3-3=3$.\\ 
  Ta có $ 
   Q=\ln (2 \mathrm{e})-\log 100=\ln 2+\ln \mathrm{e}-\log 10^2=\ln 2+1-2=\ln 2-1$.\\
   Do đó
  \begin{itemize}
   \item  $P+Q=\ln 2+2$.
    \item $Q-P=\ln 2-4$.
   \item  $3Q+P=3\ln 2$.
   \item $2Q+P=2\ln  2+1$.
  \end{itemize}
 } 
\end{ex}

% nguồn: Ok_dungsai_lop11_Lan1
\begin{ex}%[1D6H2-2]%[ma:2e1d3638]
 Cho biểu thức $B=2\ln \sqrt{\mathrm{e}x}-\ln  \dfrac{\mathrm{e}^2}{\sqrt{x}}+\ln 3\cdot\log_3\left(\mathrm{e}x^2\right)$. Vậy
 \choiceTF{\True Cho $\ln  x=2$ thì $B=7$}{\True Cho $\ln x=4$ thì $B=14$}{\True Cho $\ln x=3$ thì $B=\dfrac{15}{2}$}{Cho $\ln x=6$ thì $B=18$}
 \loigiai{ 
   Ta có
  \begin{eqnarray*}
   B  & =&2 \ln (\mathrm{e}x)^{\tfrac{1}{2}}-\left(\ln \mathrm{e}^2-\ln x^{\tfrac{1}{2}}\right)+\ln \left(\mathrm{e}x^2\right) \\
   & =&\ln (\mathrm{e}x)-\left(2-\dfrac{1}{2} \ln x\right)+\ln \left(\mathrm{e}x^2\right)=(\ln \mathrm{e}+\ln x)-2+\dfrac{1}{2} \ln  x+\ln\mathrm{e}+\ln x^2 \\
    & =&1+\ln x-2+\dfrac{1}{2} \ln x+1+2\ln x=\dfrac{7}{2}\ln x.
  \end{eqnarray*}
 } 
\end{ex}

% nguồn: Ok_dungsai_lop11_Lan1
\begin{ex}%[1D6H2-2]%[ma:f6c0d320]
 Cho $a$, $b$, $c>1$ và $m$,  $n\in \mathbb{R}$.
  \choiceTF
  {\True  $\log_a \sqrt{a\sqrt{a}}=\dfrac{3}{4}$}
  {$\log_a  b^2\cdot \log _{\sqrt{b}} c\cdot \log _{c^2} a^3=\dfrac{1}{6}$}
   {\True Cho $\log 3=m,\log 7=n$. Khi đó $\log_3 70=\dfrac{n+1}{m}$}
  {Cho $\log_5 2=m,\log_5 3=n$. Khi đó $\log  _{250} 30=\left(m+n+1\right)\left(3+m\right)$}
  \loigiai{
  \begin{itemchoice}
   \itemch  \textbf{Đúng}.\\
   Ta  có
   $\log_a \sqrt{a\sqrt{a}}=\log_a \sqrt{a^{\tfrac{3}{2}}}=\log_a  a^{\tfrac{3}{4}}=\dfrac{3}{4}$.
    \itemch \textbf{Sai}.\\
   Ta có 
   $\log_a  b^2\cdot \log _{\sqrt{b}} c\cdot \log _{c^2} a^3=2\cdot \log_a b\cdot 2\cdot \log_b c\cdot \dfrac{3}{2}\cdot \log_c a=6$.
   \itemch \textbf{Đúng}.\\ 
    Cho $\log 3=m$, $\log 7=n$.\\
   Ta có $\log_3  70=\dfrac{\log 70}{\log 3}=\dfrac{\log 7+1}{\log 3}=\dfrac{n+1}{m}$.
    \itemch \textbf{Sai}.\\
   Cho $\log_5  2=m$, $\log_5 3=n$.\\
   Khi đó 
   $\log _{250} 30=\dfrac{\log_5 \left(3\cdot 2\cdot  5\right)}{\log_5 \left(5^3\cdot 2\right)}=\dfrac{m+n+1}{3+m}$.
  \end{itemchoice}
 } 
\end{ex}

% nguồn: Ok_dungsai_lop11_Lan1
\begin{ex}%[1D6H2-2]%[ma:336f22fa]
 Xét tính đúng, sai trong các khẳng định sau 
 \choiceTF 
   {\True $\log 8-\log 27=3\log \dfrac{2}{3}$}
   {$\log _{16} 3\cdot \log_9 64\cdot \log_4 27=\dfrac{8}{9}$}
  {\True $\log _{\sqrt{2}} 5+\log_4 \dfrac{1}{\sqrt{25}}+\log _{\tfrac{1}{2}} 125=\dfrac{-3}{2}  \log_2 5$}
  {\True  $3\log_3 \sqrt[3]7+\log _{\sqrt[3]3} 21=4\log_3 7+3$}
 \loigiai{ 
  \begin{itemchoice}
    \itemch \textbf{Đúng}.\\
   Ta có $\log 8-\log 27=\log 2^3-\log  3^3=2\log 2-3\log 3=3\log \dfrac{2}{3}$.
   \itemch \textbf{Sai}.\\ 
   Ta có $\log _{16} 3\cdot \log_9 64\cdot \log_4 27=\dfrac{1}{4} \log_2 3\cdot \dfrac{6}{2} \log_3 2\cdot \dfrac{3}{2}  \log_2 3=\dfrac{9}{8} \log_2 3$.
   \itemch  \textbf{Đúng}.\\
   Ta có 
   \begin{eqnarray*}
    \log _{\sqrt{2}} 5+\log_4 \dfrac{1}{\sqrt{25}}+\log _{\tfrac{1}{2}}  125
    &=&2\log_2  5+\dfrac{1}{2}\cdot \left(-1\right)\log_2 5-3\log_2 5\\
    &=&-\dfrac{3}{2} \log_2  5.
   \end{eqnarray*}
   \itemch \textbf{Đúng}.\\ 
   Ta có $3\log_3 \sqrt[3]7+\log _{\sqrt[3]3} 21=\log_3 7+3\left(\log_3 7+1\right)=4\log_3  7+3$.
  \end{itemchoice}
  }
\end{ex}

% nguồn: Ok_dungsai_lop11_Lan1
\begin{ex}%[1D6H2-2]%[ma:122b62b1]
 Cho hai số thực $a$  và $b$, với $1< a < b$.
 \choiceTF 
  {$\log_a b  < 1< \log_b a$}
  {$1< \log_a b < \log_b a$} 
  {$\log_b a < \log_a b <  1$}
  {\True  $\log_b a < 1< \log_a b$}
  \loigiai{
  Ta có $b>a>1\Leftrightarrow \heva{&{\log_a b > \log_a a\Leftrightarrow \log_a b > 1} \\&{\log_b  b > \log_b a\Leftrightarrow 1> \log_b a}
  } \Leftrightarrow \log_b a < 1< \log_a  b$.
  \begin{itemchoice}
    \itemch \textbf{Sai}.\\
   Ta có $\log_a  b < 1< \log_b a$.
    \itemch \textbf{Sai}.\\
   Ta có $1< \log_a b <  \log_b a$.
   \itemch  \textbf{Sai}.
   $\log_b a < \log_a b < 1$. 
    \itemch \textbf{Đúng}.\\
   Ta có $\log_b a <  1< \log_a b$.
  \end{itemchoice}
 } 
\end{ex}

% nguồn: Ok_dungsai_lop11_Lan1
\begin{ex}%[1D6H2-2]%[ma:3a4d6adf]
 Cho $a$, $b$, $c  > 0$, $a\ne 1$ và $m$, $n\in \mathbb{R}$.
  \choiceTF
  {$\log_a 1=1$} 
  {\True $a^{\log_a  b}=b$}
  {$\log_a a^m=\dfrac{1}{m}$} 
  {\True $\log_a  b^m c^n=m\log_a b+n\log_a c$}
  \loigiai{
  \begin{itemchoice}
    \itemch \textbf{Sai}.\\
   Ta có $\log_a  1=1$.
    \itemch \textbf{Đúng}.\\
   Ta có  $a^{\log_a b}=b$.
   \itemch  \textbf{Sai}.\\
    Ta có $\log_a a^m=m\log_a a=m$.
    \itemch \textbf{Đúng}.\\
   Ta có $\log_a b^m c^n=\log_a b^m+\log_a c^n=m\log_a  b+n\log_a c$.
 \end{itemchoice}}
\end{ex}

% nguồn: Ok_dungsai_lop11_Lan1
\begin{ex}%[1D6H2-2]%[ma:1eb10822]
 Cho các biểu thức sau $P = \dfrac{\log_a\left(a^3b^2\right) - \log_b\left(\dfrac{b^3}{a^2}\right)}{\log_a^2 b + 1}$ và $Q = \log_a b^3 + \log_{a^2}b^6$ với $a$, $b$ và các số dương khác $1$.  Các khẳng định sau là đúng hay sai?
  \choiceTF
  {\True $Q =  6\log_a b$}
  {$P =  6\log_b a$}
  {$Q  = 3P$}
  {\True $PQ =  12$}
  \loigiai{
 \begin{itemize}
  \item Ta có $Q = 3\log_a b + 6\cdot\dfrac{1}{2}\log_a b  = 6\log_a b$.
  \item Ta  có
  \[\begin{aligned}
  P &= \dfrac{\log_a\left(a^3b^2\right) - \log_b\left(\dfrac{b^3}{a^2}\right)}{\log_a^2 b +  1} \\
  & = \dfrac{\log_a a^3 + \log_a b^2 -  \left(\log_b b^3 - \log_b a^2\right)}{\log_a^2 b + 1}\\
  & = \dfrac{3 + 2\log_a b -  3 + 2\log_b a}{\log_a^2 b + 1}\\
  & =  \dfrac{2\left(\log_a b + \dfrac{1}{\log_a b}\right)}{\log_a^2 b + 1}\\
  & = \dfrac{2\cdot\dfrac{\log_a^2 b + 1}{\log_a  b}}{\log_a^2 b + 1}\\
  &  = \dfrac{2}{\log_a b} = 2\log_b a.
 \end{aligned}\]
 \item $Q=6\log_a b=\dfrac{6}{\log_b  a}=\dfrac{12}{2\log_b a}=\dfrac{12}{P}$.
 \item $PQ = 6\log_a b\cdot 2\log_b a =  12$.
\end{itemize}
 }
\end{ex}

% nguồn: Ok_dungsai_lop11_Lan2
\begin{ex}%[1D6H2-2]%[ma:859a2f7e]
	Cho $a$, $b$, $c>1$  và $m$, $n \in \mathbb{R}$.
	\choiceTF 
		 {\True $\log_a\sqrt{a\sqrt{a}}=\dfrac{3}{4}$}
		{$\log_a b^2\cdot  \log_{\sqrt{b}}c\cdot \log_{c^2}a^3=\dfrac{1}{6}$}
		{\True Cho $\log 3=m$, $\log 7=n$. Khi  đó $\log_3 70=\dfrac{n+1}{m}$}
		{Cho  $\log_5 2=m$, $\log_5 3=n$. Khi đó $\log_{250}30=(m+n+1)(3+m)$}
	\loigiai{ 
	\begin{itemchoice}
		\itemch Ta có  $\log_a\sqrt{a\sqrt{a}}=\log_a \sqrt{a\cdot a^{\frac{1}{2}}}=\log_a \sqrt{a^{\frac{3}{2}}}=\log_a a^{\frac{3}{2}\cdot \frac{1}{2}}=\log_a a^{\frac{3}{4}}=\dfrac{3}{4}$.
		\itemch  Ta có
		$$\begin{aligned}
		\log_a b^2\cdot \log_{\sqrt{b}}c\cdot \log_{c^2}a^3&=2\log_a b\cdot  2\log_b c\cdot 3\log_{c^2}a\\
		 &=4\log_a b\cdot \log_b c\cdot \dfrac{3}{2}\log_c a\\
		&=6\cdot \log_a c\cdot \log_c a=6. 
	\end{aligned}$$
	\itemch Ta  có $\log_3 70=\dfrac{\log 70}{\log 3}=\dfrac{ \log (7\cdot 10)}{\log 3}=\dfrac{\log 7+\log 10}{\log 3}=\dfrac{\log 7+1}{\log 3}=\dfrac{n+1}{m}$.
	 \itemch Ta có $\log_{250}30=\dfrac{\log_5 250}{\log_5 30}=\dfrac{\log_5 (3\cdot 2\cdot 5)}{\log_5(5^3\cdot 2)}=\dfrac{\log_5 3+\log_5 2+\log_5 5}{\log_5 5^3+\log_52}=\dfrac{m+n+1}{3+m}$.
\end{itemchoice}
} 
\end{ex}

% nguồn: Ok_dungsai_lop11_Lan2
\begin{ex}%[1D6H2-2]%[ma:a3b466c5]
	Cho các biểu  thức sau: $P=\dfrac{\log_a (a^3b^2)-\log_b \left(\dfrac{b^3}{a^2} \right)}{\log_a^2b+1}$ và $Q=\log_a b^3+\log _{a^2} b^6$ với $a$, $b$ là các số dương và $a$ khác 1.
	 \choiceTF
		 {\True $Q=6\log_a b$}
		{$P=6\log_b  a$}
		{$Q=3P$} 
		{\True $Q\cdot P=12$} 
	\loigiai{ 
		Ta có $Q=3\log_a  b+6\cdot \dfrac{1}{2} \log_a b=6\log_a b$.\\
		Ta có $$\begin{aligned}
		P & =\dfrac{\log _a a^3+\log _a b^2-\left(\log  _b b^3-\log _b a^2\right)}{\log _a^2 b+1} \\
		 & =\dfrac{3+2 \log _a b-3+2 \log _b a}{\log _a^2 b+1}\\
		&=\dfrac{2\left(\log _a  b+\dfrac{1}{\log _a b}\right)}{\log _a^2 b+1} \\
		& =\dfrac{2\left(\dfrac{\log _a^2 b+1}{\log _a b}\right)}{\log  _a^2 b+1}\\
		 &=\dfrac{2}{\log _a b}\\
		&=2 \log _b  a .
	\end{aligned}$$
	\begin{itemchoice}
		\itemch  Vậy $Q=6\log_a b$.
		 \itemch Vậy $P=2\log_b a$.
		\itemch Vậy  $Q\neq 3P$.
		\itemch Vậy  $Q\cdot P=12$.
	\end{itemchoice}
} 
\end{ex}

% nguồn: Ok_dungsai_lop11_Lan2
\begin{ex}%[1D6H2-2]%[ma:23bcb280]
	Cho biểu thức $Q=2^{\log _{16} x^4+\log_2 x^2}$  với $x$ là số thực khác $0$.
	\choiceTF 
		{\True $Q>  0$}
		{\True Khi $x=2$ thì  $Q=8$}
		{Khi  $x=-2$ thì $Q=-8$}
		{Khi  $x=3$ thì $Q=9$}
	 \loigiai{
		 Ta có $Q=2^{\log_2 x^4+\log_2 x^2}=2^{4\cdot \tfrac{1}{4} \log_2|x|+2\log_2|x|}=2^{3\log_2|x|}=\left(2^{\log_2|x|} \right)^3=|x|^3$.
		\begin{itemchoice}
			\itemch  Vậy $Q> 0$.
			\itemch Khi $x=2$ thì $Q=2^3=8$. 
			\itemch Khi $x=-2$  thì $Q=|-2|^3=8$.
			\itemch  Khi $x=3$ thì $Q=3^3=27$.
		\end{itemchoice}
	 }
\end{ex}

% nguồn: Ok_dungsai_lop11_Lan2
\begin{ex}%[1D6H2-2]%[ma:8d9496be]
	Cho $a$, $b$ là các  số thực dương lớn hơn $1$ thỏa mãn $\log_a b=2$.
	\choiceTF 
		{$2a=b$} 
		{\True $a^2=b$} 
		{\True  $\log_b a=\dfrac{1}{2}$}
		{\True $\log _{ab^2} b^5=2$} 
	 \loigiai{
		\begin{itemchoice}
			\itemch  Với $a$, $b$ là các số thực dương lớn hơn $1$ và thỏa mãn $\log_a b=2\Leftrightarrow a^2=b$.
			\itemch Với $a$, $b$ là các số thực dương lớn hơn $1$ và thỏa  mãn $\log_a b=2\Leftrightarrow a^2=b$.
			 \itemch Ta có $\log_a b=2\Rightarrow \log_b a=\dfrac{1}{\log_a b}=\dfrac{1}{2}$.
			\itemch \[\log _{ab^2} b^5=\dfrac{1}{\log  _{b^5} ab^2}=\dfrac{1}{\log _{b^5} a+\log _{b^5} b^2}=\dfrac{1}{\dfrac{1}{5} \log_b a+\dfrac{2}{5} \log_b b}=\dfrac{1}{\dfrac{1}{5}\cdot\dfrac{1}{2}+\dfrac{2}{5}}=2.
			\] 
		\end{itemchoice}
	 }
\end{ex}

% nguồn: Ok_dungsai_lop11_Lan1
\begin{ex}%[1D6V2-2]%[ma:8dd02147]
 Cho hàm số $f(x)=\log_a(x+1)$ với $a=2024!$. Các mệnh  đề sau đúng hay sai?
  \choiceTF
   {$f(0)=1$}
  {Hàm  số xác định trên $\mathscr{D}=(0;+\infty)$}
  {$f(x)<1\Leftrightarrow  x<2024!-1$}
  {\True $f(1)+f(2)+\cdots +f(2023)=1$} 
 \loigiai{ 
  \begin{itemchoice}
   \itemch Sai. $f(0)=\log_a1=0$ nên  mệnh đề sai.
   \itemch  Sai. Hàm số xác định khi và chỉ khi $x+1>0\Leftrightarrow x >-1$.
   \itemch  Sai. $f(x)<1\Leftrightarrow\log_a(x+1)<1\Leftrightarrow 0<x+1<a\Leftrightarrow-1<x<2024!-1$.
   \itemch Đúng. \allowdisplaybreaks\begin{eqnarray*}
   &&f(1)+f(2)+\cdots +f(2023)\\ 
   &=&\log_a2+\log_a3+\cdots +\log_a2024\\ 
   &=&\log_a(2\cdot 3\cdots  2024)=\log_aa=1.
  \end{eqnarray*}
\end{itemchoice}}
\end{ex}

% nguồn: Ok_dungsai_lop11_Lan1
\begin{ex}%[1D6V2-2]%[ma:369cc583]
 Cho hàm số $y = x + \dfrac{4}{x -  3}$ có đồ thị $(C)$ và đường thẳng $d\colon y = -3x + 2$. Khi đó
 \choiceTF 
  {Điểm $N(-1;0)$ thuộc đồ thị $(C)$} 
  {Đồ thị $(C)$ và đường thẳng $d$ cắt nhau tại hai điểm có tích hoành độ giao  điểm bằng $-\dfrac{5}{2}$}
  {\True Tiếp  tuyến của đồ thị $(C)$ tại điểm có hoành độ bằng $1$ song song với trục hoành}
  {\True $M_1(2;-2),M_2(4;8)$ là hai điểm thuộc đồ thị $(C)$ mà tiếp tuyến của $(C)$ tại đó song song với đường thẳng $d$} 
 \loigiai{ 
  \begin{itemchoice}
   \itemch  Sai.\\
   Với $x=-1$  thì $y=-2$ nên điểm $N(-1;0)$ không thuộc đồ thị $(C)$.
   \itemch Sai\\ 
   Phương trình hoành độ giao điểm của $(C)$ và $(d)$  là
   \begin{eqnarray*}
     && x + \dfrac{4}{x - 3}=-3x+2\\
    &\Leftrightarrow &  \heva{&x\ne 3\\&4x-2+\dfrac{4}{x-3}=0}\\
    &\Leftrightarrow  & \heva{& x\ne 3\\& (4x-2)(x-3)+4=0}\\
    &\Leftrightarrow & \heva{&x\ne 3\\&2x^2-7x+5=0.} 
   \end{eqnarray*}
   Phương trình này có $\Delta=9>0$ nên có hai nghiệm. Theo định lý Viète tích hai  nghiệm bằng $\dfrac{5}{2}$.
   \itemch Đúng.\\ 
   Với $x=1$ thì $y=-1$ nên  điểm $K(1;-1)$ không thuộc trục $Ox$.\\
   Lại có $y'=1-\dfrac{4}{(x-3)^2}$ nên  $y'(1)=0$. Do đó tiếp tuyến của $(C)$ tại $K(1;-1)$ song song với trục hoành.
   \itemch Đúng.\\ 
   Đặt $f(x)=x +  \dfrac{4}{x - 3}$.
   Ta có $f'(x) = 1 - \dfrac{4}{\left(x - 3 \right)^2}$. \\ 
   Gọi $M\left(x_0; f\left(x_0 \right)  \right)$ là tọa độ tiếp điểm. \\
   Hệ số góc của tiếp tuyến tại điểm $\left(x_0; f\left(x_0 \right) \right)$ là $y'(x_0)  = \dfrac{x^2_0 - 6x_0 + 5}{\left(x_0 - 3 \right)^2}$. \\
   Vì tiếp tuyến song song với đường thẳng $d:y = -3x +  2 $ nên
   $$y'(x_0) = - 3 \Leftrightarrow \heva{&4x^2_0 - 24x_0 + 32 = 0 \\ &x_0 \neq 3}  \Leftrightarrow \hoac{&x_0 = 4 \\ &x_0 = 2.}$$
   Với $x_0=2$, suy ra  tiếp điểm $M_1(2;-2)$ và $f'(2)=-3$. Phương trình tiếp tuyến tại $M_1$ là
   \[y+2=-3(x-2)\Leftrightarrow  y=-3x+4.\]
    Với $x_0=4$, suy ra tiếp điểm $M_2(4;8)$ và $f'(4)=-3$. Phương trình tiếp tuyến tại $M_1$ là
    \[y-8=-3(x-4)\Leftrightarrow y=-3x+20.\]
   Cả hai tiếp  tuyến đều song song với đường thẳng $d$. Suy ra có hai tiếp điểm thỏa yêu cầu bài $M_1(2;-2),M_2(4;8)$.
  \end{itemchoice}
  }
\end{ex}

% nguồn: Ok_dungsai_lop11_Lan2
\begin{ex}%[1D6V2-2]%[ma:044fd241]
	Cho các biểu thức sau: $A=\left(a^3\sqrt{a} \right)^{\log_a b}+\left(\sqrt[3]{b^2} \right)^{\log_b a}$ với $\heva{a,b > 0 \\ a\ne 1,b\ne 1}$ và $B=\log \dfrac{a}{b}+\log \dfrac{b}{c}+\log \dfrac{c}{d}-\log \dfrac{a}{d}$ với $a$, $b$, $c$,  $d$ là các số dương.
	\choiceTF 
		{$A=\sqrt[3]a+\sqrt{b^4}$} 
		{$B=\dfrac{a}{b}$} 
		{\True  $A+B\sqrt{a}=\sqrt[3]{a^2}+\sqrt{b^7}$}
		{$A-B\sqrt{b}=2\sqrt[3]{a^2}+\sqrt{b^7}$} 
	 \loigiai{
		\begin{itemchoice}
			\itemch $$\begin{aligned}
			A & =\left(a^3\cdot a^{\tfrac{1}{2}} \right)^{\log_a b}+\left(b^{\tfrac{2}{3}} \right)^{\log_b  a} \\
			& =\left(a^{\tfrac{7}{2}} \right)^{\log_a b}+\left(b^{\tfrac{2}{3}}  \right)^{\log_b a}\\
			&= \left(a^{\log_a b} \right)^{\tfrac{7}{2}}+\left(b^{\log_b  a} \right)^{\tfrac{2}{3}}\\
			&  =b^{\tfrac{7}{2}}+a^{\tfrac{2}{3}}\\
			&=\sqrt[3]{a^2}+\sqrt{b^7}. 
		\end{aligned}$$
		\itemch Ta có $B=\log \left(\dfrac{a}{b} \cdot \dfrac{b}{c}  \cdot \dfrac{c}{d} \right)-\log \dfrac{a}{d}=\log \left(\dfrac{a}{d}:\dfrac{a}{d} \right)=\log 1=0$.
		 \itemch Vậy $A+B\sqrt{a}=\sqrt[3]{a^2}+\sqrt{b^7}$.
		 \itemch Vậy $A-B\sqrt{b}=\sqrt[3]{a^2}+\sqrt{b^7}$.
	\end{itemchoice}
 }
\end{ex}


%%%%%%%%%%%%%%%% TRẢ LỜI NGẮN %%%%%%%%%%%%%%%%

% nguồn: Ok_TraLoiNgan_Lop11_Lan1
\begin{ex}%[1D6H2-2]%[ma:fe37fa4a]
	Cho  $\log_a b=2$ và $\log_a c=3$. Tính $Q=\log_a (b^2c^3)$.
	 \shortans{13}
	\loigiai{ 
		Ta có  $Q=\log_a (b^2c^3)=\log_a b^2+\log_a c^3=2\log_a b+3\log_a c=2\cdot 2+3\cdot 3=13$.
	 }
\end{ex}

% nguồn: Ok_TraLoiNgan_Lop11_Lan1
\begin{ex}%[1D6H2-2]%[ma:6971d545]
	Với $a$, $b$ là các số  thực dương và $a$ khác $1$, ta rút gọn biểu thức $Q=\sqrt{\log_a^2(ab)-\dfrac{2\ln b}{\ln a}-1}$ thu được kết quả $m\left|\log_a b^n \right|$ với $m$; $n$ là các số tự nhiên. Tính giá trị $S=m+2n$.
	\shortans{3} 
	 \loigiai{
		 Ta có
		\begin{align*}
			Q=&\sqrt{\left(\log_a a+\log_a b\right)^2-2\log_a b-1}\\ 
			 =&\sqrt{\left(1+\log_a b\right)^2-2\log_a b-1}\\
			 =&\sqrt{1+2\log_a b+\log_a^2b-2\log_a b-1}\\
			=&\sqrt{\log_a^2b}=\left|\log_a  b\right|.
		\end{align*}
		Khi đó $\heva{&m=1\\&n=1} \Rightarrow  S=m+2n=3$.
	} 
\end{ex}

% nguồn: Ok_TraLoiNgan_Lop11_Lan1
\begin{ex}%[1D6H2-2]%[ma:650b17e9]
	Cho $a$, $b$, $c$ là các số thực khác $0$ thỏa mãn $4^a=25^b=10^c$.  Tính $T=\dfrac{c}{a}+\dfrac{c}{b}$.
	 \shortans{2}
	\loigiai{ 
		Giả sử $4^a=25^b=10^c=t\Rightarrow \heva{&{a=\log_4t}\\&{b=\log _{25} t} \\&{c=\log  _{10} t}.}$\\
		Ta  có
		\begin{align*}
			T=&\dfrac{c}{a}+\dfrac{c}{b}=\dfrac{\log _{10} t}{\log_4 t}+\dfrac{\log _{10} t}{\log  _{25} t}\\
			=&\dfrac{\log_t 4}{\log_t 10}+\dfrac{\log_t  25}{\log_t 10}\\
			=&\log  _{10} 4+\log _{10} 25\\
			 =&\log _{10} \left(4\cdot 25\right)\\
			=&\log  _{10} 100=2.
		\end{align*}
	 }
\end{ex}

% nguồn: Ok_TraLoiNgan_Lop11_Lan1
\begin{ex}%[1D6H2-2]%[ma:53b01aee]
	Cho $a=\log_2 5$, $b=\log_3 5$. Biểu diễn $\log_6 5$ theo $a$ và $b$  ta thu được kết quả dạng $\dfrac{a}{m+\dfrac{n\cdot a}{b}}$ với $m;n$ là các số tự nhiên. Tính giá trị $S=m-2n$.
	\shortans{-1} 
	 \loigiai{
		Ta có $\log_6 5=\dfrac{\log_2 5}{\log_2 6}=\dfrac{a}{\log_2 2+\log_2 3}=\dfrac{a}{1+\log_2 5\log_5  3}=\dfrac{a}{1+\dfrac{a}{b}}.\\
		\Rightarrow \heva{&{m=1}  \\ &{n=1}.}\\ \Rightarrow S=m-2n=-1$.
	 }
\end{ex}

% nguồn: Ok_TraLoiNgan_Lop11_Lan1
\begin{ex}%[1D6H2-2]%[ma:c84c9737]
	Cho số thực $a$  thỏa mãn $0< a\ne 1$. Tính giá trị của biểu thức $T=\log_a \left(\dfrac{a^2\cdot \sqrt[3]{a^2} \cdot \sqrt[5]{a^4}}{\sqrt[{15}]{a^7}} \right)$.
	 \shortans{3}
	 \loigiai{
		Ta có $T=\log_a \left(\dfrac{a^2\cdot \sqrt[3]{a^2} \cdot \sqrt[5]{a^4}}{\sqrt[{15}]{a^7}}  \right)=\log_a \dfrac{a^{2+\tfrac{2}{3}+\tfrac{4}{5}}}{a^{\tfrac{7}{15}}}=\log_a a^{2+\tfrac{2}{3}+\tfrac{4}{5}-\tfrac{7}{15}}=\log_a a^3=3$.
	 }
\end{ex}

% nguồn: Ok_TraLoiNgan_Lop11_Lan1
\begin{ex}%[1D6H2-2]%[ma:9307d664]
	Cho $a=\log 2$, $b=\ln 2$. Biểu diễn $\ln  800$ theo $a$ và $b$ ta thu được kết quả dạng $m\cdot b+\dfrac{n\cdot b}{a}$ với $m$; $n$ là các số tự nhiên. Tính giá trị $P=m\cdot n$.
	 \shortans{6}
	 \loigiai{
		Ta có 
		\begin{align*}
			\ln  800=&\ln \left(2^3\cdot 10^2\right)=3\ln 2+2\ln 10\\
			=&3\ln  2+2\ln 2\cdot \log_2 10\\
			 =&3\ln 2+\dfrac{2\ln 2}{\log 2}=3b+\dfrac{2b}{a}.
		\end{align*}
		$\Rightarrow\heva{&{m=3}  \\
		&{n=2}}  \Rightarrow P=6$.
	} 
\end{ex}

% nguồn: Ok_TraLoiNgan_Lop11_Lan2
\begin{ex}%[1D6H2-2]%[ma:6ca93a2a]
	Cho ba số thực dương $a$, $b$, $c$ đều khác $1$ thoả mãn  $\log_ab=2\log_bc=4\log_ca$ và $a+2b+3c=48$. Giá trị của $P=abc$ bằng bao nhiêu?
	\shortans{243} 
	\loigiai{ 
		Do $a,b,c$ đều khác $1$ nên  $\log_ab$, $\log_bc$, $\log_ca$ đều khác $0$. Ta có
		\begin{itemize}
			\item 
			$\log_ab=2\log_bc\Leftrightarrow\log_ac\cdot\log_cb=2\log_bc\Leftrightarrow\log_ac=2\log_b^2c$. 
			 \item $\log_ab=4\log_ca\Leftrightarrow\log_ac\cdot\log_cb=4\log_ca\Leftrightarrow\log_cb=4\log_c^2a$.
		\end{itemize}
		 Suy ra $\log_ac\cdot\log_cb=8\log_b^2c\cdot\log_c^2a\Leftrightarrow\log_ab=8\log_b^2a\Leftrightarrow\log_a^3b=8\Leftrightarrow\log_ab=2\Leftrightarrow b=a^2$.\\
		Mà  $\log_ab=2\log_bc\Leftrightarrow\log_ab=2\log_{a^2}c\Leftrightarrow b=c$.\\
		Ta lại có  $a+2b+3c=48\Leftrightarrow a+2a^2+3a^2=48\Leftrightarrow 5a^2+a-48=0\Leftrightarrow\hoac{&a=-\dfrac{16}{5}\\&a=3.}$\\
	Do $a$, $b$, $c$ đều là số  thực dương $\Rightarrow a=3\Rightarrow b=9$, $c=9\Rightarrow P=abc=243$.}
\end{ex}

% nguồn: Ok_TraLoiNgan_Lop11_Lan2
\begin{ex}%[1D6H2-2]%[ma:8b3ce600]
	Cho $\log_23=a$  và $\log_35=b$. Biết $\log_{12}150=\dfrac{2ab+ma+n}{a+2}$ với $m$, $n\in\mathbb{Z}$. Giá trị của biểu thức $L=m+n$ bằng bao nhiêu?
	\shortans{2} 
	 \loigiai{
		 Ta có $\log_23=a\Leftrightarrow\log_32=\dfrac{1}{a}$ và $\log_35=b$.\\
		Khi đó $\log_{12}150=\dfrac{\log_3150}{\log_312}=\dfrac{\log_3({2\cdot 3\cdot 5}^2)}{\log_3(2^2\cdot  3)}=\dfrac{\log_32+1+2\log_35}{2\log_32+1}=\dfrac{\dfrac{1}{a}+1+2b}{\dfrac{2}{a}+1}=\dfrac{2ab+a+1}{a+2}$.\\
	Do $m$, $n\in\mathbb{Z}\Rightarrow m=1$ và $n=1\Rightarrow L=m+n=2$.} 
\end{ex}

% nguồn: Ok_TraLoiNgan_Lop11_Lan2
\begin{ex}%[1D6H2-2]%[ma:1ca9d7ca]
	Cho $a,b$ là các số thực dương thỏa mãn $\log_ab=\dfrac{b}{8}$, $\log_2a=\dfrac{16}{b}$. Giá trị của biểu  thức $G=a+b$ bằng bao nhiêu?
	 \shortans{20}
	 \loigiai{
		Ta có $\log_ab=\dfrac{b}{8}$ và  $\log_2a=\dfrac{16}{b}\Rightarrow\log_2a\cdot\log_ab=\dfrac{16}{b}\cdot\dfrac{b}{8}\Leftrightarrow\log_2b=2\Leftrightarrow b=4$ \\
	$ \Rightarrow\log_2a=4\Leftrightarrow a=16\Rightarrow G=a+b=16+4=20  $.}
\end{ex}

% nguồn: Ok_TraLoiNgan_Lop11_Lan2
\begin{ex}%[1D6H2-2]%[ma:dfb8acfd]
	Cho $a$, $b$ là  các số thực dương thỏa mãn $\log_ab=\dfrac{b}{8}$, $\log_2a=\dfrac{16}{b}$. Giá trị của biểu thức $G=a+b$ là
	\shortans{20} 
	 \loigiai{
		 Ta có $\log_ab=\dfrac{b}{8}$ và $\log_2a=\dfrac{16}{b}$.
		\begin{align*}
			 &\bullet\log_ab=\dfrac{b}{8} \Leftrightarrow \dfrac{\log_2b}{\log_2a} =\dfrac{b}{8} \Leftrightarrow \log_2b =\dfrac{b}{8}\cdot\log_2a \Leftrightarrow \log_2b =2 \Leftrightarrow b=4.\\
			&\bullet \log_2a=\dfrac{16}{b} =4 \Leftrightarrow a=16\\ 
			&\bullet G=a+b=16+4=20. 
		\end{align*}
	 }
\end{ex}

% nguồn: Ok_TraLoiNgan_Lop11_Lan2
\begin{ex}%[1D6H2-2]%[ma:a61cc18c]
	Cho $a$, $b$, $c$ là các số  thực dương thỏa $a^{\log_37}=9$, $b^{\log_711}=7$, $c^{\log_{11}25}=11$. Giá trị của biểu thức bằng bao nhiêu? $T=a^{\log_3^27}+b^{\log_7^211}+c^{\log_{11}^225}$.
	\shortans{85} 
	\loigiai{ 
		Ta có 
		\begin{eqnarray*}
			T&=&a^{\log_3^27}+b^{\log_7^211}+c^{\log_{11}^225}\\ 
			&=&\left(a^{\log_37}\right)^{\log_37}+\left(b^{\log_711}\right)^{\log_711}+\left(c^{\log_{11}25}\right)^{\log_{11}25} \\ 
			 &=&(9)^{\log_37}+(7)^{\log_711}+(11)^{\log_{11}25}\\
			 &=&7^2+11+25\\
			&=&85. 
		\end{eqnarray*}
	 }
\end{ex}

% nguồn: Ok_TraLoiNgan_Lop11_Lan2
\begin{ex}%[1D6H2-2]%[ma:1798cba0]
	Cho $x$, $y$, $z$ là các số thực lớn hơn $1$ và gọi $w$ là số thực dương sao cho $\log_xw=12$, $\log_yw=20$ và $\log_{xyz}w=6$. Giá  trị của $\log_zw$ bằng bao nhiêu?
	 \shortans{30}
	 \loigiai{
		\begin{itemize}
			 \item $\log_xw=12\Rightarrow\log_wx=\dfrac{1}{12}$.
			 \item $\log_yw=20\Rightarrow\log_wy=\dfrac{1}{20}$.
		\end{itemize}
		\begin{eqnarray*}
			\log_{xyz}w=6&\Leftrightarrow&\dfrac{1}{\log_{_w}(xyz)}=6\\ 
			 &\Leftrightarrow&\dfrac{1}{\log_{_w}x+\log_{_w}y+\log_{_w}z}=6\\
			 &\Leftrightarrow&\dfrac{1}{\log_{_w}x+\log_{_w}y+\log_{_w}z}=6\\
			 &\Leftrightarrow&\dfrac{1}{\dfrac{1}{12}+\dfrac{1}{20}+\log_{_w}z}=6\\
			 &\Leftrightarrow&\log_{_w}z=\dfrac{1}{30}\\
			&\Rightarrow&\log_zw=30. 
		\end{eqnarray*}
	} 
\end{ex}

% nguồn: Ok_TraLoiNgan_Lop11_Lan2
\begin{ex}%[1D6H2-2]%[ma:7dce405b]
	Cho các số thực dương $x$, $y$, $z$ thỏa mãn đồng thời $\dfrac{1}{\log_2x}+\dfrac{1}{\log_2y}+\dfrac{1}{\log_2z}=\dfrac{1}{10}$ và $\log_2(xyz)=10$. Giá trị của biểu thức  $\log_2\left[xyz(x+y+z)-xy-yz-zx+1\right]$ bằng bao nhiêu?
	 \shortans{20}
	\loigiai{ 
		Đặt $a=\log_2x$; $b=\log_2y$; $c=\log_2z$. Ta có $\dfrac{1}{a}+\dfrac{1}{b}+\dfrac{1}{c}=\dfrac{1}{10}$ và  $a+b+c=10$.
		\begin{eqnarray*}
			&&\left(\dfrac{1}{a}+\dfrac{1}{b}+\dfrac{1}{c}\right)(a+b+c)=1\Leftrightarrow(a+b+c)(ab+ac+bc)=abc\\ 
			&\Leftrightarrow& a^2b+ab^2+abc+abc+b^2c+bc^2+a^2c+ac^2=0\Leftrightarrow(a+b)(b+c)(c+a)=0. 
		\end{eqnarray*}
		Vì vai trò $a$, $b$, $c$ như  nhau nên giả sử $a+b=0\Rightarrow c=10\Rightarrow z=2^{10}$ và $xy=1$.
		\begin{eqnarray*}
			\log_2\left[xyz(x+y+z)-xy-yz-zx+1\right]&=&\log_2\left[z(x+y+z)-1-yz-zx+1\right]\\ 
			&=&\log_2(z^2)=2\log_2z=2\cdot  10=20.
		\end{eqnarray*}
	 }
\end{ex}

% nguồn: Ok_TraLoiNgan_Lop11_Lan1
\begin{ex}%[1D6V2-2]%[ma:15d4b948]
	Cho  các số thực dương $a$, $b$ thỏa mãn $\log_{16}a=\log_{20}b=\log_{25}\dfrac{2a-b}{3}$. Tính tỉ số $\dfrac{a}{b}$ (Kết quả viết dưới dạng thập phân).
	 \shortans{1,5.}
	 \loigiai{
		Đặt  $\log_{16}a=\log_{20}b=\log_{25}\dfrac{2a-b}{3}=t$.\\
		Suy ra $\heva{&a=16^t\\&b=20^t\\&\dfrac{2a-b}{3}=25^t}\Rightarrow 2\cdot16^t-20^t=3\cdot25^t\Leftrightarrow  2\cdot\left(\dfrac{16}{25}\right)^t-\left(\dfrac{4}{5}\right)^t-3=0\Leftrightarrow \hoac{&\left(\dfrac{4}{5}\right)^t=\dfrac{3}{2}\\&\left(\dfrac{4}{5}\right)^t=-1 \text{ (vô lí).}}$\\
		Vậy  $\dfrac{a}{b}=\dfrac{16^t}{20^t}=\left(\dfrac{4}{5}\right)^t=\dfrac{3}{2}=1{,}5$.
	} 
\end{ex}

% nguồn: Ok_TraLoiNgan_Lop11_Lan2
\begin{ex}%[1D6V2-2]%[ma:35c82e25]
	Cho $\log_8|x|+\log_4y^2=5$  và $\log_8|y|+\log_4x^2=7$. Tìm giá trị của biểu thức $P=|x|-|y|$
	 \shortans{56}
	 \loigiai{
		Điều kiên  $x, y \ne 0$.\\
		Cộng vế với vế của hai phương trình,  ta được
		$$\log_8|xy|+\log_4x^2y^2=12  \Leftrightarrow \log_2|xy|=9\Leftrightarrow |xy|=512.$$
		Trừ vế với vế của hai  phương trình, ta được
		 $$\log_8\left|\dfrac{x}{y}\right|+\log_4\dfrac{y^2}{x^2}=-2 \Leftrightarrow \log_2\left|\dfrac{x}{y}\right|=3 \Leftrightarrow \left|\dfrac{x}{y}\right|=8 \Leftrightarrow |x|=8|y|.$$
		 Suy ra $|y|=8 \Rightarrow |x|=64 \Leftrightarrow P=56$.
	} 
\end{ex}

% nguồn: Ok_TraLoiNgan_Lop11_Lan2
\begin{ex}%[1D6V2-2]%[ma:b38a4b5c]
	Biết rằng $7^x=5$  và $5^y=7$. Không dùng máy tính cầm tay, hãy tính giá trị của $xy$
	\shortans{1} 
	 \loigiai{
		Ta có $\heva{&7^x=5\\  &5^y=7} \Leftrightarrow \heva{&x=\log_75\\ &y=\log_57} \Rightarrow xy=\log_75\cdot\log_57=1$.
	} 
\end{ex}

% nguồn: Ok_TraLoiNgan_Lop11_Lan2
\begin{ex}%[1D6V2-2]%[ma:1dd4cf7d]
	 Cho các số thực $a$, $b$ thỏa mãn $a>b>1$ và $\dfrac{1}{\log_ba}+\dfrac{1}{\log_ab}=\sqrt{2020}$. Giá trị của biểu thức $P=\dfrac{1}{\log_{ab}b}-\dfrac{1}{\log_{ab}a}$ bằng (kết quả làm tròn đến hàng phần mười)
	\shortans{44,9} 
	 \loigiai{
		Do $a>b>1$ nên $\log_ab>0$, $\log_ba>0$  và $\log_ba>\log_ab$.\\
		Ta có 
		\begin{align*}
			&\dfrac{1}{\log_ba}+\dfrac{1}{\log_ab}=\sqrt{2020}\\ 
			&\Leftrightarrow \log_ba+\log_ab=\sqrt{2020}\\ 
			&\Leftrightarrow \log_b^2a+\log_a^2b+2=2020\\ 
			 &\Leftrightarrow \log_b^2a+\log_a^2b=2018.
		\end{align*}
		Khi đó  $P=\log_bab-\log_aab=\log_ba+\log_bb-\log_aa-\log_ab=\log_ba-\log_ab$.\\
		Suy ra $P^2={{(\log_ba-\log_ab)}^2}=\log_b^2a+\log_a^2b-2=2018-2=2016 \Rightarrow  P=\sqrt{2016}$.
	 }
\end{ex}

% nguồn: Ok_TraLoiNgan_Lop11_Lan2
\begin{ex}%[1D6V2-2]%[ma:ced98727]
	Cho các số thực dương  $x,y,z$ thỏa mãn đồng thời $\dfrac{1}{\log_2x} +\dfrac{1}{\log_2y} +\dfrac{1}{\log_2z} =\dfrac{1}{2020}$ và $\log_2(xyz)=2020$. Tính $\log_2(xyz(x+y+z)-xy-yz-zx+1)$
	\shortans{4040} 
	 \loigiai{
		Đặt $a=\log_2x$; $b=\log_2y$;  $c=\log_2z$.\\
		Ta có $\dfrac{1}{a}+\dfrac{1}{b}+\dfrac{1}{c}=\dfrac{1}{2020}$  và $a+b+c=2020$.
		\begin{eqnarray*}
			&& \left(\dfrac{1}{a}+\dfrac{1}{b}+\dfrac{1}{c}\right)(a+b+c)=1\Leftrightarrow  (a+b+c)(ab+ac+bc)=abc \\
			& \Leftrightarrow &  a^2b+ab^2+abc+abc+b^2c+bc^2+a^2c+ac^2=0 \\
			& \Leftrightarrow  & (a+b)(b+c)(c+a)=0.
		\end{eqnarray*}
		Vì vai trò $a$, $b$, $c$ như nhau nên giả sử $a+b=0 \Rightarrow c=2020 \Rightarrow z=2^{2020}$  và $xy=1$.
		\begin{eqnarray*}
			\log_2(xyz(x+y+z)-xy-yz-zx+1)  &=& \log_2(z(x+y+z)-1-yz-zx+1) \\
			& = &  \log_2(z^2)=2\log_2z=4040.
		\end{eqnarray*}
	 }
\end{ex}

% nguồn: Ok_TraLoiNgan_Lop11_Lan2
\begin{ex}%[1D6V2-2]%[1D6V2-3]%[ma:3b01ea26]
	Cho $2$ số thực dương $a$, $b$ thỏa mãn $\sqrt{a}\ne b$, $a\ne 1$, $\log_ab=2$. Tính $T={\log_{\tfrac{\sqrt{a}}{b}}}\sqrt[3]{ba}$. (kết  quả làm tròn đến hàng phần chục)
	\shortans{-0,7} 
	\loigiai{ 
		Ta  có
		\begin{align*}
			&  \log_ab=2\Rightarrow \log_ba=\dfrac{1}{2}.\\
			 & T={\log_{\tfrac{\sqrt{a}}{b}}}\sqrt[3]{ba} ={\log_{\tfrac{\sqrt{a}}{b}}}\sqrt[3]{b} +\log_{\tfrac{\sqrt{a}}{b}}\sqrt[3]{a} =\dfrac{1}{{\log_{\sqrt[3]{b}}}\tfrac{\sqrt{a}}{b}} +\dfrac{1}{{\log_{\sqrt[3]{a}}}\tfrac{\sqrt{a}}{b}}\\
			& =\dfrac{1}{{\log_{\sqrt[3]{b}}}\sqrt{a}-{\log_{\sqrt[3]{b}}}b} +\dfrac{1}{{\log_{\sqrt[3]{a}}}\sqrt{a}-{\log_{\sqrt[3]{a}}}b} =\dfrac{1}{\tfrac{3}{2}\log_ba-3}+\dfrac{1}{\tfrac{3}{2}-3\log_ab}\\ 
			& =\dfrac{1}{\dfrac{3}{2}\cdot\dfrac{1}{2}-3}+\dfrac{1}{\dfrac{3}{2}-3\cdot  2} =-\dfrac{2}{3}.
		\end{align*}
	 }
\end{ex}

% nguồn: Ok_TraLoiNgan_Lop11_Lan2
\begin{ex}%[1D6V2-2]%[ma:921f6f34]
	Cho $\log_ax=2$, $\log_bx=3$ với $a$, $b$ là các số thực lớn hơn $1$.  Tính $P=\log_{\tfrac{a}{b^2}}x$
	\shortans{-6} 
	 \loigiai{
		Vì $a$, $b$ là các số thực lớn hơn $1$  nên ta có
		\begin{align*}
			&\heva{&\log_ax=2\\ &\log_bx=3} \Leftrightarrow \heva{&x=a^2\\ &x=b^3} \Leftrightarrow  a^2=b^3 \Leftrightarrow a=\sqrt{b^3} \Leftrightarrow a=b^{\tfrac{3}{2}}.\\
			&P=\log_{\tfrac{a}{b^2}}x =\log_{\tfrac{b^{\frac{3}{2}}}{b^2}}x =\log_{b^{\tfrac{-1}{2}}}x =-2\log_bx=-2\cdot3=-6. 
		\end{align*}
	} 
\end{ex}

% nguồn: Ok_TraLoiNgan_Lop11_Lan2
\begin{ex}%[1D6V2-2]%[ma:06b05bf4]
	Cho $a$, $b$ là  các số thực dương thỏa mãn $\log_3a+\log_4b^2=5$, $\log_3a^3-\log_4b=1$. Giá trị của biểu thức $H=a+b$ là
	 \shortans{19}
	\loigiai{ 
		Ta có 
		\begin{align*}
			\heva{&\log_3a+\log_4b^2=5\\ &\log_3a^3-\log_4b=1} \Leftrightarrow  \heva{&\log_3a+2\log_4b=5\\ &3\log_3a-\log_4b=1} \Leftrightarrow \heva{&\log_3a=1\\ &\log_4b=2} \Leftrightarrow \heva{&a=3\\ &b=16.}\\
		\end{align*}
		 Suy ra $H=a+b=3+16=19$.
	} 
\end{ex}

% nguồn: Ok_TraLoiNgan_Lop11_Lan2
\begin{ex}%[1D6V2-2]%[ma:5148caa2]
	Cho hàm số $f(x)=\dfrac{1}{2}\log_2\left(\dfrac{2x}{1-x}\right)$  và hai số thực $m$, $n$ thuộc khoảng $(0;1)$ sao cho $m+n=1$. Tính $f(m)+f(n)$
	\shortans{1} 
	\loigiai{ 
		Ta có 
		\begin{align*}
			 f(m)+f(n)&=\dfrac{1}{2}\log_2\left(\dfrac{2m}{1-m}\right)+\dfrac{1}{2}\log_2\left(\dfrac{2n}{1-n}\right)\\
			&=\dfrac{1}{2}\left[\log_2\left(\dfrac{2m}{1-m}\right)+\log_2\left(\dfrac{2n}{1-n}\right)\right]\\ 
			&=\dfrac{1}{2}\log_2\left(\dfrac{2m}{1-m}\cdot\dfrac{2n}{1-n}\right)\\ 
			&=\dfrac{1}{2}\log_2\left(\dfrac{4mn}{1-m-n+mn}\right),  \text{ vì } m+n=1\\
			 &=\dfrac{1}{2}\log_2\left(\dfrac{4mn}{mn}\right)=\dfrac{1}{2}\log_24 =\dfrac{1}{2}\cdot2=1.
		\end{align*}
	 }
\end{ex}

% nguồn: Ok_TraLoiNgan_Lop11_Lan2
\begin{ex}%[1D6V2-2]%[ma:b0a06f17]
	Cho $a$, $b$, $c$  là các số thực dương thỏa $a^{\log_37}=9$, $b^{\log_711}=7$, $c^{\log_{11}25}=11$. Tính giá trị biểu thức $T=a^{\log_3^27}+b^{\log_7^211}+c^{\log_{11}^225}$.
	\shortans{85} 
	 \loigiai{
		Ta  có
		\begin{align*}
			T&=a^{\log_3^27}+b^{\log_7^211}+c^{\log_{11}^225}\\ 
			 &=\left(a^{\log_37}\right)^{\log_37} +\left(b^{\log_711}\right)^{\log_711} +\left(c^{\log_{11}25}\right)^{\log_{11}25}\\
			&=(9)^{\log_37}+(7)^{\log_711}+(11)^{\log_{11}25}\\ 
			 &=7^2+11+25=85.
		\end{align*}
	 }
\end{ex}

% nguồn: Ok_TraLoiNgan_Lop11_Lan2
\begin{ex}%[1D6V2-2]%[ma:1de5f3d7]
	Cho hàm số $f(x)$ xác định trên $\mathbb{R}$ và thỏa mãn $f(1)=1$, $f(m+n)=f(m)+f(n)+mn$ với mọi $m$, $n\in\mathbb{N}^*$. Giá trị của biểu thức $T=\log\left[\dfrac{f(96)-f(69)-241}{2}\right]$ bằng bao  nhiêu?
	 \shortans{3}
	\loigiai{ 
		 Có $f(1)=1$, $f(m+n)=f(m)+f(n)+mn$.\\
		 Suy ra
		\begin{eqnarray*}
			 f(96)&=&f(95+1)=f(95)+f(1)+95=f(95)+96=f(94)+95+96\\
			 &=&\cdots \\
			&=&f(1)+2+\cdots  +95+96\\
			&\Rightarrow&  f(96)=1+2+\cdots +95+96=\dfrac{96\cdot 97}{2}=4\,656.
		\end{eqnarray*}
		Tương tự $f(69)=1+2+\cdots +68+69=\dfrac{69\cdot 70}{2}=2\,415$.\\ 
	 Vậy $T=\log\left[\dfrac{f(96)-f(69)-241}{2}\right]=\log\left(\dfrac{4\,656-2\,415-241}{2}\right)=\log 1\,000=3$.}
\end{ex}

% nguồn: Ok_TraLoiNgan_Lop11_Lan2
\begin{ex}%[1D6V2-2]%[ma:09e3091a]
	Tính giá trị của biểu thức $P=\log_{a^2}(a^{10}b^2) +\log_{\sqrt{a}}\left(\dfrac{a}{\sqrt{b}}\right) +\log_{\sqrt[3]{b}}b^{-2}$ 
	\shortans{1} 
	\loigiai{ 
		Sử dụng các quy tắc biến đổi  logarit\\
		\begin{eqnarray*} \allowdisplaybreaks
			P&  =&\log_{a^2}(a^{10}b^2) +\log_{\sqrt{a}}\left(\dfrac{a}{\sqrt{b}}\right) +\log_{\sqrt[3]{b}}b^{-2} \\
			 & =&\dfrac{1}{2}\left[\log_aa^{10}+\log_ab^2\right] +2\left[\log_aa-\log_a\sqrt{b}\right]+3\cdot(-2)\log_bb \\
			& =&\dfrac{1}{2}\left[10+2\log_ab\right]+2\left[1-\dfrac{1}{2}\log_ab\right]-6=1. 
		\end{eqnarray*}
	} 
\end{ex}

% nguồn: Ok_TraLoiNgan_Lop11_Lan2
\begin{ex}%[1D6V2-2]%[ma:dba08909]
	Cho $a$, $b$ là các số thực  dương lớn hơn 1 thỏa mãn $\log_ab=2$. Tính giá trị biểu thức $P=\log_{a^2}b+\log_{ab^2}b^5$
	\shortans{3} 
	\loigiai{ 
		 Ta có
		\begin{align*}
			P&=\log_{a^2}b+\log_{ab^2}b^5 =\log_{a^2}b+\dfrac{1}{\log_{b^5}ab^2}  =\dfrac{1}{2}\log_ab+\dfrac{1}{\log_{b^5}a+\log_{b^5}b^2}\\
			 &=\dfrac{1}{2}\log_ab+\dfrac{1}{\tfrac{1}{5}\log_ba+\tfrac{2}{5}} =\dfrac{1}{2}\cdot2 +\dfrac{1}{\tfrac{1}{5}\cdot\tfrac{1}{2}+\tfrac{2}{5}}=1+2=3.
		\end{align*}
	 }
\end{ex}

% nguồn: Ok_TraLoiNgan_Lop11_Lan2
\begin{ex}%[1D6V2-2]%[ma:b06e30eb]
	Cho $x$, $y$ và $z$ là các số thực lớn hơn $1$ và gọi ${w}$ là  số thực dương sao cho $\log_xw=24$, $\log_yw=40$ và $\log_{xyz}w=12$. Tính $\log_zw$
	\shortans{60} 
	\loigiai{ 
		\begin{align*}
			&\log_xw=24 \Rightarrow  \log_wx=\dfrac{1}{24}\\
			&\log_yw=40 \Rightarrow \log_wy=\dfrac{1}{40}. 
		\end{align*}
		Lại  do
		\begin{align*}
			 \log_{xyz}w=12 &\Leftrightarrow \dfrac{1}{\log_{_w}(xyz)}=12\\
			&  \Leftrightarrow \dfrac{1}{\log_{_w}x+\log_{_w}y+\log_{_w}z}=12\\
			&  \Leftrightarrow \dfrac{1}{\log_{_w}x+\log_{_w}y+\log_{_w}z}=12\\
			 & \Leftrightarrow \dfrac{1}{\dfrac{1}{24}+\dfrac{1}{40}+\log_{_w}z}=12\\
			&  \Leftrightarrow \log_{_w}z=\dfrac{1}{60}\\
			&  \Rightarrow \log_zw=60.
		\end{align*}
	} 
\end{ex}

% nguồn: Ok_TraLoiNgan_Lop11_Lan2
\begin{ex}%[1D6C2-2]%[ma:59d08e00]
	Cho $x$, $y$, $z$ là ba số thực dương lập thành cấp số nhân; $\log_ax$, $\log_{\sqrt{a}}y$, $\log_{\sqrt[3]{a}}z$ lập thành cấp số cộng, với $a$ là số thực dương  khác $1$. Giá trị của $p=\dfrac{9x}{y}+\dfrac{y}{z}+\dfrac{3{z}}{x}$ là
	 \shortans{13}
	\loigiai{ 
		$x$, $y$, $z$ là ba số thực dương lập thành  cấp số nhân nên ta có $xz=y^2$.\\
		$\log_ax$, $\log_{\sqrt{a}}y$, $\log_{\sqrt[3]{a}}z$ lập thành cấp số  cộng nên:\\
		$\log_ax+\log_{\sqrt[3]{a}}z=2\log_{\sqrt{a}}y$  $\Leftrightarrow \log_ax+3\log_az=4\log_ay$ $\Leftrightarrow x{z^3}=y^4$.\\
		Từ và  ta suy ra $x=y=z$.\\
		Vậy $p=\dfrac{9x}{y}+\dfrac{y}{z}+\dfrac{3z}{x}=9+1+3=13$. 
	} 
\end{ex}

% nguồn: Ok_TraLoiNgan_Lop11_Lan1
\begin{ex}%[1D6G2-2]%[ma:0a5cf1d5]
	Cho $\log_2 x=\sqrt{2}$. Tính giá trị biểu thức $P=\log_2 x^2+\log _{\tfrac{1}{2}} x^3+\log_4 x$. (Làm tròn đến hàng  hàng phần mười).
	\shortans{-0,7} 
	 \loigiai{
		Ta có $P=2\log_2 x-3\log_2 x+\dfrac{1}{2} \log_2 x=-\dfrac{1}{2} \log_2 x=-\dfrac{1}{2}\cdot \sqrt{2}=-\dfrac{\sqrt{2}}{2}  \approx-0{,}7$.
	 }
\end{ex}

% nguồn: Ok_TraLoiNgan_Lop11_Lan2
\begin{ex}%[1D6K2-2]%[ma:60425a20]
	Có bao nhiêu số nguyên dương $n$ để $\log_{n} 256$ là một  số nguyên dương?
	 \shortans{4}
	\loigiai{ 
		$\log_{n} 256 = 8.\log_{n} 2 = \dfrac{8}{\log_{2} n}$ là  số nguyên dương $\Leftrightarrow \log_{2} n \in \{1;2;4;8\} \Leftrightarrow n \in \{2;4;16;256\}$.\\
	 Vậy, có $4$ số nguyên dương.}
\end{ex}


%%%%%%%%%%%%%%%% TỰ LUẬN %%%%%%%%%%%%%%%%

% nguồn: Ok_TuanLuan_Toan11_Lan1
\begin{bt}%[1D6N2-2]%[ma:36491e91]
	Đặt $a=\log_23$ và $b=\log_53$. Hãy biểu diễn $\log_6{25}$ theo $a$ và  $b$.
	 \loigiai{$\log_6{45}=\dfrac{\log_3{\left(5\cdot 3^2\right)}}{\log_3\left(2\cdot3\right)}=\dfrac{\log_35+2}{\log_32+1}=\dfrac{\dfrac{1}{b}+2}{\dfrac{1}{a}+1}=\dfrac{a+2ab}{ab+b}$.}
\end{bt}

% nguồn: Ok_TuanLuan_Toan11_Lan1
\begin{bt}%[1D6N2-2]%[ma:25491146]
	 Cho $\log_2 6=a$; $\log_2 7=b$. Hãy biểu diễn $\log_{18}{42}$ theo $a$ và $b$.
	\loigiai{ 
		 Ta có $\log_{18}{42}=\dfrac{\log_2{42}}{\log_2{18}}=\dfrac{\log_2\left(6\cdot 7\right)}{\log_2\left(\dfrac{36}{2}\right)}=\dfrac{\log_26+\log_27}{\log_26^2-\log_22}=\dfrac{a+b}{2a-1}$.
	 }
\end{bt}

% nguồn: Ok_TuanLuan_Toan11_Lan2
\begin{ex}%[1D6N2-2]%[ma:aebc13b7]
	Đặt $a=\log_23$ và $b=\log_53$. Hãy biểu diễn $\log_6{25}$ theo $a$  và $b$.
	\loigiai{$\log_6{45}=\dfrac{\log_3{\left(5\cdot  3^2\right)}}{\log_3\left(2\cdot3\right)}=\dfrac{\log_35+2}{\log_32+1}=\dfrac{\dfrac{1}{b}+2}{\dfrac{1}{a}+1}=\dfrac{a+2ab}{ab+b}$.}
\end{ex}

% nguồn: Ok_TuanLuan_Toan11_Lan2
\begin{ex}%[1D6N2-2]%[ma:a760a9f7]
	Cho $\log_2 6=a$; $\log_2 7=b$.  Hãy biểu diễn $\log_{18}{42}$ theo $a$ và $b$.
	\loigiai{ 
		Ta có  $\log_{18}{42}=\dfrac{\log_2{42}}{\log_2{18}}=\dfrac{\log_2\left(6\cdot 7\right)}{\log_2\left(\dfrac{36}{2}\right)}=\dfrac{\log_26+\log_27}{\log_26^2-\log_22}=\dfrac{a+b}{2a-1}$.
	 }
\end{ex}

% nguồn: Ok_TuanLuan_Toan11_Lan1
\begin{bt}%[1D6H2-2]%[ma:8e31fdb6]
	Cho $\log_a b=3$ và $\log _a c=4$ với $a$;  $b$; $c > 0$; $a \neq 1$. Tính giá trị của $P=\log _a\left(\dfrac{a^2\cdot \sqrt{b}}{c^3}\right)$.
	 \loigiai{
		Xét  biểu thức
		\begin{eqnarray*}
			P &=& \log _a\left(\dfrac{a^2  \cdot \sqrt{b}}{c^3}\right) \\
			&=& \log _a  a^2+\log _a \sqrt{b}-\log _a c^3 \\
			&=& 2+\log _a b^{\tfrac{1}{2}}-3 \log _a c \\ 
			&=&  2+\dfrac{1}{2} \log _a b-3 \log _a c \\
			&=& 2+\dfrac{3}{2}-12  \\
			&=&  - \dfrac{17}{2}.
		\end{eqnarray*}
	 }
\end{bt}

% nguồn: Ok_TuanLuan_Toan11_Lan1
\begin{bt}%[1D6H2-2]%[ma:eed010a5]
	Cho  $\log_{a} x=2$, $\log_{b} x=3$. Giá trị của $P=\log_{\tfrac{a^2}{b}} x$ bằng bao nhiêu?
	 \loigiai{
		Ta có $\log_{\tfrac{a^2}{b}}x=\dfrac{1}{\log_{x}{\tfrac{a^2}{b}}}=\dfrac{1}{\log_{x} a^2+\log_{x} b}=\dfrac{1}{2\log_{x} a+\log_{x} b}=\dfrac{1}{2\cdot  \dfrac{1}{2}+\dfrac{1}{3}}=\dfrac{3}{4}$.
	} 
\end{bt}

% nguồn: Ok_TuanLuan_Toan11_Lan1
\begin{bt}%[1D6H2-2]%[ma:f5a47937]
	\begin{enumerate}
		\item Cho số  dương $a$ khác $1$. Tính giá trị của biểu thức $\log_a\left(a^4a^5\right) $.
		\item Đặt $\log  _3 2=a$, $\log _3 5=b$. Biểu diễn $\log _{15} 4$ theo $a$ và $b$.
	\end{enumerate}
	\loigiai{ 
		\begin{enumerate}
			 \item Ta có $\log_a\left(a^4a^5\right) =\log_aa^9=9\log_aa=9$.
			\item Ta có $\log  _{15} 4=\dfrac{\log_34}{\log_315}=\dfrac{\log_32^2}{\log_3 3 +\log_35}=\dfrac{2a}{1+b}$.
		\end{enumerate}
	 }
\end{bt}

% nguồn: Ok_TuanLuan_Toan11_Lan1
\begin{bt}%[1D6H2-2]%[ma:5a71ae41]
	Cho  $a=\log _2 5, b=\log _3 5$. Hãy biểu diễn $\log _6 5$ theo $a$ và $b$.
	\loigiai{ 
		$\log_{6}5 = \dfrac{1}{\log_{5}6} = \dfrac{1}{\log_{5}2+\log_{5}3} =  \dfrac{1}{\dfrac{1}{a}+\dfrac{1}{b}} = \dfrac{ab}{a+b}$.
	 }
\end{bt}

% nguồn: Ok_TuanLuan_Toan11_Lan1
\begin{bt}%[1D6H2-2]%[ma:b60ad01c]
	Đặt $\log_23=a$ và $\log_25=b$. Biểu  thị $\log_{12}45$ theo $a$ và $b$.
	\loigiai{Ta  có $\log_{12}45=\dfrac{\log_245}{\log_212}=\dfrac{\log_29+\log_25}{\log_24+\log_23}=\dfrac{2\log_23+\log_25}{2+\log_23}=\dfrac{2a+b}{2+a}$.}
\end{bt}

% nguồn: Ok_TuanLuan_Toan11_Lan1
\begin{bt}%[1D6H2-2]%[ma:d75eadf5]
	Cho $a=\log_2 14$.  Biểu diễn $\log_{49}16 $ theo $a$.
	 \loigiai{Ta có $\log_{49}16=\log_{7^2}{2^4}=2\log_72=2\cdot \dfrac{1}{\log_27}=\dfrac{2}{\log_2{\left(\dfrac{14}{2}\right)}}=\dfrac{2}{\log_2{14}-\log_22}=\dfrac{2}{a-1}$.}
\end{bt}

% nguồn: Ok_TuanLuan_Toan11_Lan2
\begin{ex}%[1D6H2-2]%[ma:bd155980]
	Cho $\log_a b=3$ và $\log _a c=4$ với $a$; $b$; $c > 0$; $a \neq 1$. Tính giá  trị của $P=\log _a\left(\dfrac{a^2\cdot \sqrt{b}}{c^3}\right)$.
	\loigiai{ 
		Xét  biểu thức
		\begin{eqnarray*}
			P &=& \log _a\left(\dfrac{a^2 \cdot \sqrt{b}}{c^3}\right)  \\
			&=&  \log _a a^2+\log _a \sqrt{b}-\log _a c^3 \\
			&=& 2+\log _a b^{\tfrac{1}{2}}-3 \log _a c  \\
			 &=& 2+\dfrac{1}{2} \log _a b-3 \log _a c \\
			&=& 2+\dfrac{3}{2}-12 \\ 
			&=& - \dfrac{17}{2}. 
		\end{eqnarray*}
	 }
\end{ex}

% nguồn: Ok_TuanLuan_Toan11_Lan2
\begin{ex}%[1D6H2-2]%[ma:3d13edb7]
	Cho $\log_{a} x=2$, $\log_{b} x=3$. Giá trị của $P=\log_{\tfrac{a^2}{b}} x$ bằng  bao nhiêu?
	 \loigiai{
		Ta có $\log_{\tfrac{a^2}{b}}x=\dfrac{1}{\log_{x}{\tfrac{a^2}{b}}}=\dfrac{1}{\log_{x} a^2+\log_{x} b}=\dfrac{1}{2\log_{x} a+\log_{x} b}=\dfrac{1}{2\cdot \dfrac{1}{2}+\dfrac{1}{3}}=\dfrac{3}{4}$. 
	 }
\end{ex}

% nguồn: Ok_TuanLuan_Toan11_Lan2
\begin{ex}%[1D6H2-2]%[ma:82770569]
	\begin{enumerate}
		\item Cho số dương $a$ khác $1$. Tính giá trị của biểu thức $\log_a\left(a^4a^5\right) $. 
		\item Đặt $\log _3 2=a$, $\log _3 5=b$. Biểu diễn $\log _{15}  4$ theo $a$ và $b$.
	\end{enumerate}
	\loigiai{ 
		\begin{enumerate}
			\item Ta có $\log_a\left(a^4a^5\right)  =\log_aa^9=9\log_aa=9$.
			\item Ta có $\log _{15}  4=\dfrac{\log_34}{\log_315}=\dfrac{\log_32^2}{\log_3 3 +\log_35}=\dfrac{2a}{1+b}$.
		\end{enumerate}
	} 
\end{ex}

% nguồn: Ok_TuanLuan_Toan11_Lan2
\begin{ex}%[1D6H2-2]%[ma:2c328f9b]
	Cho  $a=\log _2 5, b=\log _3 5$. Hãy biểu diễn $\log _6 5$ theo $a$ và $b$.
	\loigiai{ 
		$\log_{6}5 = \dfrac{1}{\log_{5}6} = \dfrac{1}{\log_{5}2+\log_{5}3} = \dfrac{1}{\dfrac{1}{a}+\dfrac{1}{b}} =  \dfrac{ab}{a+b}$.
	} 
\end{ex}

% nguồn: Ok_TuanLuan_Toan11_Lan2
\begin{ex}%[1D6H2-2]%[ma:3aec94e8]
	Đặt $\log_23=a$ và $\log_25=b$.  Biểu thị $\log_{12}45$ theo $a$ và $b$.
	\loigiai{Ta  có $\log_{12}45=\dfrac{\log_245}{\log_212}=\dfrac{\log_29+\log_25}{\log_24+\log_23}=\dfrac{2\log_23+\log_25}{2+\log_23}=\dfrac{2a+b}{2+a}$.}
\end{ex}

% nguồn: Ok_TuanLuan_Toan11_Lan2
\begin{ex}%[1D6H2-2]%[ma:b22634bc]
	Cho $a=\log_2  14$. Biểu diễn $\log_{49}16 $ theo $a$.
	\loigiai{Ta  có $\log_{49}16=\log_{7^2}{2^4}=2\log_72=2\cdot \dfrac{1}{\log_27}=\dfrac{2}{\log_2{\left(\dfrac{14}{2}\right)}}=\dfrac{2}{\log_2{14}-\log_22}=\dfrac{2}{a-1}$.}
\end{ex}

% nguồn: Ok_TuanLuan_Toan11_Lan2
\begin{ex}%[1D6H2-2]%[ma:25d29f6b]
	Đặt $a=\log _2 5$, $b=\log _3 5$. Hãy biểu diễn $\log _6 5$ theo  $a$ và $b$.
	\loigiai{ 
		Ta có  $a=\log _2 5 \Rightarrow \log _5 2=\dfrac{1}{a}$ và $b=\log _3 5 \Rightarrow \log _5 3=\dfrac{1}{b}$.\\
		Suy ra $\log _6 5=\dfrac{1}{\log _5 6}=\dfrac{1}{\log  _5(2\cdot 3)}=\dfrac{1}{\log _5 2+\log _5 3}=\dfrac{1}{\dfrac{1}{a}+\dfrac{1}{b}}=\dfrac{a b}{a+b}$.
	 }
\end{ex}

% nguồn: Ok_TuanLuan_Toan11_Lan2
\begin{ex}%[1D6H2-2]%[ma:a9b8c867]
	Với mọi số $a$, $b$ thỏa mãn $\log _a\sqrt{b}=3$. Tính giá trị của biểu thức  $\log _a\left(a^3b^2\right)$.
	\loigiai{ 
		Từ giả thiết $\log _a\sqrt{b}=3\Leftrightarrow \dfrac{1}{2}\log _ab=3\Leftrightarrow \log  _ab=6$.\\
		 Khi đó $\log _a\left(a^3b^2\right)=\log _aa^3+\log _ab^2=3+2\log _ab=15$.\\
	 }
\end{ex}

% nguồn: Ok_TuanLuan_Toan11_Lan2
\begin{ex}%[1D6H2-2]%[ma:e9dda5d3]
	Cho $a$ và  $b$ là hai số thực dương thỏa mãn $a^3b^2=32$. Tính giá trị của $3\log _2a+2\log _2b$.
	\loigiai{ 
		Ta  có
		$a^3b^2=32\Leftrightarrow \log _2\left(a^3b^2\right)=\log _232\Leftrightarrow \log _2a^3+\log _2b^2=5\Leftrightarrow  3\log _2a+2\log _2b=5$.
	 }
\end{ex}

% nguồn: Ok_TuanLuan_Toan11_Lan2
\begin{ex}%[1D6H2-2]%[ma:0fb3fc43]
	Đặt  $a=\log _2 3$, $b=\log _5 3$. Hãy biểu diễn $\log _6 45$ theo $a$ và $b$.
	 \loigiai{
		 Ta có $\log _6 45=\log _6 9+\log _6 5$.\\
		Lúc đó 
		$$ 
		\begin{aligned}
			& \log _6 9=2 \log _6 3=\dfrac{2}{\log _3 6}=\dfrac{2}{1+\log _3 2}=\dfrac{2}{1+\dfrac{1}{a}}=\dfrac{2 a}{a+1}.\\ 
			& \log _6 5=\dfrac{1}{\log _5 6}=\dfrac{1}{\log _5 3+\log _5 2}=\dfrac{a}{b(a+1)}, \text { vì }  \log _5 2=\dfrac{b}{a}.
		\end{aligned}
		$$ 
		 Vậy $\log _6 45=\dfrac{2 a}{a+1}+\dfrac{a}{b(a+1)}=\dfrac{a+2 a b}{a b+b}$.
	 }
\end{ex}

% nguồn: Ok_TuanLuan_Toan11_Lan2
\begin{ex}%[1D6H2-2]%[ma:9f48d2a4]
	Cho hai số thực $a$, $b$ thỏa  mãn $2 \log _3(a-3 b)=\log _3 a+\log _3(4 b)$ và $a>3 b>0$. Tính $\dfrac{a}{b}$.
	\loigiai{ 
		Ta  có
		 \allowdisplaybreaks
		\begin{eqnarray*}
			2 \log _3(a-3  b)=\log _3 a+\log _3(4 b) &\Leftrightarrow& \log_3(a-3 b)^2=\log _3(4 a b)\\
			&\Leftrightarrow& (a-3  b)^2=4 a b\\
			&\Leftrightarrow& a^2-10 a b+9 b^2=0\\ 
			&\Leftrightarrow&\left(\dfrac{a}{b}\right)^2-10  \dfrac{a}{b}+9=0\\ &\Leftrightarrow&\hoac{&
			\dfrac{a}{b}=9  \\&
			\dfrac{a}{b}=1 \text { (loại).}} 
		\end{eqnarray*}
	 }
\end{ex}

% nguồn: Ok_TuanLuan_Toan11_Lan2
\begin{ex}%[1D6H2-2]%[ma:d480fab8]
	Cho $a=\log _23$. Tính  $\log _672$.
	 \loigiai{
		Ta có $\log _672=\dfrac{\log _272}{\log _26}=\dfrac{\log _2\left(2^3\cdot  3^2\right)}{\log _2(2.3)}=\dfrac{3\log _22+2\log _23}{\log _22+\log _23}=\dfrac{2a+3}{a+1}$.\\
		Vậy $\log  _672=\dfrac{2a+3}{a+1}$.
	 }
\end{ex}

% nguồn: Ok_TuanLuan_Toan11_Lan2
\begin{ex}%[1D6H2-2]%[ma:818b38a3]
	Cho $a$,  $b$ là các số thực dương khác 1 thỏa mãn $\log _ab=\sqrt{3}$. Giá trị của $\log _{\tfrac{\sqrt{b}}{a}}\left(\dfrac{\sqrt[3]{b}}{\sqrt{a}}\right)$ là
	\loigiai{ 
		Ta  có $\log _ab=\sqrt{3}\Leftrightarrow b=a^{\sqrt{3}}\Rightarrow \log _{\tfrac{\sqrt{b}}{a}}\left(\dfrac{\sqrt[3]{b}}{\sqrt{a}}\right)=\log _{\tfrac{a^{\tfrac{\sqrt{3}}{2}}}{a}}\left(\dfrac{a^{\tfrac{\sqrt{3}}{3}}}{a^{\tfrac{1}{2}}}\right)=\dfrac{\dfrac{\sqrt{3}}{3}-\dfrac{1}{2}}{\dfrac{\sqrt{3}}{2}-1}=-\dfrac{1}{\sqrt{3}}$.
	} 
\end{ex}

% nguồn: Ok_TuanLuan_Toan11_Lan2
\begin{ex}%[1D6H2-2]%[ma:460baf2a]
	Tính  $\log _{81}25$ theo $\log _53=b$
	 \loigiai{
		Ta có $\log _{81}25=\dfrac{\log _525}{\log _581}=\dfrac{\log _55^2}{\log _53^4}  =\dfrac{2}{4\log _53}=\dfrac{1}{2b}$.
	 }
\end{ex}

% nguồn: Ok_TuanLuan_Toan11_Lan2
\begin{ex}%[1D6H2-2]%[ma:df44009c]
	Cho các số  dương $a$, $b$. Tìm điều kiện của $a$, $b$ để $\log_a\dfrac{1}{2}>\log_a\dfrac{1}{3}$ và $b^{\tfrac{2}{3}}>b^{\tfrac{3}{4}}.$
	\loigiai{ 
		Ta có $\log_a\dfrac{1}{2}>\log_a\dfrac{1}{3}$ và do $\dfrac{1}{2}>\dfrac{1}{3}$ nên $a>1$.\\ 
	Vì ${b^{\tfrac{2}{3}}}>{b^{\tfrac{3}{4}}}$ và  $\dfrac{2}{3}<\dfrac{3}{4}$ nên $0<b<1$.}
\end{ex}

% nguồn: Ok_TuanLuan_Toan11_Lan2
\begin{ex}%[1D6H2-2]%[ma:9b02cf68]
	 Cho $a=\log _{25}7$; $b=\log _25$. Tính $\log _5\dfrac{49}{8}$ theo $a$, $b$.
	 \loigiai{
		Ta  có $a=\log _{25}7=\log _{5^2}7=\dfrac{1}{2}\log _57\Rightarrow \log _57=2a$ và $b=\log _25\Rightarrow \log _52=\dfrac{1}{b}$.\\
		Suy ra $\log _5\dfrac{49}{8}=\log _549-\log _58=\log _57^2-\log _52^3=2\log _57-3\log _52=2\cdot 2a-3\cdot \dfrac{1}{b}=\dfrac{4ab-3}{b}$. 
	 }
\end{ex}

% nguồn: Ok_TuanLuan_Toan11_Lan2
\begin{ex}%[1D6H2-2]%[ma:b0f7f919]
	Cho $\log _ax=2$, $\log _bx=3$ với $a$,  $b$ là các số thực lớn hơn $1$. Tính $P=\log _{\tfrac{a}{b^2}}x$.
	 \loigiai{
		Với $a$, $b$ là các số thực lớn hơn $1$ và $x>0$,  $x\neq 1$, ta có
		\begin{eqnarray*}
			 P&=&\log _{\tfrac{a}{b^2}} x=\dfrac{1}{\log _x \dfrac{a}{b^2}}\\
			&=& \dfrac{1}{\log  _x a-2 \log _x b} \\
			&=&\dfrac{1}{\dfrac{1}{\log _a x}-\dfrac{2}{\log  _b x}}\\
			&=&\dfrac{1}{\dfrac{1}{2}-\dfrac{2}{3}}=-6. 
		\end{eqnarray*}
	} 
\end{ex}

% nguồn: Ok_TuanLuan_Toan11_Lan2
\begin{ex}%[1D6H2-2]%[ma:4635b1fc]
	Cho $\log _ab=3$, $\log _ac=-4$. Tính giá trị  của biểu thức $P=\log _a\left(\dfrac{a^3\sqrt{c}}{b^2}\right)$.
	 \loigiai{
		 Ta có
		\begin{eqnarray*}
			P&=&\log  _a\left(\dfrac{a^3 \sqrt{c}}{b^2}\right)\\
			&=&\log _a a^3  \sqrt{c}-\log _a b^2\\
			&=&\log _a a^3+\log _a  c^{\tfrac{1}{2}}-\log _a b^2\\
			&=&3+\dfrac{1}{2} \log _a c-2  \log _a b \\
			 &=&3+\dfrac{1}{2}(-4)-2\cdot3\\
			 &=&-5.
		\end{eqnarray*}
	 }
\end{ex}

% nguồn: Ok_TuanLuan_Toan11_Lan2
\begin{ex}%[1D6H2-2]%[ma:ef72d994]
	Cho các số dương $a$, $b$. Tìm điều kiện của $a$, $b$  để $\log_b\left(\sqrt{2}+\sqrt{5}\right)>\log_b\left(2+\sqrt{3}\right)$\\ và ${a^{\tfrac{13}{7}}}<{a^{\tfrac{15}{8}}.}$
	\loigiai{ 
		Ta có $\heva{& \dfrac{13}{7}<\dfrac{15}{8} \\ & a^{\tfrac{13}{7}}<a^{\tfrac{15}{8}}.}$\\ Suy ra $a>1$.\\ 
	Mặt khác  $\heva{& \sqrt{2}+\sqrt{5}<2+\sqrt{3} \\ & \log_b\left(\sqrt{2}+\sqrt{5}\right)>\log_b\left(2+\sqrt{3}\right).}$\\Suy ra $0<b<1$.}
\end{ex}

% nguồn: Ok_TuanLuan_Toan11_Lan2
\begin{ex}%[1D6H2-2]%[ma:e84146c1]
	 Cho $a$, $b$ là các số thực dương lớn hơn 1 thỏa mãn $\log _ab=3$. Tính gái trị biểu thức $P=\log _{a^2b}a^3-3\log _{a^2}2\cdot \log _4\left(\dfrac{a}{b}\right)$.
	\loigiai{ 
		Ta có  $\log _ab=3\Leftrightarrow b=a^3$.\\
		Suy  ra
		\begin{eqnarray*}
			P&=&\log _{a^2 b} a^3-3 \log _{a^2}  2 \cdot \log _4\left(\dfrac{a}{b}\right)\\
			&=&\log _{a^5} a^3-3 \log _{a^2} 2 \cdot  \log _4\left(\dfrac{a}{a^3}\right)\\
			&=&\dfrac{3}{5}-3 \cdot \dfrac{1}{2} \cdot \log  _a 2 \cdot \log _{2^2}\left(\dfrac{1}{a^2}\right). \\
			&=&\dfrac{3}{5}-3 \cdot \dfrac{1}{2} \cdot \log _a 2 \cdot \dfrac{1}{2} \cdot  \log _2 a^{-2}\\
			&=&\dfrac{3}{5}+\dfrac{3}{2} \cdot  \log _a 2 \cdot \log _2 a\\
			 &=&\dfrac{3}{5}+\dfrac{3}{2}\\
			 &=&\dfrac{21}{10}.
		\end{eqnarray*}
	 }
\end{ex}

% nguồn: Ok_TuanLuan_Toan11_Lan2
\begin{ex}%[1D6H2-2]%[ma:7b48a4f7]
	Cho hai số dương $a$, $b$, $a\neq 1$, thỏa mãn  $\log _{a^2}b+\log _ab^2=2$. Tính $\log _ab$.
	\loigiai{ 
		Ta  có $\log _{a^2}b+\log _ab^2=2\Leftrightarrow \dfrac{1}{2}\log _ab+2\log _ab=2\Leftrightarrow \log _ab=\dfrac{4}{5}$.
	} 
\end{ex}

% nguồn: Ok_TuanLuan_Toan11_Lan2
\begin{ex}%[1D6H2-2]%[ma:3e47a7eb]
	Cho $a$, $b$  là các số thực dương và $a$ khác 1, thỏa mãn $\log _{a^2}\left(\dfrac{a^3}{\sqrt[5]{b^3}}\right)=3$. Tính $\log _ab$.
	\loigiai{ 
		Ta có $\log _{a^2}\left(\dfrac{a^3}{\sqrt[5]{b^3}}\right)=3\Leftrightarrow \dfrac{1}{2}\left(\log _aa^3-\log _ab^{\tfrac{3}{5}}\right)=3\Leftrightarrow 3-\dfrac{3}{5}\log _ab=6\Leftrightarrow  \log _ab=-5$.
	 }
\end{ex}

% nguồn: Ok_TuanLuan_Toan11_Lan2
\begin{ex}%[1D6H2-2]%[ma:e0868db6]
	Cho $\log _23=a$. Tính $P=\log _612$. 
	\loigiai{ 
		Ta có $P=\log _612=\log _6(6\cdot2)=1+\log _62=1+\dfrac{1}{\log _26}=1+\dfrac{1}{1+\log _23}=1+\dfrac{1}{1+a}=\dfrac{2+a}{1+a}$. 
	 }
\end{ex}

% nguồn: Ok_TuanLuan_Toan11_Lan2
\begin{ex}%[1D6H2-2]%[ma:768441e1]
	Với mọi số thực $a$ dương  thoả mãn $\log _3a=5$. Tính $\log _3(3a)$.
	\loigiai{ 
		Ta có $\log _3(3a)=\log _33+\log _3a=1+\log _3a=6.$ 
	} 
\end{ex}

% nguồn: Ok_TuanLuan_Toan11_Lan2
\begin{ex}%[1D6H2-2]%[ma:956c72c0]
	Cho $\log  _23=a$. Hãy tính $\log _92$ theo $a$.
	 \loigiai{
		Ta có $\log _92=\log _{3^2}a=\dfrac{1}{2}\log  _32=\dfrac{1}{2\log _23}=\dfrac{1}{2a}$.
	} 
\end{ex}

% nguồn: Ok_TuanLuan_Toan11_Lan2
\begin{ex}%[1D6H2-2]%[ma:c6830fb6]
	 Cho $a=\log _25$ và $b=\log _57$. Tính $\log _{\sqrt[3]{5}}\dfrac{49}{8}$ theo $a$ và $b$.
	 \loigiai{
		Ta có $\log _{\sqrt[3]{5}}\dfrac{49}{8}=3\log  _5\dfrac{49}{8}=6\log _57-9\log _52=6b-\dfrac{9}{a}=3\left(2b-\dfrac{3}{a}\right)$.
	 }
\end{ex}

% nguồn: Ok_TuanLuan_Toan11_Lan2
\begin{ex}%[1D6H2-2]%[ma:800c9885]
	Xét các số thực dương $a$, $b$ thỏa mãn $\log _5a=5$ và $\log _3b=\dfrac{2}{3}$.  Tính giá trị biểu thức $I=2\log _6\left[\log _5(5a)\right]+\log _{\tfrac{1}{9}}b^3$.
	\loigiai{ 
		Ta có $I=2\log _6\left[\log _5(5a)\right]+\log _{\tfrac{1}{9}}b^3=2\log _6\left(1+\log _5a\right)-\dfrac{3}{2}\log _3b=2\log _66-\dfrac{3}{2}\cdot \dfrac{2}{3}=2-1=1$. 
	} 
\end{ex}

% nguồn: Ok_TuanLuan_Toan11_Lan2
\begin{ex}%[1D6H2-2]%[ma:c865b307]
	Cho $a$, $b$ là hai số thực dương thỏa mãn $4^{\log _2\left(a^2\cdot b\right)}=3a^3$. Tinh giá  trị của $ab^2$.
	\loigiai{ 
		 Ta có $4^{\log _2\left(a^2\cdot b\right)}=3a^3\Leftrightarrow\left(a^2\cdot b\right)^2=3a^3\Leftrightarrow ab^2=3$.
	} 
\end{ex}

% nguồn: Ok_TuanLuan_Toan11_Lan2
\begin{ex}%[1D6H2-2]%[ma:d26a2b72]
	Biết $a=\log _23$, $b=\log _35$. Tính $\log _25$ theo $a$  và $b$.
	\loigiai{ 
		Ta có $\log _25=\log  _23\cdot \log _35=ab$.
	 }
\end{ex}

% nguồn: Ok_TuanLuan_Toan11_Lan2
\begin{ex}%[1D6H2-2]%[ma:b43ca2ba]
	Cho $a$, $b$, $c>0$, $a\neq 1$ và  $\log _ab=\sqrt{2\,022}$. Tính $\log _{\sqrt[6]{a}}\left(a^{\tfrac{7}{4}}\cdot \sqrt[6]{b}\right)$.
	\loigiai{ 
		Ta có $\log _{\sqrt[6]{a}}\left(a^{\tfrac{7}{4}}\cdot \sqrt[6]{b}\right)=\log _{\sqrt[6]{a}}a^{\tfrac{7}{4}}+\log  _{\sqrt[6]{a}}\sqrt[6]{b}=6\cdot \dfrac{7}{4}+\sqrt{2\,022}=\dfrac{21}{2}+\sqrt{2\,022}$.
	} 
\end{ex}

% nguồn: Ok_TuanLuan_Toan11_Lan2
\begin{ex}%[1D6H2-2]%[ma:6f5f7ea2]
	Cho $a$ và $b$ là hai số thực dương thỏa mãn $ab^2=9$. Tính $\log _3a+2\log  _3b$.
	 \loigiai{
		Ta có $ab^2=9\Rightarrow \log _3\left(ab^2\right)=\log _39\Rightarrow \log _3a+2\log  _2b=2$.
	} 
\end{ex}

% nguồn: Ok_TuanLuan_Toan11_Lan2
\begin{ex}%[1D6H2-2]%[ma:dfedffff]
	Cho hai số thực $a$, $b$ biết $0<a<b<1$. Hãy so sánh  $\log_ab$ và $\log_ba$.
	\loigiai{ 
	Do $0<a<b<1$ nên ta có $1=\log_aa>\log_ab$;  $\log_ba>\log_bb=1$.\\ Vậy ta có $\log_ab<\log_ba$. }
\end{ex}

% nguồn: Ok_TuanLuan_Toan11_Lan2
\begin{ex}%[1D6H2-2]%[ma:2d1cc714]
	Cho $\log _35=a$; $\log _57=b$. Tính $\log _{45}175$. 
	\loigiai{ 
		Ta có $\log _{45}175=\dfrac{\log _55^2\cdot 7}{\log _53^2\cdot 5}=\dfrac{2+b}{1+2\log  _53}=\dfrac{2+b}{1+\dfrac{2}{a}}=\dfrac{a(2+b)}{2+a}$.
	 }
\end{ex}

% nguồn: Ok_TuanLuan_Toan11_Lan2
\begin{ex}%[1D6H2-2]%[ma:405defbb]
	Tính giá trị của  biểu thức $P=\log _ab^3+\log _{a^3}b^6$ trong đó $a$, $b$ là các số thực dương tùy ý và $a\neq 1$.
	\loigiai{ 
		Ta có $P=\log _ab^3+\log _{a^3}b^6=3\log _ab+6\cdot \dfrac{1}{3}\log _ab=5\cdot \log _ab$. 
	} 
\end{ex}

% nguồn: Ok_TuanLuan_Toan11_Lan1
\begin{bt}%[1D6V2-2]%[ma:e0a640b2]
	Cho $\log_{a}x=3$, $\log_{b}x=4$ với $a>1,b>1,x>1$. Tính  $P=\log_{ab}x$.
	 \loigiai{
	Ta có $P=\log_{ab}x=\dfrac{1}{\log_{x}(ab)}=\dfrac{1}{\log_{x}a+\log_{x}b}=1:\left(\dfrac{1}{3}+\dfrac{1}{4} \right)=\dfrac{12}{7}$. 
	 }
\end{bt}

% nguồn: Ok_TuanLuan_Toan11_Lan1
\begin{bt}%[1D6V2-2]%[ma:6c05dcca]
	Đặt $a=\log_3 4$, $b=\log_{5} 4$. Hãy biểu diễn $\log_{12}  80$ theo $a$ và $b$.
	 \loigiai{
		 Ta có
		\allowdisplaybreaks 
		\begin{eqnarray*}
			\log_{12}  80
			 = \log_{12} (5 \cdot 4^2)
			&=& \log_{12} 5 + 2\cdot\log_{12} 4\\ 
			&= &\dfrac{1}{\log_{5} (3  \cdot 4)} + \dfrac{2}{\log_{4} (3 \cdot 4)}\\
			&=& \dfrac{1}{\log_{5} 3 +  \log_{5} 4} + \dfrac{2}{\log_{4} 3+1}\\
			&=& \dfrac{1}{\dfrac{\log_4 3}{\log_4 5}+\log_5 4} +  \dfrac{2}{\dfrac{1}{\log_3 4}+1}\\
			&=&\dfrac{1}{\dfrac{\log_5 4}{\log_3 4}+\log_5 4} + \dfrac{2}{\dfrac{1}{\log_3 4}+1}\\ 
			 &=&\dfrac{1}{\dfrac{b}{a}+b} + \dfrac{2}{\dfrac{1}{a}+1}\\
			&=& \dfrac{2ab+a}{ab+b}. 
	\end{eqnarray*}}
\end{bt}

% nguồn: Ok_TuanLuan_Toan11_Lan1
\begin{bt}%[1D6V2-2]%[ma:62d95ef3]
	\begin{enumEX}{1}
		\item Rút gọn biểu thức $\dfrac{2^{1-\sqrt{5}}  \cdot 6^{2+\sqrt{5}}}{3^{\sqrt{5}+3}}$.
		\item Cho $m=\log 2$; $n=\log 3$. Tính $Q = \log_23 - \log  _{20} 6$ theo $m$, $n$.
	\end{enumEX}
	\loigiai{ 
		\begin{enumEX}{1}
			\item  Ta có
			\begin{eqnarray*}
				\dfrac{2^{1-\sqrt{5}} \cdot 6^{2+\sqrt{5}}}{3^{\sqrt{5}+3}} & = & \dfrac{2^{1-\sqrt{5}} \cdot 2^{2+\sqrt{5}} \cdot 3^{2+\sqrt{5}}}{3^{\sqrt{5}+3}}  \\
				& = & 2^{1-\sqrt{5} + 2+\sqrt{5}} \cdot 3^{2+\sqrt{5}  - \sqrt{5} - 3} \\
				 & = & 2^3 \cdot 3^{-1} \\
				&  = & \dfrac{8}{3}.
			\end{eqnarray*}
			\item Ta có $\log_2 3  = \dfrac{\log 3}{\log 2} = \dfrac{n}{m}$. \\
			Lại có $\log_{20} 6 = \dfrac{\log 6}{\log 20} = \dfrac{\log (2 \cdot 3)}{\log (2 \cdot 10)} = \dfrac{\log 2 + \log 3}{\log 2 + \log 10} = \dfrac{m + n}{m + 1}$. \\ 
			Vậy $Q = \log_2 3 - \log_{20} 6 = \dfrac{n}{m}  - \dfrac{m + n}{m + 1} = \dfrac{nm + n - m^2 - mn}{m(m+1)} = \dfrac{n - m^2}{m(m+1)}$.
		\end{enumEX}
	} 
\end{bt}

% nguồn: Ok_TuanLuan_Toan11_Lan1
\begin{bt}%[1D6V2-2]%[ma:1b031add]
	Cho $a=\log2$,  $b=\ln2$. Hãy biểu diễn $\ln800$ theo $a$ và $b$.
	\loigiai{ 
		Ta  có
		\begin{eqnarray*}
			 &\ln800 & =\ln\left(2^3\cdot10^2\right)\\
			 & & =3\ln2+2\ln10\\
			 & & =3\ln2+2\ln2\cdot\log_210\\
			& &  =3\ln2+\dfrac{2\ln2}{\log2}\\
			&  & =3b+\dfrac{2b}{a}.
		\end{eqnarray*}
	} 
\end{bt}

% nguồn: Ok_TuanLuan_Toan11_Lan2
\begin{ex}%[1D6V2-2]%[ma:3435959b]
	Cho $\log_{a}x=3$, $\log_{b}x=4$ với $a>1,b>1,x>1$.  Tính $P=\log_{ab}x$.
	\loigiai{ 
	Ta có $P=\log_{ab}x=\dfrac{1}{\log_{x}(ab)}=\dfrac{1}{\log_{x}a+\log_{x}b}=1:\left(\dfrac{1}{3}+\dfrac{1}{4}  \right)=\dfrac{12}{7}$.
	} 
\end{ex}

% nguồn: Ok_TuanLuan_Toan11_Lan2
\begin{ex}%[1D6V2-2]%[ma:d8097c75]
	Đặt  $a=\log_3 4$, $b=\log_{5} 4$. Hãy biểu diễn $\log_{12} 80$ theo $a$ và $b$.
	\loigiai{ 
		Ta có 
		 \allowdisplaybreaks
		\begin{eqnarray*}
			\log_{12}  80
			= \log_{12}  (5 \cdot 4^2)
			&=& \log_{12} 5 + 2\cdot\log_{12} 4\\ 
			&= &\dfrac{1}{\log_{5} (3 \cdot 4)} +  \dfrac{2}{\log_{4} (3 \cdot 4)}\\
			&=& \dfrac{1}{\log_{5} 3 + \log_{5} 4} + \dfrac{2}{\log_{4} 3+1}\\ 
			&=& \dfrac{1}{\dfrac{\log_4 3}{\log_4 5}+\log_5 4} + \dfrac{2}{\dfrac{1}{\log_3  4}+1}\\
			&=&\dfrac{1}{\dfrac{\log_5 4}{\log_3 4}+\log_5 4}  + \dfrac{2}{\dfrac{1}{\log_3 4}+1}\\
			&=&\dfrac{1}{\dfrac{b}{a}+b} + \dfrac{2}{\dfrac{1}{a}+1}\\ 
			&=&  \dfrac{2ab+a}{ab+b}.
	\end{eqnarray*}}
\end{ex}

% nguồn: Ok_TuanLuan_Toan11_Lan2
\begin{ex}%[1D6V2-2]%[ma:c29f11d4]
	\begin{enumEX}{1}
		\item Rút  gọn biểu thức $\dfrac{2^{1-\sqrt{5}} \cdot 6^{2+\sqrt{5}}}{3^{\sqrt{5}+3}}$.
		\item Cho $m=\log 2$;  $n=\log 3$. Tính $Q = \log_23 - \log _{20} 6$ theo $m$, $n$.
	\end{enumEX}
	\loigiai{ 
		\begin{enumEX}{1}
			 \item Ta có
			\begin{eqnarray*}
				\dfrac{2^{1-\sqrt{5}} \cdot 6^{2+\sqrt{5}}}{3^{\sqrt{5}+3}} & = & \dfrac{2^{1-\sqrt{5}} \cdot 2^{2+\sqrt{5}}  \cdot 3^{2+\sqrt{5}}}{3^{\sqrt{5}+3}} \\
				& =  & 2^{1-\sqrt{5} + 2+\sqrt{5}} \cdot 3^{2+\sqrt{5} - \sqrt{5} - 3} \\
				 & = & 2^3 \cdot 3^{-1} \\
				& =  & \dfrac{8}{3}.
			\end{eqnarray*}
			\item Ta có $\log_2  3 = \dfrac{\log 3}{\log 2} = \dfrac{n}{m}$. \\
			Lại có $\log_{20} 6 = \dfrac{\log 6}{\log 20} = \dfrac{\log (2 \cdot 3)}{\log (2 \cdot 10)} =  \dfrac{\log 2 + \log 3}{\log 2 + \log 10} = \dfrac{m + n}{m + 1}$. \\
			Vậy $Q = \log_2 3 - \log_{20}  6 = \dfrac{n}{m} - \dfrac{m + n}{m + 1} = \dfrac{nm + n - m^2 - mn}{m(m+1)} = \dfrac{n - m^2}{m(m+1)}$.
		\end{enumEX}
	 }
\end{ex}

% nguồn: Ok_TuanLuan_Toan11_Lan2
\begin{ex}%[1D6V2-2]%[ma:14d303b2]
	Cho  $a=\log2$, $b=\ln2$. Hãy biểu diễn $\ln800$ theo $a$ và $b$.
	 \loigiai{
		Ta  có
		\begin{eqnarray*}
			&\ln800 &  =\ln\left(2^3\cdot10^2\right)\\
			&  & =3\ln2+2\ln10\\
			& &  =3\ln2+2\ln2\cdot\log_210\\
			& & =3\ln2+\dfrac{2\ln2}{\log2}\\ 
			&  & =3b+\dfrac{2b}{a}.
		\end{eqnarray*}
	 }
\end{ex}

% nguồn: Ok_TuanLuan_Toan11_Lan2
\begin{ex}%[1D6V2-2]%[ma:f9451fff]
	Cho hai số thực dương $m$, $n$ \;$(n\ne  1)$ thỏa mãn $\dfrac{\log_7m\cdot\log_27}{\log_210-1}=3+\dfrac{1}{\log_n5}$. Tìm mối quan hệ giữa $m$ và $n$.
	\loigiai{ 
		 Ta có \allowdisplaybreaks
		\begin{align*}
			&\dfrac{\log_7m\cdot \log_27}{\log_210-1}=3+\dfrac{1}{\log_n5}\\\Leftrightarrow &\dfrac{\log_2m}{\log_25}=3+\log_5n\\ 
			\Leftrightarrow &\log_5m=3+\log_5n\\\Leftrightarrow  &\log_5m-\log_5n=3\\\Leftrightarrow &\log_5\dfrac{m}{n}=3\\\Leftrightarrow& \dfrac{m}{n}=5^3=125\Leftrightarrow m=125n.
		\end{align*}
		 Vậy $m=125n$.
	} 
\end{ex}

% nguồn: Ok_TuanLuan_Toan11_Lan2
\begin{ex}%[1D6V2-2]%[ma:39eac13d]
	Đặt $a=\log _3 2$, khi đó $\log _{72} 768$ được  biểu diễn dưới dạng $\dfrac{m a+n}{p a+2}$, với $m$, $n$, $p$ là các số nguyên. Tính $m+n^2+p^3$.
	\loigiai{ 
		 Ta có
		 \allowdisplaybreaks
		\begin{eqnarray*}
			\log _{72} 768=\log  _{72}\left(2^8 \cdot 3\right)=\log _{72} 2^8+\log _{72} 3&=&\dfrac{8}{\log _2\left(2^3 \cdot 3^2\right)}+\dfrac{1}{\log _3\left(2^3 \cdot 3^2\right)}\\
			&=&\dfrac{8}{3+2 \log _2 3}+\dfrac{1}{2+3 \log  _3 2}\\
			&=&\dfrac{8}{3+\dfrac{2}{a}}+\dfrac{1}{2+3  a}\\
			 &=&\dfrac{8 a+1}{3 a+2}.
		\end{eqnarray*}
		Suy ra $m=8$,  $n=1$, $p=3$.\\
		Vậy $m+n^2+p^3=8+1^2+3^3=36$. 
	 }
\end{ex}

% nguồn: Ok_TuanLuan_Toan11_Lan2
\begin{ex}%[1D6V2-2]%[ma:fb59fa3a]
	Cho  $a$, $b>0$;$a$, $b\ne 1$ thỏa mãn $\log_a^2b-8\log_b\left(a\cdot \sqrt[3]{b}\right)=\dfrac{-8}{3}$.\\ Tính giá trị $P=\log_a\left(a\cdot \sqrt[3]{ab}\right)+2019$.
	\loigiai{ 
		\allowdisplaybreaks 
		\begin{eqnarray*}
			\log_a^2b-8\log_b\left(a\cdot\sqrt[3]{b}\right) =  \dfrac{-8}{3} &\Leftrightarrow& \log_a^2b-8\log_ba = 0 \\
			&\Leftrightarrow &\log_a^3b-8=0  \Leftrightarrow \log_ab=2\\
			&\Leftrightarrow& a^2=b. 
		\end{eqnarray*}
		 Khi đó
		\begin{align*}
			P&=\log_a\left(a\cdot\sqrt[3]{ab}\right)+2019\\&=1+\dfrac{1}{3}\cdot  (1+\log_aa^2)+2019\\&=1+\dfrac{1}{3}\cdot (1+2)+2019=2021.
		\end{align*}
	 }
\end{ex}

% nguồn: Ok_TuanLuan_Toan11_Lan2
\begin{ex}%[1D6V2-2]%[ma:6753b09b]
	Cho các số thực dương $x\neq  1$, $y\neq 1$ thỏa mãn $\log _2x=\log _y16$ và tích $xy=64$. Tính giá trị của biểu thức $\left(\log _2\dfrac{y}{x}\right)^2$.
	\loigiai{ 
		Đặt $\log  _2x=\log _y16=t$.\\
		Suy  ra $\heva{&x=2^t\\&4\log _y2=t}\Leftrightarrow\heva{&x=2^t\\&\log _y2=\dfrac{t}{4}}\Leftrightarrow\heva{&x=2^t\\&\log _2y=\dfrac{4}{t}}\Leftrightarrow\heva{&x=2^t\\&y=2^{\tfrac{4}{t}}.}$.\\
		Ta  có $xy=64\Leftrightarrow 2^t\cdot 2^{\tfrac{4}{t}}=2^6\Leftrightarrow t+\dfrac{4}{t}=6$.\\
		Suy ra $\left(\log  _2\dfrac{y}{x}\right)^2=\left(\log _2y-\log _2x\right)^2=\left(\dfrac{4}{t}-t\right)^2=\dfrac{16}{t^2}-8+t^2=\left(t+\dfrac{4}{t}\right)^2-16=6^2-16=20$.
	 }
\end{ex}

% nguồn: Ok_TuanLuan_Toan11_Lan2
\begin{ex}%[1D6V2-2]%[ma:9f4b1511]
	Cho các số thực $a$, $b$ thỏa mãn $a>b>1$ và $\dfrac{1}{\log_{b}a}+\dfrac{1}{\log_{a}b}=\sqrt{2\,022}.$ Tính giá trị của biểu thức  $P=\dfrac{1}{\log_{ab}b}-\dfrac{1}{\log_{ab}a}.$
	 \loigiai{
		\begin{itemize}
			\item Ta  có $\dfrac{1}{\log_{b}a}+\dfrac{1}{\log_{a}b}=\sqrt{2\,022}\Leftrightarrow \log_{a}b+\log_{b}a=\sqrt{2\,022}\left(^{*}\right)$.
			\item Biến đổi biểu thức $ P=\dfrac{1}{\log_{ab}b}-\dfrac{1}{\log_{ab}a} 
			=\log_{b}\left(ab\right)-\log_{a}\left(ab\right) 
			=\log_{b}a-\log_{a}b$. 
			\item Đặt  $t=\log_{a}b$ thì (*) trở thành\\
			$t+\dfrac{1}{t}=\sqrt{2\,022}\Leftrightarrow t^{2}-\sqrt{2\,022}t+1=0\Leftrightarrow\hoac{&t=\dfrac{\sqrt{2\,022}+\sqrt{2\,018}}{2}\Rightarrow  P=\dfrac{1}{t}-t=-\sqrt{2\,018}\\ &t=\dfrac{\sqrt{2\,022}-\sqrt{2\,018}}{2}\Rightarrow P=\dfrac{1}{t}-t=\sqrt{2\,018}}$.
			\item Vì $a>b>1\Leftrightarrow0<\log_{a}b<1$ nên $0<t<1\Rightarrow\dfrac{1}{t}>1\Rightarrow P=\dfrac{1}{t}-t>0\Rightarrow P=\sqrt{2\,018}$. 
		\end{itemize}
	} 
\end{ex}

% nguồn: Ok_TuanLuan_Toan11_Lan2
\begin{ex}%[1D6V2-2]%[ma:1c3e5ec6]
	Cho các số dương $a$, $b$, $c$ khác $1$ thỏa mãn $\log_a(bc)=2$,  $\log_b(ca)=4$. Tính giá trị của biểu thức $\log_c(ab)$.
	\loigiai{ 
		 Ta có \allowdisplaybreaks
		\begin{eqnarray*}
			\heva{&\log_a(bc)=2\\  &\log_b(ca)=4}
			&\Leftrightarrow& \heva{&  \dfrac{\log_c(bc)}{\log_ca}=2 \\ & \dfrac{\log_c(ca)}{\log_c b}=4} \\
			&\Leftrightarrow& \heva{& \log_cb+1=2\log_ca \\ &  1+\log_ca=4\log_cb} \\
			&\Leftrightarrow& \heva{& 2\log_ca-\log_cb=1 \\ &  \log_ca-4\log_cb=-1}\\
			&\Leftrightarrow &\heva{& \log_ca=\dfrac{5}{7} \\ &  \log_cb=\dfrac{3}{7}.}
		\end{eqnarray*}
	Vậy  $\log_c(ab)=\log_ca+\log_cb=\dfrac{8}{7}$.}
\end{ex}

% nguồn: Ok_TuanLuan_Toan11_Lan2
\begin{ex}%[1D6V2-2]%[ma:97e6c12d]
	Cho các số thực dương $a,b$ khác $1$ thoả mãn $\log_2a=\log_b16$ và $ab=64$. Tính  giá trị của biểu thức $\left(\log_2\dfrac{a}{b}\right)^2$.
	 \loigiai{
		Ta  có $$\log_2a=\log_b16\Leftrightarrow \log_2a=\dfrac{1}{\log_{16}b}\Leftrightarrow \log_2a\cdot \log_2b=4$$.
		Mặt khác $$ab=64\Rightarrow  \log_2ab=\log_264\Rightarrow \log_2a+\log_2b=6$$.
	Khi đó $$\left(\log_2\dfrac{a}{b}\right)^2=(\log_2a-\log_2b)^2=(\log_2a+\log_2b)^2-4\log_2a\cdot \log_2b=6^2-4\cdot 4=20.$$  }
\end{ex}

% nguồn: Ok_TuanLuan_Toan11_Lan2
\begin{ex}%[1D6V2-2]%[ma:46ca3544]
	Cho các số dương $a$, $b$ khác 1 sao cho $\log _{16}\sqrt[3]{a}=\log _{a^2}\sqrt[9]{b}=\log _b2$ . Tính giá  trị của biểu thức $\dfrac{b}{a^3}$.
	\loigiai{ 
		Đặt $\log _{16}\sqrt[3]{a}=\log _{a^2}\sqrt[9]{b}=\log _b2=t$. Suy  ra $\heva{&a=16^{3t}\\&b=a^{18t}\\&2=b^t}\Leftrightarrow\heva{&a=16^{3t}\\&b=16^{54t^2}\\&2=16^{54t^3}.}(*)$\\
		Từ $(*)$ suy ra $2=2^{216t^3}\Leftrightarrow 216t^3=1\Leftrightarrow t=\dfrac{1}{6}$.\\ 
		Suy ra  $a=4$, $b=4^3$.\\
		 Vậy $\dfrac{b}{a^3}=1$.
	} 
\end{ex}

% nguồn: Ok_TuanLuan_Toan11_Lan2
\begin{ex}%[1D6V2-2]%[ma:2cdc834f]
	Tính giá  trị biểu thức
	$P=\dfrac{1}{\log _22\,021!}+\dfrac{1}{\log _32\,021!}+\cdots+\dfrac{1}{\log  _{2\,021}2\,021!}$.
	 \loigiai{
		 Ta có
		\begin{eqnarray*}
			 P&=&\dfrac{1}{\log _22\,021!}+\dfrac{1}{\log _32\,021!}+\cdots+\dfrac{1}{\log _{2\,021}2\,021!}\\
			&=&\log _{2\,021!}2+\log _{2\,021!}3+\cdots  +\log _{2\,021!}2\,021\\
			 &=&\log _{2\,021!}2\cdot 3\cdots 2\,021\\
			&=&\log  _{2\,021!}2\,021!=1.
		\end{eqnarray*}
	} 
\end{ex}

% nguồn: Ok_TuanLuan_Toan11_Lan2
\begin{ex}%[1D6V2-2]%[ma:02064f25]
	Gọi $a$, $b$ là các số thực lớn hơn $1$ sao cho biểu thức $T=\left(\log  _{\sqrt[3]{a}}b\right)^3+\log _ba$ đạt giá trị nhỏ nhất. Tính giá trị của biểu thức $P=\log _{\sqrt{a}}\sqrt[4]{ab}$.
	\loigiai{ 
		Do $a$, $b$ lớn hơn $1$ nên $\log  _ab>0$. Khi đó
		\begin{eqnarray*}
			T&=&\left(\log _{\sqrt[3]{a}}b\right)^3+\log  _ba\\
			 &=&27\log _a^3b+\dfrac{1}{3\log _ab}+\dfrac{1}{3\log _ab}+\dfrac{1}{3\log _ab}\\
			&\geq& 4\sqrt[4]{27\log _a^3b\cdot \dfrac{1}{3\log  _ab}\cdot \dfrac{1}{3\log _ab}\cdot \dfrac{1}{3\log _ab}}=4.
		\end{eqnarray*}
		Dấu bằng xảy ra khi  $\log _ab=\dfrac{1}{3}\Leftrightarrow a=b^3$.\\
		 Khi đó $P=\log _{\sqrt{a}}\sqrt[4]{ab}=\log _{\sqrt{b^3}}b=\dfrac{2}{3}$.
	} 
\end{ex}

% nguồn: Ok_TuanLuan_Toan11_Lan2
\begin{ex}%[1D6V2-2]%[ma:b274cea1]
	Cho số thực dương $a$,  $b$ thỏa mãn $\log _{16}a=\log _{20}b=\log _{25}\dfrac{2a-b}{3}$. Tính $\dfrac{a}{b}$.
	\loigiai{ 
		Đặt $\log  _{16}a=\log _{20}b=\log _{25}\dfrac{2a-b}{3}=x\Leftrightarrow\heva{&a=16^x\\&b=20^x\\&2a-b=3\cdot25^x.}$\\
		Suy ra 
		\begin{eqnarray*}
			 2\cdot 16^x-20^x=3\cdot 25^x
			 &\Leftrightarrow& 2\cdot \left(\dfrac{16}{20}\right)^x-3\cdot \left(\dfrac{25}{20}\right)^x-1=0\\
			&\Leftrightarrow& 2\cdot  \left(\dfrac{4}{5}\right)^x-3\cdot \left(\dfrac{5}{4}\right))^x-1=0\\
			 &\Leftrightarrow& 2\cdot \left(\dfrac{4}{5}\right)^{2x}-\left(\dfrac{4}{5}\right)^x-3=0\\
			&\Leftrightarrow&\hoac{&\left(\dfrac{4}{5}\right)^x=-1(\text {vô  nghiệm})\\&\left(\dfrac{4}{5}\right)^x=\dfrac{3}{2}.}
		\end{eqnarray*}
		Vậy $\dfrac{a}{b}=\dfrac{16^x}{20^x}=\left(\dfrac{4}{5}\right)^x=\dfrac{3}{2}\in(1;2)$. 
	 }
\end{ex}

% nguồn: Ok_TuanLuan_Toan11_Lan2
\begin{ex}%[1D6V2-2]%[ma:90be3e57]
	Biết $x$ và $y$ là hai  số thực thỏa mãn $\log _4x=\log _9y=\log _6(x-2y)$. Tính $\dfrac{x}{y}$ bằng
	 \loigiai{
		Đặt $\log _4x=\log  _9y=\log _6(x-2y)=t\Rightarrow\heva{&x=4^t\\&y=9^t\\&x-2y=6^t.}$\\
		Khi đó $4^t-2\cdot  9^t=6^t\Leftrightarrow\left(\dfrac{2}{3}\right)^{2t}-\left(\dfrac{2}{3}\right)^t-2=0\Leftrightarrow\hoac{&\left(\dfrac{2}{3}\right)^t=-1\quad\text {(loại)}\\&\left(\dfrac{2}{3}\right)^t=2.}$.\\
		Suy ra  $\dfrac{x}{y}=\left(\dfrac{2}{3}\right)^{2t}=2^2=4$.
	} 
\end{ex}

% nguồn: Ok_TuanLuan_Toan11_Lan2
\begin{ex}%[1D6V2-2]%[ma:8c6918ce]
	Xét $a$ và $b$ là hai số thực dương tùy ý. Đặt $x=\ln  (a^2-ab+b^2)^{1\,000}$; \break$y=1\,000\ln a-\ln \dfrac{1}{b^{1\,000}}$. Hãy so sánh $x$ và $y$.
	 \loigiai{
		Ta có $a^2-ab+b^2={{\left(a-\dfrac{1}{2}b\right)}^2}+\dfrac{3}{4}b^2>0$, $\forall a>0$;\; $b>0$.\\ 
		$x=\ln  {(a^2-ab+b^2)^{1\,000}}=1\,000\ln (a^2ab+b^2)$.\\
		$y=1\,000\ln a-\ln \dfrac{1}{b^{1\,000}}=1\,000\ln  a+1\,000\ln b=1\,000\ln (ab)$.\\
		Với $\forall a>0$; $b>0$  thì $(a-b)^2\geq 0\Leftrightarrow a^2-2ab+b^2\geq 0\Leftrightarrow a^2-ab+b^2\geq ab$. \\
		Suy ra $\ln  (a^2-ab+b^2)\geq \ln (ab)$.\\
	Nên $1\,000\ln (a^2-ab+b^2)\geq 1\,000\ln (ab)$ khi và chỉ khi $x\geq y$.\\Vậy $x\geq y$.} 
\end{ex}

% nguồn: Ok_TuanLuan_Toan11_Lan2
\begin{ex}%[1D6V2-2]%[ma:c175d364]
	Cho các số thực $a$,  $b$ thỏa mãn $a>b>1$ và $\dfrac{1}{\log _ba}+\dfrac{1}{\log _ab}=\sqrt{2\,022}$. Tính giá trị của biểu thức $P=\dfrac{1}{\log _{a b} b}-\dfrac{1}{\log _{a b} a}$.
	 \loigiai{
		Ta  có
		\begin{itemize}
			\item $\dfrac{1}{\log _ba}+\dfrac{1}{\log _ab}=\sqrt{2\,022}\Leftrightarrow  \log _ab+\log _ba=\sqrt{2\,022}$. (*)
			\item $P=\dfrac{1}{\log _{ab}b}-\dfrac{1}{\log _{ab}a}=\log _b(ab)-\log _a(ab)=\log _ba-\log  _ab$.
		\end{itemize}
		Đặt $t=\log _ab$ thì (*) trở  thành
		\[t+\dfrac{1}{t}=\sqrt{2\,022}  \Leftrightarrow t^2-\sqrt{2\,022} t+1=0 \Leftrightarrow\hoac{&
		t=\dfrac{\sqrt{2\,022}+\sqrt{2\,018}}{2}  \\&
		t=\dfrac{\sqrt{2\,022}-\sqrt{2\,018}}{2}.}\] 
		\begin{itemize}
			\item Với $t=\dfrac{\sqrt{2\,022}+\sqrt{2\,018}}{2}$, ta có  $P=\dfrac{1}{t}-t=-\sqrt{2\,018}$.
			\item Với $t=\dfrac{\sqrt{2\,022}-\sqrt{2\,018}}{2}$, ta có $P=\dfrac{1}{t}-t=\sqrt{2\,018}$. 
		\end{itemize}
		Mặt khác $a>b>1\Leftrightarrow  0<\log _ab<1$ nên $0<t<1\Rightarrow \dfrac{1}{t}>1\Rightarrow P=\dfrac{1}{t}-t>0\Rightarrow P=\sqrt{2\,018}$.
	} 
\end{ex}

% nguồn: Ok_TuanLuan_Toan11_Lan2
\begin{ex}%[1D6V2-2]%[ma:585ee26a]
	Cho hai số thực $a$, $b$ dương khác $1$ thỏa mãn $\dfrac{1}{\log _a b}+\dfrac{1}{\log _{a^2} b}+\dfrac{1}{\log _{a^n} b}=\dfrac{1}{\log _{a^8}  b}$. Tìm $n$.
	\loigiai{ 
		Vì $a$, $b$ dương khác $1$ nên ta có $\dfrac{1}{\log _a b}+\dfrac{1}{\log _{a^2} b}+\dfrac{1}{\log _{a^n} b}=\dfrac{1}{\log _{a^8}  b}$.\\
		Lúc đó $\log _b a+\log _b a^2+\log _b a^n=\log _b a^8 \Leftrightarrow(1+2+n-8) \log _b  a=0$.\\
		Vì hai số thực $a$, $b$ dương khác $1$ nên $\log _b a \neq 0$ nên từ $(1)$ suy  ra $1+2+n-8=0 \Leftrightarrow n=5$.
	 }
\end{ex}

% nguồn: Ok_TuanLuan_Toan11_Lan2
\begin{ex}%[1D6V2-2]%[ma:6a6c3095]
	Cho các số nguyên dương $x$, $y$, $z$ đôi một  nguyên tố cùng nhau và thỏa mãn biểu thức dưới đây $x \log _{200} 5+y \log _{200} 2=z$. Tính giá trị của biểu thức $29 x-y-2021 z$.
	\loigiai{ 
		Ta  có $x \log _{200} 5+y \log _{200} 2=z \Leftrightarrow \log _{200} 5^x+\log _{200} 2^y=z$.\\
		Lúc đó 
		\allowdisplaybreaks 
		\begin{eqnarray*}
			\log _{200} 5^x \cdot 2^y=z &\Leftrightarrow& 5^x \cdot 2^y=200^z\\ 
			&\Leftrightarrow& 5^x  \cdot 2^y=(25 \cdot 4 \cdot 2)^z\\
			& \Leftrightarrow &5^x \cdot  2^y=5^{2 z} \cdot 2^{3 z} \\
			 &\Rightarrow&\heva{&x=2 z \\&y=3z.}
		\end{eqnarray*}
		Do các số nguyên dương  $x$, $y$, $z$ đôi một nguyên tố cùng nhau nên $\heva{&x=2 \\& y=3 \\& z=1.}$\\
		Giá trị  của biểu thức $29 x-y-2021 z=29\cdot 2-3-2\,021=-1\,966$.
	} 
\end{ex}

% nguồn: Ok_TuanLuan_Toan11_Lan2
\begin{ex}%[1D6V2-2]%[ma:26bf61d5]
	 Cho $x$, $y$ là các số thực dương thỏa mãn $\log _{\tfrac{4}{3}} x=\log _3 y=\log _2(2 x-3 y)$. Tính $\dfrac{x}{y}$.
	 \loigiai{
		 Đặt $\log _{\tfrac{4}{3}} x=\log _3 y=\log _2(2 x-3 y)=t$.\\
		Suy ra $\heva{&x=\left(\dfrac{4}{3}\right)^t \\ &y=3^t \\& 2 x-3 y=2^t} \Rightarrow 2 \cdot\left(\dfrac{4}{3}\right)^t-3\cdot 3^t=2^t \Leftrightarrow 2 \cdot \left(\dfrac{2}{3}\right)^t-3  \cdot \left(\dfrac{3}{2}\right)^t-1=0$. \quad $(1)$\\
		Đặt $\left(\dfrac{2}{3}\right)^t=a$, $(a>0)$.\\ 
		Khi đó phương trình $(1)$ trở thành $2 a-\dfrac{3}{a}-1=0 \Leftrightarrow 2 a^2-a-3=0 \Leftrightarrow\hoac{&a=-1 & \text {  (loại)} \\ &a=\dfrac{3}{2} & \text { (nhận).}}$.\\
		Do đó $\dfrac{x}{y}=\left(\dfrac{4}{9}\right)^t=\left(\dfrac{2}{3}\right)^{2 t}=a^2=\dfrac{9}{4}$. 
	 }
\end{ex}

% nguồn: Ok_TuanLuan_Toan11_Lan2
\begin{ex}%[1D6V2-2]%[ma:9c905991]
	Cho các số thực $a$, $b$, $c$, $d$ thoả mãn $\dfrac{1}{2^a}+\dfrac{1}{4^b}+\dfrac{1}{8^c}+\dfrac{1}{16^d}=\dfrac{1}{4}$. Gọi  $m$ là giá trị nhỏ nhất của biểu thức $S=a+2b+3c+4d$. Giá trị của biểu thức $\log_2m$ bằng
	\loigiai{ 
		Ta  có
		\begin{eqnarray*} \dfrac{1}{4}=\dfrac{1}{2^a}+\dfrac{1}{4^b}+\dfrac{1}{8^c}+\dfrac{1}{16^d}&\geq& 4\sqrt[4]{\dfrac{1}{2^a}\cdot\dfrac{1}{4^b}\cdot\dfrac{1}{8^c}\cdot\dfrac{1}{16^d}}\\\Leftrightarrow \dfrac{1}{16}&\geq& \sqrt[4]{\dfrac{1}{2^{a+2b+3c+4d}}}\\\Leftrightarrow 2^{a+2b+3c+4d}&\geq& 2^{16}\\
			 \Leftrightarrow a+2b+3c+4d&\geq& 16.
		\end{eqnarray*}
		Dấu “=” xảy ra khi $\dfrac{1}{2^a}=\dfrac{1}{4^b}=\dfrac{1}{8^c}=\dfrac{1}{16^d}=\dfrac{1}{16}\Leftrightarrow \heva{&  a=4 \\ & b=2 \\ & c=\dfrac{4}{3} \\ & d=1. \\
	 }$\\ Khi đó $m=16\Rightarrow \log_2m=4$.}
\end{ex}

% nguồn: Ok_TuanLuan_Toan11_Lan2
\begin{ex}%[1D6V2-2]%[ma:19e34625]
	 Cho $a=\log _75$, $b=\log _35$. Tính $M=\log _{21}5$.
	 \loigiai{
		Ta có \begin{eqnarray*}
		M&=&\log _{21}5=\dfrac{1}{\log _521}=\dfrac{1}{\log _53+\log _57}=\dfrac{1}{\dfrac{1}{\log _35}+\dfrac{1}{\log _75}}\\ 
		&=&\dfrac{1}{\dfrac{\log _75+\log _35}{\log _75\cdot \log _35}}=\dfrac{\log _75\cdot \log _35}{\log _75+\log  _35}=\dfrac{ab}{a+b}.
	\end{eqnarray*}
 }
\end{ex}

% nguồn: Ok_TuanLuan_Toan11_Lan2
\begin{ex}%[1D6V2-2]%[ma:b297e1e7]
	Cho $x$, $y$ là hai số thực dương, $x\neq 1$ thỏa mãn $\log _{\sqrt{x}}y=\dfrac{2y}{5}$, $\log _{25}x=\dfrac{5}{2y}$. Tính  giá trị của $P=y^2-2x^2$.
	 \loigiai{
		Ta có $\heva{&\log _{\sqrt{x}}y=\dfrac{2y}{5}\\&\log  _{25}x=\dfrac{5}{2y}}\Leftrightarrow\heva{&\log _xy^2=\dfrac{2y}{5}\\&\log _x25=\dfrac{2y}{5}}\Leftrightarrow\heva{&y^2=25\\&\log _{25}x=\dfrac{5}{2y}}\Leftrightarrow\heva{&y=5\\&\log _{25}x=\dfrac{1}{2}}\Leftrightarrow\heva{&y=5\\&x=5.}$\\
		Vậy  $P=y^2-2x^2=-25$.
	 }
\end{ex}

% nguồn: Ok_TuanLuan_Toan11_Lan2
\begin{ex}%[1D6V2-2]%[ma:a1b39e05]
	Cho biết hai số thực dương $a$ và $b$  thỏa mãn $\log_a^2(ab)=4$; với $b>1>a>0$. Tính giá trị của biểu thức $\log_a^3(ab^2)$.
	 \loigiai{
		Với $b>1>a>0$ ta có 
		\allowdisplaybreaks 
		\begin{eqnarray*}
			\log_a^2(ab)=4&\Leftrightarrow& (\log_aa+\log_ab)^2=4\\ 
			&\Leftrightarrow&  (1+\log_ab)^2=4\\
			&\Leftrightarrow& \hoac{& 1+\log_ab=2 \\  & 1+\log_ab=-2}\\
			&\Leftrightarrow& \hoac{& \log_ab=1 \\ &  \log_ab=-3.}
		\end{eqnarray*}
		Vì $\heva{& 0<a<1 \\  & b>1 \\
		}$nên $\log_ab=-3$.\\ Khi  đó
		 $\log_a^3(ab^2)={(\log_aa+2\log_ab)}^3=\left(1+2\cdot(-3)\right)^3=-125$.
	 }
\end{ex}

% nguồn: Ok_TuanLuan_Toan11_Lan2
\begin{ex}%[1D6C2-2]%[ma:1ae8bb8c]
	Cho ba số thực dương $a$, $b$,  $c$ đều khác 1 thoả mãn $\log _ab=2\log _bc=4\log _ca$ và $a+2b+3c=48$. Tính $P=abc$.
	 \loigiai{
		 Do $a$, $b$, $c$ đều khác 1 nên $\log _ab$, $\log _bc$, $\log _ca$ đều khác $0$.\\
		 Ta có
		\begin{itemize}
			\item $\log _ab=2\log  _bc\Leftrightarrow \log _ac\cdot \log _cb=2\log _bc\Leftrightarrow \log _ac=2\log _b^2c$.
			\item $\log _ab=4\log _ca\Leftrightarrow \log  _ac\cdot \log _cb=4\log _ca\Leftrightarrow \log _cb=4\log _c^2a$.
		\end{itemize}
		Suy ra  $\log _ac\cdot \log _cb=8\log _b^2c\cdot \log _c^2a\Leftrightarrow \log _ab=8\log _b^2a\Leftrightarrow \log _a^3b=8\Leftrightarrow \log _ab=2\Leftrightarrow b=a^2$.\\
		Mặt khác $\log _ab=2\log _bc\Leftrightarrow \log _ab=2\log _{a^2}c\Leftrightarrow  b=c$.\\
		Ta lại có  $a+2b+3c=48\Leftrightarrow a+2a^2+3a^2=48\Leftrightarrow 5a^2+a-48=0\Leftrightarrow\hoac{&a=-\dfrac{16}{5}\\&a=3.}$\\
		Do $a$, $b$, $c$ đều là số thực dương nên $a=3\Rightarrow  b=c=9$.\\
		Vậy  $P=abc=243$.
	 }
\end{ex}

% nguồn: Ok_TuanLuan_Toan11_Lan2
\begin{ex}%[1D6C2-2]%[ma:b9238b4e]
	Cho $\log _95=a$, $\log _47=b$, $\log _23=c$. Biết $\log _{24}175=\dfrac{mb+nac}{pc+q}$ với $m$, $n$, $p$, $q\in \mathbb{Z}$ và $q$ là số nguyên  tố. Tính $A=mnpq$.
	\loigiai{ 
		Ta  có $\heva{&\log_95=a\\&\log_47=b\\&\log_23=c}\Leftrightarrow
		\heva{&\log_35=2a\\&\log_27=2b\\&\log_23=c} 
		\Leftrightarrow  \heva{&\log _7 3=\dfrac{c}{2 b} \\&\log _5 3=\dfrac{1}{2 a} \\&\log _5 2=\dfrac{1}{2ac}.}$\\
		Suy ra 
		\begin{eqnarray*}
			 \log _{24}175&=&\log _{2^3\cdot 3}5^2\cdot 7\\
			 &=&\log _{2^3\cdot 3}5^2+\log _{2^3\cdot 3}7\\
			&=&\dfrac{2}{\log _5 2^3 \cdot  3}+\dfrac{1}{\log _7 2^3 \cdot 3}\\
			&=&\dfrac{2}{3 \cdot \log _5 2+\log _5  3}+\dfrac{1}{3 \log _7 2+\log _7 3}.\\
			 &=&\dfrac{2}{\dfrac{3}{2ac}+\dfrac{1}{2a}}+\dfrac{1}{\dfrac{3}{2b}+\dfrac{c}{2b}}\\
			 &=&\dfrac{2}{\dfrac{3+c}{2ac}}+\dfrac{1}{\dfrac{3+c}{2b}}\\
			 &=&\dfrac{4ac}{3+c}+\dfrac{2b}{3+c}=\dfrac{4ac+2b}{c+3}.
		\end{eqnarray*}
		Do đó 
		$m=2$,  $n=4$, $p=1$, $q=3$.\\
		Vậy $mnpq=2\cdot4\cdot1\cdot3=24$. 
	} 
\end{ex}
